% Lesson 39 — Vocabulary: Words and expressions related to volunteering in gender equality promotion activities — Anglais Tle A — Cours signé OnBuch+
% Programme officiel MINESEC. Compiler avec : tectonic EN39-vocabulary-words-and-expressions-related-to-volunteerin.tex
\newcommand{\DOCMATIERE}{Anglais}
\newcommand{\DOCNIVEAU}{Tle A}
\newcommand{\DOCMODULE}{Chapter 4 : Citizenship and Human Rights}
\newcommand{\DOCLECON}{EN39}
\newcommand{\DOCTITRE}{Lesson 39 — Vocabulary: Words and expressions related to volunteering in gender equality promotion activities}
% OnBuch+ — préambule « cours signés » ENRICHI (compilé avec tectonic / XeLaTeX)
% Système visuel « Pop » d'OnBuch+ : fond crème (jamais blanc), bordures encre
% épaisses, ombres « dures » décalées, titres Archivo Black en capitales.
% Version MATHÉMATIQUES : graphes (pgfplots/tikz), boîtes pédagogiques
% (définition, propriété, méthode, attention, le savais-tu, exercice résolu,
% exercice, TP, bilan), page de garde et signature OnBuch+ ancrée en bas.
%
% Macros attendues dans main.tex AVANT \input{preamble} :
%   \DOCMATIERE  \DOCNIVEAU  \DOCMODULE  \DOCLECON (numéro)  \DOCTITRE
\documentclass[11pt]{article}

\usepackage{fontspec}
\usepackage[english,french]{babel} % français = langue principale ; l'anglais via \eng{..} / textebox
\usepackage[a4paper,margin=2cm,top=3.4cm,bottom=2.8cm,headheight=1.4cm,footskip=1.3cm]{geometry}
\usepackage{amsmath,amssymb,amsfonts}
\usepackage{mathtools} % \xmapsto, \coloneqq... souvent utilisés par les modèles
\usepackage{cancel} % simplifications barrées (bilans de forces)
% \abs{z} (module, valeur absolue) et \norm{u} (norme), utilisés par les modèles
\providecommand{\abs}{}\let\abs\relax\DeclarePairedDelimiter{\abs}{\lvert}{\rvert}
\providecommand{\norm}{}\let\norm\relax\DeclarePairedDelimiter{\norm}{\lVert}{\rVert}
\usepackage[table]{xcolor} % [table] : fournit aussi \rowcolors (alternance de lignes)
\usepackage{pagecolor}
\usepackage{tikz}
\usetikzlibrary{arrows.meta,positioning,calc,shapes.geometric,shapes.misc,shapes.arrows,decorations.pathmorphing,decorations.pathreplacing,decorations.markings,patterns,shadows,mindmap,trees,fit,backgrounds,matrix,intersections,angles,quotes}
\usepackage{pgfplots}
\usepackage[version=4]{mhchem} % équations nucléaires \ce{^{235}_{92}U}
\usepackage[european,siunitx]{circuitikz} % schémas électriques
\pgfplotsset{compat=1.18}
\usepgfplotslibrary{fillbetween,groupplots} % aires entre courbes (calcul intégral)
\usepackage{enumitem}
\usepackage{fancyhdr}
% [most] charge toutes les bibliothèques tcolorbox (listings, minted, raster,
% documentation...) et gonfle inutilement la mémoire TeX sur un document long
% (~300 boîtes) : on ne charge que ce qui est réellement utilisé.
\usepackage[skins,breakable]{tcolorbox}
\usepackage{graphicx}
\usepackage{colortbl}
\usepackage{booktabs}
\usepackage{tabularx}
\usepackage{diagbox} % case d'en-tête diagonale des tableaux à double entrée
% Colonnes extensibles centrée / gauche / droite (C, L, R), souvent utilisées par les modèles
\newcolumntype{C}{>{\centering\arraybackslash}X}
\newcolumntype{L}{>{\raggedright\arraybackslash}X}
\newcolumntype{R}{>{\raggedleft\arraybackslash}X}
\usepackage{array}
\usepackage{multicol}
\usepackage{multirow}
\usepackage{siunitx}
% parse-numbers=false : accepte aussi \SI{2,5 \times 10^{-3}}{...}
\sisetup{locale=FR,output-decimal-marker={,},group-minimum-digits=4,parse-numbers=false}
\DeclareSIUnit\ohm{\ensuremath{\symup{\Omega}}} % U+2126 absent de la police
\DeclareSIUnit\division{div} % réglages d'oscilloscope (ms/div, V/div)
\DeclareSIUnit\tr{tr} % tours (vitesse de rotation, ex. \SI{1500}{\tr\per\minute})
\DeclareSIUnit\molar{M} % concentration molaire (ex. \SI{3}{\molar}, \SI{1}{\micro\molar})
\DeclareSIUnit\an{an} % année (le modèle confond parfois avec \year, un compteur TeX) — ex. \SI{5}{\kilo\meter\per\an}
\DeclareSIUnit\FCFA{FCFA} % franc CFA (ex. \SI{500}{\FCFA\per\kilo\gram})
\DeclareSIUnit\degreeBrix{\textdegree Brix} % teneur en sucre d'un sirop/jus (réfractométrie)
\DeclareSIUnit\month{mois} % le modèle utilise parfois \month, un compteur TeX, comme unité de durée
\usepackage{qrcode}
% \degree : commande fréquemment utilisée par les modèles pour un angle ou une
% température en dehors de \SI{}{} (non fournie par les paquets chargés).
\providecommand{\degree}{^{\circ}}
% \qed (fin de démonstration), utilisé par les modèles sans amsthm
\providecommand{\qed}{\hfill$\square$}
% \barre : double barre des valeurs interdites dans un tableau de variations (tkz-tab)
\providecommand{\barre}{\ensuremath{\|}}
% environnement proof (amsthm non chargé) utilisé par les modèles
\ifdefined\proof\else\newenvironment{proof}[1][Démonstration]{\par\noindent\textit{#1.}\ }{\qed\par}\fi
\usepackage{lastpage}
\usepackage{hyperref}
\hypersetup{hidelinks}

% ============================================================
% Palette « Pop » OnBuch+ — mêmes hex que la classe P (plus_ui.dart)
% ============================================================
\definecolor{poporange}{HTML}{F59321}
\definecolor{popdark}{HTML}{E07A0C}
\definecolor{popcream}{HTML}{FFF6E8}
\definecolor{popink}{HTML}{1C1714}
\definecolor{poppurple}{HTML}{7A5AE0}
\definecolor{popgreen}{HTML}{2E9E6A}
\definecolor{poppink}{HTML}{E0457B}
\definecolor{popblue}{HTML}{2F7DE1}
\definecolor{popgold}{HTML}{C98A12}
\definecolor{popmuted}{HTML}{8A7F76}
\definecolor{popline}{HTML}{EADFD2}
\definecolor{popgray}{HTML}{8A7F76}
\colorlet{popgrey}{popgray}
% Alias tolérants (le modèle nomme parfois les couleurs différemment)
\colorlet{popwhite}{white}
\colorlet{popblack}{popink}
\colorlet{popred}{poppink}
\colorlet{popyellow}{popgold}
% Teintes claires (fonds de tableaux/encadrés) — le modèle nomme parfois
% les couleurs avec un suffixe L (ex. popblueL) sans les avoir définies.
\colorlet{poporangeL}{poporange!20}
\colorlet{popdarkL}{popdark!20}
\colorlet{poppurpleL}{poppurple!20}
\colorlet{popgreenL}{popgreen!20}
\colorlet{poppinkL}{poppink!20}
\colorlet{popblueL}{popblue!20}
\colorlet{popgoldL}{popgold!20}
\colorlet{popredL}{popred!20}
\colorlet{popgrayL}{popgray!20}
% Teintes claires pour fonds de tableaux / zones
\colorlet{popblueL}{popblue!10!white}
\colorlet{poporangeL}{poporange!14!white}
\colorlet{popgreenL}{popgreen!12!white}
\colorlet{poppurpleL}{poppurple!12!white}
\colorlet{poppinkL}{poppink!12!white}
\colorlet{popgoldL}{popgold!14!white}

\pagecolor{popcream}

% ============================================================
% Typographie
% ============================================================
\defaultfontfeatures{Ligatures=TeX}
\setmainfont{PlusJakartaSans-Medium.ttf}[Path=../../../fonts/,BoldFont=PlusJakartaSans-Bold.ttf]
% Symboles, API et emojis absents de la police principale : repli sur DejaVu Sans (caractères actifs)
\newfontface\symfont{DejaVuSans.ttf}[Path=../../../fonts/]
\makeatletter
\newcommand\symactive[1]{\catcode#1=13 \begingroup\lccode`\~=#1 \lowercase{\endgroup\def~}{{\symfont\char#1}}}
\newcommand\symblank[1]{\catcode#1=13 \begingroup\lccode`\~=#1 \lowercase{\endgroup\def~}{}}
\newcommand\symhyph[1]{\catcode#1=13 \begingroup\lccode`\~=#1 \lowercase{\endgroup\def~}{\mbox{-}}}
\newcount\symc
\newcommand\symrange[2]{\symc=#1 \@whilenum\symc<#2 \do{\expandafter\symactive\expandafter{\the\symc}\advance\symc by 1 }}
\newcommand\symblankrange[2]{\symc=#1 \@whilenum\symc<#2 \do{\expandafter\symblank\expandafter{\the\symc}\advance\symc by 1 }}
\makeatother
\symrange{"0250}{"02FF}\symactive{"2713}\symactive{"2714}\symactive{"2717}\symactive{"2718}\symactive{"2610}\symactive{"2611}\symactive{"2612}
\symhyph{"2011}\symhyph{"2E1A}\symblankrange{"1F300}{"1FAFF}\symblank{"FE0F}
% Maths en sans-serif (Fira Math) avec chiffres et lettres droites de Plus
% Jakarta, pour que les formules se fondent dans le texte.
\usepackage[math-style=ISO,bold-style=ISO]{unicode-math}
\setmathfont{FiraMath-Regular.otf}[Path=../../../fonts/]
\setmathfont{PlusJakartaSans-Medium.ttf}[Path=../../../fonts/,range={up/num,up/{Latin,latin},\mathup}]
\usepackage{icomma} % virgule décimale sans espace en mode math
% Caractères mathématiques Unicode absents des polices de texte (Plus Jakarta,
% Archivo Black) : rendus par leur équivalent mathématique, en texte comme en
% formule — sinon ils s'affichent en carré vide (ex. « Arithmétique dans ℤ »).
\usepackage{newunicodechar}
\newunicodechar{ℝ}{\ensuremath{\mathbb{R}}}
\newunicodechar{ℤ}{\ensuremath{\mathbb{Z}}}
\newunicodechar{ℂ}{\ensuremath{\mathbb{C}}}
\newunicodechar{ℕ}{\ensuremath{\mathbb{N}}}
\newunicodechar{ℚ}{\ensuremath{\mathbb{Q}}}
\newunicodechar{𝔻}{\ensuremath{\mathbb{D}}}
\newunicodechar{α}{\ensuremath{\alpha}}
\newunicodechar{β}{\ensuremath{\beta}}
\newunicodechar{γ}{\ensuremath{\gamma}}
\newunicodechar{δ}{\ensuremath{\delta}}
\newunicodechar{ε}{\ensuremath{\varepsilon}}
\newunicodechar{θ}{\ensuremath{\theta}}
\newunicodechar{λ}{\ensuremath{\lambda}}
\newunicodechar{μ}{\ensuremath{\mu}}
\newunicodechar{σ}{\ensuremath{\sigma}}
\newunicodechar{τ}{\ensuremath{\tau}}
\newunicodechar{φ}{\ensuremath{\varphi}}
\newunicodechar{ψ}{\ensuremath{\psi}}
\newunicodechar{ω}{\ensuremath{\omega}}
\newunicodechar{Σ}{\ensuremath{\Sigma}}
% majuscules produites par \MakeUppercase dans les titres (α-aminés -> Α)
\newunicodechar{Α}{\ensuremath{\alpha}}
\newunicodechar{Β}{\ensuremath{\beta}}
\newunicodechar{Γ}{\ensuremath{\Gamma}}
\newunicodechar{Δ}{\ensuremath{\Delta}}
\newunicodechar{Ε}{\ensuremath{\varepsilon}}
\newunicodechar{Θ}{\ensuremath{\Theta}}
\newunicodechar{Λ}{\ensuremath{\Lambda}}
\newunicodechar{Μ}{\ensuremath{\mu}}
\newunicodechar{Τ}{\ensuremath{\tau}}
\newunicodechar{Φ}{\ensuremath{\Phi}}
\newunicodechar{Ψ}{\ensuremath{\Psi}}
\newunicodechar{Ω}{\ensuremath{\Omega}}
\newunicodechar{↦}{\ensuremath{\mapsto}}
\newunicodechar{∈}{\ensuremath{\in}}
\newunicodechar{∉}{\ensuremath{\notin}}
\newunicodechar{∧}{\ensuremath{\wedge}}
\newunicodechar{⇔}{\ensuremath{\Leftrightarrow}}
\newunicodechar{⟺}{\ensuremath{\Longleftrightarrow}}
\newunicodechar{⇒}{\ensuremath{\Rightarrow}}
\newunicodechar{⟹}{\ensuremath{\Longrightarrow}}
\newunicodechar{∘}{\ensuremath{\circ}}
\newunicodechar{∪}{\ensuremath{\cup}}
\newunicodechar{∩}{\ensuremath{\cap}}
\newunicodechar{≡}{\ensuremath{\equiv}}
\newunicodechar{⊂}{\ensuremath{\subset}}
\newunicodechar{⊥}{\ensuremath{\perp}}
\newunicodechar{∃}{\ensuremath{\exists}}
\newunicodechar{∀}{\ensuremath{\forall}}
\newunicodechar{∛}{\ensuremath{\sqrt[3]{\,}}}
\newunicodechar{∜}{\ensuremath{\sqrt[4]{\,}}}
\newunicodechar{⃗}{\ensuremath{{}^{\to}}}
\newunicodechar{ⁿ}{\ifmmode^{n}\else\textsuperscript{\ensuremath{n}}\fi}
\newunicodechar{ˣ}{\ifmmode^{x}\else\textsuperscript{\ensuremath{x}}\fi}
\newunicodechar{ᵇ}{\ifmmode^{b}\else\textsuperscript{\ensuremath{b}}\fi}
\newunicodechar{ʳ}{\ifmmode^{r}\else\textsuperscript{\ensuremath{r}}\fi}
\newunicodechar{ᵘ}{\ifmmode^{u}\else\textsuperscript{\ensuremath{u}}\fi}
\newunicodechar{ʸ}{\ifmmode^{y}\else\textsuperscript{\ensuremath{y}}\fi}
\newunicodechar{ᵃ}{\ifmmode^{a}\else\textsuperscript{\ensuremath{a}}\fi}
\newunicodechar{ᵉ}{\ifmmode^{e}\else\textsuperscript{\ensuremath{e}}\fi}
\newunicodechar{ᵏ}{\ifmmode^{k}\else\textsuperscript{\ensuremath{k}}\fi}
\newunicodechar{ᵖ}{\ifmmode^{p}\else\textsuperscript{\ensuremath{p}}\fi}
\newunicodechar{ᴺ}{\ifmmode^{N}\else\textsuperscript{\ensuremath{N}}\fi}
\newunicodechar{ᶜ}{\ifmmode^{c}\else\textsuperscript{\ensuremath{c}}\fi}
\newunicodechar{ʰ}{\ifmmode^{h}\else\textsuperscript{\ensuremath{h}}\fi}
\newunicodechar{ᵠ}{\ifmmode^{q}\else\textsuperscript{\ensuremath{q}}\fi}
\newunicodechar{ₙ}{\ifmmode_{n}\else\textsubscript{\ensuremath{n}}\fi}
\newunicodechar{ₐ}{\ifmmode_{a}\else\textsubscript{\ensuremath{a}}\fi}
\newunicodechar{ᵢ}{\ifmmode_{i}\else\textsubscript{\ensuremath{i}}\fi}
\newunicodechar{ᵣ}{\ifmmode_{r}\else\textsubscript{\ensuremath{r}}\fi}
\newunicodechar{ₖ}{\ifmmode_{k}\else\textsubscript{\ensuremath{k}}\fi}
\newunicodechar{ᵦ}{\ifmmode_{\beta}\else\textsubscript{\ensuremath{\beta}}\fi}
\newunicodechar{ₓ}{\ifmmode_{x}\else\textsubscript{\ensuremath{x}}\fi}
\newfontfamily\popdisplay{ArchivoBlack-Regular.ttf}[Path=../../../fonts/]
\newfontfamily\poplogo{SpaceGrotesk-Bold.ttf}[Path=../../../fonts/]
\newfontfamily\popsemi{PlusJakartaSans-SemiBold.ttf}[Path=../../../fonts/]
\newfontfamily\popxbold{PlusJakartaSans-ExtraBold.ttf}[Path=../../../fonts/]
\newfontfamily\popxboldls{PlusJakartaSans-ExtraBold.ttf}[Path=../../../fonts/,LetterSpace=3.0]

\setlist[enumerate]{leftmargin=1.7em,itemsep=5pt,topsep=4pt}
\setlist[itemize]{leftmargin=1.7em,itemsep=4pt,topsep=4pt,label={\color{poporange}\textbullet}}
\setlength{\parindent}{0pt}
\setlength{\parskip}{5pt}
\renewcommand{\arraystretch}{1.25}

% pgfplots : style Pop par défaut
\pgfplotsset{
  every axis/.append style={axis line style={popink,line width=1pt},tick style={popink},
    grid style={popline,line width=0.5pt},label style={font=\small},tick label style={font=\footnotesize}},
  popcurve/.style={poporange,line width=1.8pt,no markers,smooth},
}
% Même style disponible dans un \draw TikZ simple (hors pgfplots)
\tikzset{popcurve/.style={poporange,line width=1.8pt}}
\tikzset{popgrid/.style={popline,very thin}}
\colorlet{popcurve}{poporange} % le nom du style est parfois utilisé comme couleur

% ============================================================
% Pill / squiggle
% ============================================================
\newcommand{\poppill}[3]{%
  \tikz[baseline=-0.55ex]\node[draw=popink,line width=1.2pt,rounded corners=2.6pt,%
    fill=#2,inner xsep=6.5pt,inner ysep=3pt]{%
    \popxboldls\fontsize{7}{7}\selectfont\color{#3}\MakeUppercase{#1}};%
}
\newcommand{\popsquiggle}[1]{%
  \tikz[baseline=-0.4ex]{
    \draw[#1,line width=2.6pt,line cap=round]
      plot [smooth] coordinates {(0,0.09) (0.34,0.01) (0.68,0.21) (1.02,0.01) (1.36,0.10)};
  }%
}
% Mot-clé surligné (marqueur orange sous le texte)
\newcommand{\cle}[1]{{\popxbold\color{popdark}#1}}

% ============================================================
% En-tête / pied
% ============================================================
\pagestyle{fancy}
\fancyhf{}
\renewcommand{\headrulewidth}{0pt}
\renewcommand{\footrulewidth}{0pt}
\lhead{\raisebox{-0.15ex}{\poplogo\fontsize{13}{13}\selectfont\color{popink}OnBuch\color{poporange}+}}
\rhead{\poppill{\DOCMATIERE\ \textbullet\ \DOCNIVEAU\ \textbullet\ Leçon \DOCLECON}{white}{popink}}
\lfoot{\popxbold\fontsize{7.6}{7.6}\selectfont\color{popmuted}\MakeUppercase{Cours signé OnBuch+}\ \textbullet\ {\popsemi\DOCTITRE}}
\rfoot{\tikz[baseline=-0.6ex]\node[rounded corners=8.5pt,fill=popink,minimum height=17pt,minimum width=17pt,inner xsep=5pt,inner sep=0]{\popxbold\fontsize{8}{8}\selectfont\color{popcream}\thepage\,/\,\pageref*{LastPage}};}

% ============================================================
% Titres
% ============================================================
\newcommand{\chaptitre}[2]{%
  \begin{tcolorbox}[enhanced,colback=poporange,colframe=popink,boxrule=2.6pt,arc=11pt,
      drop shadow=popink,left=15pt,right=15pt,top=12pt,bottom=11pt,before skip=3pt,after skip=18pt]
    \poppill{#2}{popink}{popcream}\par\vspace{9pt}
    {\raggedright\hyphenpenalty=10000\exhyphenpenalty=10000\popdisplay\fontsize{22}{24}\selectfont\color{popcream}\MakeUppercase{#1}\par}
  \end{tcolorbox}
}
% Section numérotée : pastille orange avec le numéro + titre Archivo
\newcounter{popsec}
\newcommand{\coursec}[1]{%
  \stepcounter{popsec}\setcounter{popsub}{0}%
  \par\vspace{16pt}\noindent
  \begin{minipage}{\linewidth}
  \tikz[baseline=-0.7ex]\node[draw=popink,line width=1.6pt,fill=poporange,rounded corners=4pt,minimum size=22pt,inner sep=0]{\popdisplay\fontsize{12}{12}\selectfont\color{popcream}\thepopsec};%
  \hspace{8pt}{\popdisplay\fontsize{15}{17}\selectfont\color{popink}\MakeUppercase{#1}}\par
  \vspace{4pt}\noindent{\color{poporange}\rule{36pt}{3.6pt}}
  \end{minipage}\par\nopagebreak\vspace{8pt}
}
\newcounter{popsub}
\newcommand{\courssub}[1]{%
  \stepcounter{popsub}%
  \par\vspace{9pt}\noindent{\popxbold\fontsize{11.5}{13}\selectfont\color{popink}{\color{poporange}\thepopsec.\thepopsub}\hspace{6pt}#1}\par\nopagebreak\vspace{4pt}
}

% ============================================================
% Boîtes pédagogiques — même recette : fond blanc, bordure épaisse colorée,
% ombre dure encre, badge plein. Titre optionnel [#1].
% ============================================================
\tcbset{popbase/.style={breakable,enhanced,colback=white,boxrule=2.3pt,arc=9pt,
  shadow={3pt}{-3pt}{0pt}{popink},left=14pt,right=14pt,top=11pt,bottom=11pt,before skip=11pt,after skip=15pt,
  fontupper=\popsemi\fontsize{10.6}{14.5}\selectfont\color{popink},
  fontlower=\popsemi\fontsize{10.6}{14.5}\selectfont\color{popink}}}
% Coche dessinée (le glyphe ✓ n'existe pas dans les polices du document)
\newcommand{\popcheck}[1]{\tikz[baseline=-0.5ex]\draw[#1,line width=1.6pt,line cap=round,line join=round](-0.13,0.0)--(-0.04,-0.09)--(0.13,0.1);}
\providecommand{\coche}{\popcheck{popgreen}}
\providecommand{\resultat}[1]{\par\smallskip\noindent\fcolorbox{popgreen}{popgreenL}{\parbox{\dimexpr\linewidth-2\fboxsep-2\fboxrule\relax}{\textbf{Résultat.} #1}}\par\smallskip}
\newcommand{\popboxhead}[3]{\poppill{#1}{#2}{white}\ifx\relax#3\relax\else\hspace{6pt}{\popxbold\fontsize{10.6}{12}\selectfont\color{popink}#3}\fi\par\vspace{7pt}}

\newtcolorbox{aretenir}[1][]{popbase,colframe=popgreen,before upper={\popboxhead{À retenir}{popgreen}{#1}}}
% Texte en anglais (document de lecture, dialogue, extrait) : cadre
% d'étude. Un titre optionnel (source, auteur) s'affiche dans l'en-tête.
\newtcolorbox{textebox}[1][]{popbase,colframe=popblue,before upper={\popboxhead{Text}{popblue}{#1}\begin{otherlanguage}{english}},after upper={\end{otherlanguage}}}
% Marques ✓ / ✗ (forme correcte / fautive) souvent utilisées par les modèles
\providecommand{\cross}{\textcolor{poppink}{\ensuremath{\times}}}
\providecommand{\xmark}{\cross}\providecommand{\crossmark}{\cross}\providecommand{\cmark}{\ensuremath{\checkmark}}
% Ligne de réponse / justification à compléter (inventée par les modèles)
\providecommand{\justification}[1]{\par\noindent\textit{Justification~:} \dotfill\par}
% \textipa (package tipa non chargé) : la police ne couvre pas l'API -> simple passage
\providecommand{\textipa}[1]{#1}
% Soulignement (ulem non chargé) : \uline, \ul
\providecommand{\uline}[1]{\underline{#1}}\providecommand{\ul}[1]{\underline{#1}}
% \glossaire{\item ...} inventée par les modèles : liste de glossaire
\providecommand{\glossaire}[1]{\begin{itemize}#1\end{itemize}}
% \alert (beamer) inventée par les modèles : texte mis en évidence
\providecommand{\alert}[1]{\textbf{\textcolor{poppink}{#1}}}
% variantes ASCII d'ordinaux écrites en macro
\providecommand{\iere}{\textsuperscript{re}}\providecommand{\ieme}{\textsuperscript{e}}\providecommand{\ere}{\textsuperscript{re}}\providecommand{\eme}{\textsuperscript{e}}\providecommand{\er}{\textsuperscript{er}}\providecommand{\ie}{\textsuperscript{e}}
% Ordinaux français écrits en macro par les modèles : 1\ière, 3\ième
\newcommand{\ière}{\textsuperscript{re}}\newcommand{\ième}{\textsuperscript{e}}
% \exemple (inventée par les modèles) : étiquette « Exemple : »
\providecommand{\exemple}{\textbf{Exemple~:}\ }
% \square utilisable aussi hors mode math (cases à cocher)
\AtBeginDocument{\let\squareMath\square\renewcommand{\square}{\ensuremath{\squareMath}}}
% Mot ou phrase en anglais à l'intérieur d'un paragraphe français
\newcommand{\eng}[1]{\ifmmode\text{\textit{#1}}\else\begingroup\selectlanguage{english}\itshape #1\endgroup\fi}
\let\esp\eng % alias toléré
% flèches d'intonation inventées par les modèles
\providecommand{\textsearrow}{}\renewcommand{\textsearrow}{\ensuremath{\searrow}}\providecommand{\textnearrow}{}\renewcommand{\textnearrow}{\ensuremath{\nearrow}}
% \fr{..} : mot français (inventé par les modèles)
\providecommand{\fr}[1]{\textit{#1}}
\newtcolorbox{exemplebox}[1][]{popbase,colframe=poporange,before upper={\popboxhead{Exemple}{poporange}{#1}}}
\newtcolorbox{definition}[1][]{popbase,colframe=popblue,before upper={\popboxhead{Définition}{popblue}{#1}}}
\newtcolorbox{propriete}[1][]{popbase,colframe=poppurple,before upper={\popboxhead{Propriété}{poppurple}{#1}}}
\newtcolorbox{methode}[1][]{popbase,colframe=popgold,before upper={\popboxhead{Méthode}{popgold}{#1}}}
\newtcolorbox{attention}[1][]{popbase,colframe=poppink,before upper={\popboxhead{Attention}{poppink}{#1}}}
\newtcolorbox{experience}[1][]{popbase,colframe=popblue,colback=popblueL,before upper={\popboxhead{Expérience}{popblue}{#1}}}
% « Le savais-tu ? » : carte violette pleine (contexte, culture, Cameroun)
\newtcolorbox{savaistu}[1][]{popbase,colback=poppurple,colframe=popink,
  fontupper=\popsemi\fontsize{10.4}{14.2}\selectfont\color{white},
  before upper={\poppill{Le savais-tu\,?}{white}{poppurple}\ifx\relax#1\relax\else\hspace{6pt}{\popxbold\color{white}#1}\fi\par\vspace{7pt}}}
% Exercice résolu : énoncé en haut, \tcblower puis la solution sur fond crème
\newcounter{popexo}\newcounter{popexr}
\newtcolorbox{exoresolu}[1][]{popbase,colframe=poporange,colbacklower=popcream,
  segmentation style={popink,line width=1.2pt,dashed},
  before upper={\stepcounter{popexr}\popboxhead{Exercice résolu \thepopexr}{poporange}{#1}},
  before lower={\poppill{Solution}{popink}{popcream}\par\vspace{6pt}}}
% Exercice (non résolu en place) — la correction est dans la partie Corrigés
\newtcolorbox{exercice}[1][]{popbase,colframe=popink,
  before upper={\stepcounter{popexo}\popboxhead{Exercice \thepopexo}{popink}{#1}}}
% Corrigé
\newtcolorbox{corrige}[1][]{popbase,colframe=popgreen,colback=popgreenL,
  before upper={\popboxhead{Corrigé}{popgreen}{#1}}}
% Objectifs / prérequis / bilan
\newtcolorbox{objectifs}{popbase,colframe=popink,colback=poporangeL,
  before upper={\popboxhead{Objectifs de la leçon}{popink}{}}}
\newtcolorbox{prerequis}{popbase,colframe=popink,colback=popblueL,
  before upper={\popboxhead{Prérequis}{popblue}{}}}
\newtcolorbox{bilan}{popbase,colframe=popink,colback=white,boxrule=2.8pt,
  before upper={\popboxhead{Fiche bilan}{poporange}{L'essentiel à connaître pour l'examen}}}
% Figure encadrée avec légende
\newtcolorbox{popfigure}[1][]{enhanced,breakable=false,colback=white,colframe=popline,boxrule=1.4pt,arc=8pt,
  left=10pt,right=10pt,top=10pt,bottom=8pt,before skip=10pt,after skip=12pt,halign=center,
  lower separated=false,halign lower=center,
  fontlower=\popsemi\fontsize{9}{11.5}\selectfont\color{popmuted},
  #1}
\newcommand{\legende}[1]{\par\vspace{6pt}{\popsemi\fontsize{9}{11.5}\selectfont\color{popmuted}#1}}

% ============================================================
% Signature OnBuch+
% ============================================================
\newcommand{\poplogoinline}[1]{{\poplogo\fontsize{#1}{#1}\selectfont\color{popink}OnBuch\color{poporange}+}}
\newcommand{\signatureonbuch}{%
  \begin{tcolorbox}[enhanced,colback=popink,colframe=popink,boxrule=2.2pt,arc=10pt,
      drop shadow=poporange,left=14pt,right=14pt,top=10pt,bottom=10pt,before skip=0pt,after skip=0pt]
    \tikz[baseline=-0.7ex]\node[circle,fill=popgreen,draw=popcream,line width=1.3pt,minimum size=22pt,inner sep=0]{\popcheck{white}};%
    \hspace{9pt}\begin{minipage}[c]{0.62\linewidth}
      {\popxbold\fontsize{12}{13}\selectfont\color{popcream}Signé\ \textbullet\ {\poplogo OnBuch\color{poporange}+}}\par\vspace{2pt}
      {\raggedright\popsemi\fontsize{8}{10.5}\selectfont\color{popline}\MakeUppercase{Équipe pédagogique OnBuch+}\\\MakeUppercase{Conforme au programme officiel MINESEC}\par}
    \end{minipage}\hfill
    {\poplogo\fontsize{22}{22}\selectfont\color{popcream}OnBuch\color{poporange}+}
  \end{tcolorbox}%
}

% ============================================================
% Page de garde
% #1 titre · #2 matière · #3 niveau · #4 module · #5 n° leçon · #6 durée ·
% #7 description / objectifs (texte ou itemize)
% ============================================================
\newcommand{\pagegarde}[7]{%
  \thispagestyle{empty}
  \newgeometry{a4paper,margin=1.7cm,top=1.6cm,bottom=1.4cm}%
  \begin{tikzpicture}[remember picture,overlay]
    \fill[poporange!18] ([xshift=-3.2cm,yshift=-3.6cm]current page.north east) circle (4.6cm);
    \fill[poppurple!14] ([xshift=2.4cm,yshift=7.6cm]current page.south west) circle (3.8cm);
  \end{tikzpicture}%
  {\poplogo\fontsize{30}{30}\selectfont\color{popink}OnBuch\color{poporange}+}\hfill
  \raisebox{0.6ex}{\poppill{Cours signé \textbullet\ Premium}{popink}{popcream}}\par
  \vspace{1pt}\popsquiggle{poppurple}\par
  \vspace{8pt}
  \begin{tcolorbox}[enhanced,colback=poporange,colframe=popink,boxrule=2.8pt,arc=13pt,
      drop shadow=popink,left=18pt,right=18pt,top=12pt,bottom=13pt,after skip=10pt]
    \poppill{#2 \textbullet\ #3}{popink}{popcream}\hspace{5pt}\poppill{#4}{white}{popink}\par\vspace{8pt}
    {\popdisplay\fontsize{14}{14}\selectfont\color{popink}LEÇON #5}\par\vspace{4pt}
    {\raggedright\hyphenpenalty=10000\exhyphenpenalty=10000\popdisplay\fontsize{25}{27}\selectfont\color{popcream}\MakeUppercase{#1}\par}
    \vspace{8pt}
    {\popxbold\fontsize{9.5}{9.5}\selectfont\color{popink}\MakeUppercase{Durée conseillée : #6}}
  \end{tcolorbox}
  \begin{tcolorbox}[enhanced,colback=white,colframe=popink,boxrule=2.2pt,arc=10pt,
      drop shadow=popink,left=16pt,right=16pt,top=10pt,bottom=10pt,after skip=8pt]
    \poppill{Dans cette leçon}{poppurple}{white}\par\vspace{5pt}
    \popsemi\fontsize{9.3}{12.6}\selectfont\color{popink}#7
  \end{tcolorbox}
  \noindent\begin{minipage}[t]{0.487\linewidth}
    \begin{tcolorbox}[enhanced,colback=popgreen,colframe=popink,boxrule=2.2pt,arc=10pt,
        drop shadow=popink,left=11pt,right=11pt,top=10pt,bottom=10pt,height=3.1cm,valign=center]
      \tcbox[colback=white,colframe=popink,boxrule=1.3pt,arc=5pt,boxsep=0pt,nobeforeafter,
        left=3pt,right=3pt,top=3pt,bottom=3pt,valign=center]{\qrcode[height=1.9cm]{https://play.google.com/store/apps/details?id=com.luvvix.onbuch&hl=fr}}%
      \hspace{8pt}\begin{minipage}[c]{0.5\linewidth}
        {\popdisplay\fontsize{11}{12.5}\selectfont\color{white}\MakeUppercase{Télécharge OnBuch}}\par\vspace{4pt}
        {\raggedright\popsemi\fontsize{8.4}{11}\selectfont\color{white}Gratuit \textbullet\ Google Play\\Tuteur IA, annales, concours, résultats\par}
      \end{minipage}
    \end{tcolorbox}
  \end{minipage}\hfill
  \begin{minipage}[t]{0.487\linewidth}
    \begin{tcolorbox}[enhanced,colback=poppurple,colframe=popink,boxrule=2.2pt,arc=10pt,
        drop shadow=popink,left=11pt,right=11pt,top=10pt,bottom=10pt,height=3.1cm,valign=center]
      \tikz[baseline=-0.7ex]\node[circle,fill=white,draw=popink,line width=1.3pt,minimum size=1.6cm,inner sep=0]{\popdisplay\fontsize{24}{24}\selectfont\color{poppurple}+};%
      \hspace{8pt}\begin{minipage}[c]{0.6\linewidth}
        {\popdisplay\fontsize{11}{12.5}\selectfont\color{white}\MakeUppercase{Passe à OnBuch+}}\par\vspace{4pt}
        {\raggedright\popsemi\fontsize{8.4}{11}\selectfont\color{white}Tous les cours signés, Tuteur IA illimité, fiches PDF — onglet OnBuch+ de l'app\par}
      \end{minipage}
    \end{tcolorbox}
  \end{minipage}
  \vspace{8pt}
  \popcontenu
  \vfill
  \signatureonbuch
  \restoregeometry
  \newpage
}

% Rangée « Ce cours contient » (page de garde)
\newcommand{\poptile}[3]{% #1 couleur #2 gros texte #3 légende
  \begin{tcolorbox}[enhanced,colback=#1,colframe=popink,boxrule=2pt,arc=9pt,shadow={2.5pt}{-2.5pt}{0pt}{popink},
      left=6pt,right=6pt,top=9pt,bottom=9pt,halign=center,valign=center,height=1.75cm,width=\linewidth,nobeforeafter]
    {\popdisplay\fontsize{12.5}{13}\selectfont\color{white}\MakeUppercase{#2}}\par\vspace{3pt}
    {\centering\popsemi\fontsize{7.8}{9.5}\selectfont\color{white}#3\par}
  \end{tcolorbox}}
\newcommand{\popcontenu}{%
  \par\noindent{\popxboldls\fontsize{7.5}{7.5}\selectfont\color{popmuted}CE COURS SIGNÉ CONTIENT}\par\vspace{7pt}
  \noindent\begin{minipage}[t]{0.235\linewidth}\poptile{popblue}{Cours}{complet, illustré, pas à pas}\end{minipage}\hfill
  \begin{minipage}[t]{0.235\linewidth}\poptile{poporange}{Méthodes}{et exercices résolus}\end{minipage}\hfill
  \begin{minipage}[t]{0.235\linewidth}\poptile{poppink}{Exercices}{type examen + corrigés}\end{minipage}\hfill
  \begin{minipage}[t]{0.235\linewidth}\poptile{popgreen}{Fiche bilan}{l'essentiel à retenir}\end{minipage}\par}

% Sommaire
\newtcolorbox{sommaire}{enhanced,colback=white,colframe=popink,boxrule=2.2pt,arc=10pt,
  drop shadow=popink,left=18pt,right=18pt,top=13pt,bottom=13pt,before skip=6pt,after skip=20pt,
  before upper={\poppill{Sommaire}{poppurple}{white}\par\vspace{9pt}\popsemi\fontsize{10.6}{14.5}\selectfont\color{popink}}}

% Signature de fin de cours — poussée tout en bas de la dernière page
\newcommand{\findecours}{%
  \par\vfill\nopagebreak
  \begin{center}\popsquiggle{poporange}\end{center}\vspace{4pt}
  \signatureonbuch
}
\AtBeginDocument{\def\llbracket{\mathopen{[\mkern-3.5mu[}}\def\rrbracket{\mathclose{]\mkern-3.5mu]}}} % intervalles d'entiers (glyphe ⟦ absent de la police)
% Glyphes absents de Fira Math : redéfinis à partir de glyphes disponibles
\AtBeginDocument{%
  \def\setminus{\mathbin{\backslash}}\let\smallsetminus\setminus
  \def\vdots{\mathord{\vbox{\baselineskip3.2pt\lineskiplimit0pt\kern4pt\hbox{.}\hbox{.}\hbox{.}}}}%
  \def\ddots{\mathinner{\mkern1mu\raise6.5pt\hbox{.}\mkern2mu\raise3.6pt\hbox{.}\mkern2mu\raise0.7pt\hbox{.}\mkern1mu}}%
  \def\triangle{\mathord{\scalebox{0.9}{$\symup{\Delta}$}}}%
  \def\nabla{\mathord{\raisebox{\depth}{\scalebox{1}[-1]{$\symup{\Delta}$}}}}%
  \def\checkmark{\popcheck{popdark}}%
  \def\star{\mathbin{\ast}}\def\bigstar{\mathord{\ast}}%
  \def\diamond{\mathbin{\scalebox{0.6}{$\Diamond$}}}%
  \def\bigcup{\mathop{\mathchoice{\raisebox{-0.3em}{\scalebox{1.7}{$\cup$}}}{\scalebox{1.25}{$\cup$}}{\cup}{\cup}}\displaylimits}%
  \def\bigcap{\mathop{\mathchoice{\raisebox{-0.3em}{\scalebox{1.7}{$\cap$}}}{\scalebox{1.25}{$\cap$}}{\cap}{\cap}}\displaylimits}%
  \def\lesseqqgtr{\mathrel{\lessgtr}}\def\bigoplus{\mathop{\oplus}\displaylimits}%
}
\providecommand{\ssi}{si et seulement si} % abréviation fréquente
\AtBeginDocument{\providecommand{\wideparen}{\overparen}} % arc de cercle

%% FIGFIX-BEGIN v5 — correctifs globaux des figures (docs/onbuch-generation/tools/figfix)
\usepackage{adjustbox}\usepackage{etoolbox}\tcbuselibrary{raster}
% toute figure TikZ/pgfplots plus large que la ligne est réduite à la largeur disponible
\BeforeBeginEnvironment{tikzpicture}{\begin{adjustbox}{max width=\linewidth}}
\AfterEndEnvironment{tikzpicture}{\end{adjustbox}}
% légendes pgfplots lisibles par-dessus les courbes
\pgfplotsset{every axis/.append style={legend style={fill=white,fill opacity=0.92,text opacity=1,draw=black!15,inner sep=3pt}}}
% étiquettes d'arêtes (syntaxe edge["..."]) avec fond blanc
\tikzset{every edge quotes/.append style={fill=white,inner sep=1.2pt}}
%% FIGFIX-END
\begin{document}

\pagegarde{Lesson 39 — Vocabulary: Words and expressions related to volunteering in gender equality promotion activities}{Anglais}{Tle A}{Chapter 4 : Citizenship and Human Rights}{EN39}{2 h}{Cette leçon de vocabulaire présente les mots, collocations et expressions liés au volontariat dans les activités de promotion de l'égalité des sexes. Elle permet aux élèves de Terminale A d'identifier, de prononcer et de réemployer ce lexique à l'oral comme à l'écrit, en évitant les faux amis et les calques du français.
\vspace{4pt}
\begin{multicols}{2}\small
\begin{itemize}
\item À la fin de cette leçon, je sais identifier et définir le vocabulaire clé du volontariat (volunteer, commitment, community service, outreach...).
\item À la fin de cette leçon, je sais employer le lexique de l'égalité des sexes (gender equality, empowerment, discrimination, advocacy...).
\item À la fin de cette leçon, je sais utiliser les collocations fréquentes du thème (to raise awareness, to empower women, to promote gender equality...).
\item À la fin de cette leçon, je sais distinguer les faux amis et éviter les calques du français (to assist ≠ assister à, sensible ≠ sensible...).
\item À la fin de cette leçon, je sais former des mots de la même famille (equal → equality → unequal → inequality ; volunteer → voluntary → voluntarily...).
\item À la fin de cette leçon, je sais prononcer correctement les mots difficiles du thème (volunteer, equality, discrimination, advocacy).
\end{itemize}\end{multicols}}

\begin{sommaire}
\begin{multicols}{2}
\begin{enumerate}
\item Champ lexical 1 : Volunteering and community service
\item Champ lexical 2 : Gender equality and women's rights
\item Collocations et expressions idiomatiques
\item Faux amis et pièges des francophones
\item Mots de la même famille, prononciation et orthographe
\item Le lexique en contexte : lecture et production guidée
\item Activité d'intégration
\item Méthodes et exercices résolus
\item Exercices
\item Corrigés des exercices
\item Fiche bilan
\end{enumerate}
\end{multicols}
\end{sommaire}

\begin{objectifs}
\begin{itemize}
\item À la fin de cette leçon, je sais identifier et définir le vocabulaire clé du volontariat (volunteer, commitment, community service, outreach...).
\item À la fin de cette leçon, je sais employer le lexique de l'égalité des sexes (gender equality, empowerment, discrimination, advocacy...).
\item À la fin de cette leçon, je sais utiliser les collocations fréquentes du thème (to raise awareness, to empower women, to promote gender equality...).
\item À la fin de cette leçon, je sais distinguer les faux amis et éviter les calques du français (to assist ≠ assister à, sensible ≠ sensible...).
\item À la fin de cette leçon, je sais former des mots de la même famille (equal → equality → unequal → inequality ; volunteer → voluntary → voluntarily...).
\item À la fin de cette leçon, je sais prononcer correctement les mots difficiles du thème (volunteer, equality, discrimination, advocacy).
\item À la fin de cette leçon, je sais réemployer ce lexique dans des phrases, des dialogues et un court texte sur le volontariat.
\item À la fin de cette leçon, je sais rédiger un court paragraphe ou une lettre de motivation de volontaire en utilisant le lexique étudié.
\end{itemize}
\end{objectifs}

\begin{prerequis}
\begin{itemize}
\item Vocabulaire général des droits de l'homme vu en Première (rights, freedom, justice, respect).
\item Notions sur les organisations et la vie associative (association, NGO, member, project).
\item Formation des mots : préfixes et suffixes courants (-tion, -ment, -ity, un-, dis-).
\item Expression de l'opinion et de l'accord (I think, I believe, I agree) pour les activités orales.
\item Structures de base de la lettre formelle (formules d'appel et de clôture).
\end{itemize}
\end{prerequis}

\newpage

\coursec{Champ lexical 1 : Volunteering and community service}

Every year in Bamenda, Yaoundé and Buea, hundreds of young Cameroonians join NGOs like UN Women volunteers or local associations to promote girls' education and fight gender-based violence. Imagine your school is launching a “Gender Equality Club” and a British partner organisation is recruiting English-speaking volunteers. To apply, you must understand and use the right vocabulary: words like \eng{empowerment}, \eng{advocacy} or \eng{outreach}. This lesson equips you with the exact words and expressions you need to talk about volunteering for gender equality — in class, at the Bac, and in real life.

\courssub{Les acteurs du volontariat : noms et désignations}

Commençons par observer un court témoignage authentique.

\begin{textebox}{Témoignage d'une volontaire à Bamenda}
“Last year, I joined a local NGO as a \textbf{volunteer}. Every Saturday, we organise \textbf{workshops} in the \textbf{community centre} to teach young girls about their rights. It is unpaid \textbf{voluntary work}, but the \textbf{commitment} of the team is amazing. Our \textbf{coordinator} says that \textbf{activism} starts with small actions.”
\end{textebox}

\begin{definition}[Les acteurs clés]
Un \textbf{volunteer} (nom) est une personne qui travaille de son plein gré, sans être payée, pour aider les autres. L'organisation qui l'accueille est souvent une \textbf{NGO} (Non-Governmental Organisation / Organisation non gouvernementale) ou une \textbf{charity} (association caritative). Une personne très engagée dans une cause est un \textbf{activist} (militant). Celui qui encadre les volontaires est un \textbf{coordinator} (coordinateur) ou un \textbf{supervisor} (superviseur).
\end{definition}

\begin{center}
\begin{tabularx}{\linewidth}{|X|X|X|}
\hline
\rowcolor{popblueL} \textbf{English} & \textbf{Français} & \textbf{Classe / Note} \\
\hline
a \textbf{volunteer} & un volontaire, un bénévole & Nom (personne) \\
\textbf{volunteering} & le volontariat, le bénévolat & Nom d'action (gérondif) \\
\textbf{voluntary work} / \textbf{community service} & le travail bénévole / le service communautaire & Locutions nominales \\
an \textbf{NGO} (Non-Governmental Organisation) & une ONG & Sigle (se prononce \eng{en-dji-ou}) \\
a \textbf{charity} & une association caritative, une œuvre de charité & Nom \\
an \textbf{activist} & un militant, une militante & Nom (personne engagée) \\
a \textbf{helper} & un aide, un assistant & Nom (terme plus familier) \\
a \textbf{coordinator} & un coordinateur, une coordinatrice & Nom (encadrement) \\
\hline
\end{tabularx}
\end{center}

\courssub{Les actions du volontariat : verbes et collocations}

En anglais, les actions s'expriment par des verbes précis et des \textbf{collocations} (associations de mots figées) qu'il faut mémoriser par blocs.

\begin{propriete}[Verbes et collocations du volontariat]
\begin{itemize}
    \item \textbf{to volunteer} (verbe intransitif) : être volontaire, faire du bénévolat. \eng{She volunteers at the local clinic.} (Elle fait du bénévolat à la clinique locale.)
    \item \textbf{to volunteer for} + nom / \textbf{to volunteer to} + infinitif : se porter volontaire pour. \eng{He volunteered for the outreach programme.} / \eng{She volunteered to teach.}
    \item \textbf{to do voluntary work} / \textbf{to do community service} : faire du travail bénévole / effectuer un service communautaire. \textbf{Attention :} on ne dit pas \eng{*to make volunteering}.
    \item \textbf{to join an organisation} : rejoindre / adhérer à une organisation.
    \item \textbf{to take part in} / \textbf{to participate in} : participer à (un projet, une campagne).
    \item \textbf{to help out} (particule \emph{out} = aide ponctuelle/dépannage) : donner un coup de main, aider. \eng{Can you help out this Saturday?}
    \item \textbf{to give one's time} : donner de son temps.
    \item \textbf{to serve the community} : servir la communauté.
    \item \textbf{to empower} : autonomiser, donner du pouvoir à (terme clé pour l'égalité des genres). \eng{The project empowers rural women.}
    \item \textbf{to advocate for} : plaider pour, défendre (une cause). \eng{They advocate for girls' education.}
    \item \textbf{to raise awareness (about)} : sensibiliser (à).
\end{itemize}
\end{propriete}

\begin{exemplebox}[Collocations en contexte]
\begin{enumerate}
    \item \eng{Many students \textbf{volunteer for} the annual \textbf{outreach programme}.} (Beaucoup d'élèves se portent volontaires pour le programme de sensibilisation annuel.)
    \item \eng{We \textbf{help out} at the \textbf{community centre} during holidays.} (Nous donnons un coup de main au centre communautaire pendant les vacances.)
    \item \eng{The NGO \textbf{empowers} women through \textbf{training sessions}.} (L'ONG autonomise les femmes grâce à des sessions de formation.)
    \item \eng{Activists \textbf{advocate for} gender equality.} (Les militants plaident pour l'égalité des sexes.)
\end{enumerate}
\end{exemplebox}

\courssub{Les lieux et cadres d'intervention}

\begin{center}
\begin{tabularx}{\linewidth}{|X|X|X|}
\hline
\rowcolor{popgreenL} \textbf{English} & \textbf{Français} & \textbf{Contexte d'usage} \\
\hline
a \textbf{community centre} & un centre communautaire & Lieu physique principal (quartier, village) \\
\textbf{field work} & le travail de terrain & Actions hors des bureaux, dans les villages \\
an \textbf{outreach programme} & un programme de sensibilisation / de proximité & Aller vers les populations (écoles, marchés) \\
a \textbf{campaign} & une campagne & Action limitée dans le temps (ex: \eng{16 Days of Activism}) \\
a \textbf{workshop} & un atelier & Formation pratique, interactive \\
a \textbf{training session} & une session de formation & Apprentissage de compétences spécifiques \\
a \textbf{meeting} / a \textbf{gathering} & une réunion / un rassemblement & Échange d'informations, mobilisation \\
\hline
\end{tabularx}
\end{center}

\courssub{Les qualités du volontaire : lexique des valeurs}

Ce vocabulaire est indispensable pour rédiger une lettre de motivation (\eng{cover letter / motivation letter}) ou un CV (\eng{résumé / CV}) pour un poste de volontaire.

\begin{center}
\begin{tabularx}{\linewidth}{|X|X|X|}
\hline
\rowcolor{poppurpleL} \textbf{Nom (Qualité)} & \textbf{Adjectif correspondant} & \textbf{Traduction / Nuance} \\
\hline
\textbf{commitment} & \textbf{committed} & engagement / dévoué (fermeté, sérieux) \\
\textbf{dedication} & \textbf{dedicated} & dévouement / dévoué (passion, don de soi) \\
\textbf{selflessness} & \textbf{selfless} & abnégation, altruisme / désintéressé \\
\textbf{team spirit} & — & esprit d'équipe (invariable) \\
\textbf{motivation} & \textbf{motivated} & motivation / motivé \\
\textbf{reliability} & \textbf{reliable} & fiabilité / fiable (on peut compter sur lui) \\
\textbf{empathy} & \textbf{empathetic} & empathie / empathique \\
\textbf{leadership} & — & leadership, capacité à diriger \\
\hline
\end{tabularx}
\end{center}

\begin{exemplebox}[Qualités en phrases]
\eng{Volunteers must show strong \textbf{commitment} and \textbf{reliability}.} (Les volontaires doivent faire preuve d'un fort engagement et de fiabilité.) \\
\eng{Her \textbf{dedication} to the cause is inspiring.} (Son dévouement à la cause est inspirant.) \\
\eng{We need \textbf{empathetic} people for this \textbf{outreach programme}.} (Nous avons besoin de personnes empathiques pour ce programme de proximité.)
\end{exemplebox}

\courssub{Distinction cruciale : \texorpdfstring{\cle{volunteer}}{volunteer} / \texorpdfstring{\cle{voluntary}}{voluntary} / \texorpdfstring{\cle{voluntarily}}{voluntarily}}

C'est le point de grammaire lexicale le plus fréquent aux examens (Bac, Probatoire) et la source n°1 d'erreurs pour les francophones.

\begin{propriete}[Formes de la même famille : Nom / Adjectif / Adverbe]
\begin{center}
\begin{tabular}{|l|l|l|l|}
\hline
\rowcolor{poporangeL} \textbf{Mot} & \textbf{Classe grammaticale} & \textbf{Prononciation (approx. FR)} & \textbf{Traduction} \\
\hline
\textbf{volunteer} & Nom (personne) / Verbe & « volon-\textbf{TIRE} » (accent dernière syllabe) & un volontaire / faire du bénévolat \\
\textbf{voluntary} & Adjectif & « volon-\textbf{TE-ri} » (accent sur l'antépénultième) & bénévole, volontaire (non payé / libre) \\
\textbf{voluntarily} & Adverbe & « volon-\textbf{TE-ri-li} » & bénévolement, volontairement \\
\hline
\end{tabular}
\end{center}

\end{propriete}

\begin{attention}[Pièges majeurs pour francophones]
\begin{itemize}
    \item \textbf{Erreur :} \eng{*He is a voluntary.} \rightarrow \textbf{Correct :} \eng{He is a \textbf{volunteer}.} (\eng{voluntary} est un adjectif, jamais un nom pour une personne).
    \item \textbf{Erreur :} \eng{*She does voluntary.} \rightarrow \textbf{Correct :} \eng{She does \textbf{voluntary work}.} ou \eng{She \textbf{volunteers}.}
    \item \textbf{Erreur :} \eng{*He worked voluntary.} \rightarrow \textbf{Correct :} \eng{He worked \textbf{voluntarily}.} (Adverbe en \textit{-ly} pour modifier le verbe).
    \item \textbf{Faux ami :} \eng{Voluntary} ne signifie pas « volontaire » au sens de « intentionnel » seulement (comme \fr{volontaire} en français), mais surtout « non rémunéré, fait par choix ». \eng{Voluntary redundancy} = départ volontaire (licenciement accepté).
\end{itemize}
\end{attention}

\begin{exemplebox}[Les trois formes en contexte]
\begin{tabular}{ll}
\eng{1. A \textbf{volunteer} helps the community.} & (Un \textbf{bénévole} aide la communauté.) — \textit{Nom sujet} \\
\eng{2. I want to \textbf{volunteer} this summer.} & (Je veux \textbf{faire du bénévolat} cet été.) — \textit{Verbe} \\
\eng{3. It is a \textbf{voluntary} organisation.} & (C'est une organisation \textbf{bénévole} / à but non lucratif.) — \textit{Adjectif} \\
\eng{4. She \textbf{voluntarily} gave her time.} & (Elle a \textbf{bénévolement / volontairement} donné son temps.) — \textit{Adverbe} \\
\end{tabular}
\end{exemplebox}

\courssub{Prononciation : points de vigilance}

Le stress (accent tonique) change le sens ou la classe grammaticale. En anglais, les mots de 3 syllabes ou plus voient souvent leur accent se déplacer selon le suffixe.

\begin{center}
\begin{tabular}{|l|l|l|}
\hline
\rowcolor{poppinkL} \textbf{Mot} & \textbf{Transcription simplifiée (FR)} & \textbf{Règle / Astuce} \\
\hline
\textbf{volunteer} (N/V) & « volon-\textbf{TIRE} » & Accent sur la \textbf{dernière} syllabe (son fort ). Pensez au verbe \eng{engin-\textbf{EER}}. \\
\textbf{voluntary} (Adj) & « volon-\textbf{TE-ri} » & Accent sur l'\textbf{antépénultième} (3e à partir de la fin) à cause du suffixe \textit{-ary}. \\
\textbf{voluntarily} (Adv) & « volon-\textbf{TE-ri-li} » & Même accent que l'adjectif + \textit{-ly}. \\
\textbf{charity} & « \textbf{TCHA}-ri-ti » & Accent sur la 1ère syllabe. \textit{ch} = /tʃ/ (tch). Pas « cha-ri-té » à la française. \\
\textbf{commitment} & « ko-\textbf{MIT}-ment » & Accent sur la 2ème syllabe. Double \textit{mm} = son \textit{m} court. Ne pas dire « com-mi-tment ». \\
\textbf{NGO} & « \textbf{EN}-\textbf{DI}-\textbf{OU} » & Épeler les lettres. Pas « éngo ». \\
\textbf{activist} & « \textbf{AC}-ti-vist » & Accent sur la 1ère syllabe. \textit{ti} = /tɪ/ (ti), pas /si/. \\
\hline
\end{tabular}
\end{center}

\begin{methode}[Comment bien prononcer \eng{volunteer} et sa famille]
\begin{enumerate}
    \item \textbf{Isoler le son fort :} \eng{vol-un-TEER} → tapez dans les mains sur \textbf{TEER}.
    \item \textbf{Chaîner :} \eng{A volunteer} (liaison \textit{r} final + voyelle) → « a-vo-lun-ti-reur ».
    \item \textbf{Transformer :} Dites \eng{voluntary} en gardant la base \eng{volun-} mais en déplaçant le coup de poing sur \textbf{TARY}.
    \item \textbf{S'enregistrer :} Comparez avec un modèle natif (Forvo, Cambridge Dictionary en ligne).
\end{enumerate}
\end{methode}

\courssub{Carte mentale récapitulative : VOLUNTEERING}

\begin{popfigure}
\begin{center}\tcbox[colback=popink,colframe=popink,coltext=white,arc=9pt,boxsep=4pt,left=6pt,right=6pt]{\bfseries\small VOLUNTEERING}\end{center}
\vspace{-2pt}
\begin{tcbraster}[raster columns=2,raster equal height=rows,raster column skip=6pt,raster row skip=6pt,enhanced,blankest]
\begin{tcolorbox}[enhanced,colback=white,colframe=popgreen,boxrule=0.9pt,arc=3pt,title={PEOPLE},fonttitle=\bfseries\scriptsize,coltitle=white,colbacktitle=popgreen,left=4pt,right=4pt,top=2pt,bottom=2pt,boxsep=1pt,toptitle=1.5pt,bottomtitle=1.5pt,halign=left]
\begin{itemize}[nosep,leftmargin=*,itemsep=1pt,font=\scriptsize]
\item volunteer
\item activist
\item coordinator
\item helper
\end{itemize}
\end{tcolorbox}
\begin{tcolorbox}[enhanced,colback=white,colframe=poporange,boxrule=0.9pt,arc=3pt,title={ACTIONS},fonttitle=\bfseries\scriptsize,coltitle=white,colbacktitle=poporange,left=4pt,right=4pt,top=2pt,bottom=2pt,boxsep=1pt,toptitle=1.5pt,bottomtitle=1.5pt,halign=left]
\begin{itemize}[nosep,leftmargin=*,itemsep=1pt,font=\scriptsize]
\item to volunteer
\item to help out
\item to serve
\item to empower
\item to advocate for
\end{itemize}
\end{tcolorbox}
\begin{tcolorbox}[enhanced,colback=white,colframe=poppurple,boxrule=0.9pt,arc=3pt,title={PLACES},fonttitle=\bfseries\scriptsize,coltitle=white,colbacktitle=poppurple,left=4pt,right=4pt,top=2pt,bottom=2pt,boxsep=1pt,toptitle=1.5pt,bottomtitle=1.5pt,halign=left]
\begin{itemize}[nosep,leftmargin=*,itemsep=1pt,font=\scriptsize]
\item community centre
\item field work
\item outreach programme
\item workshop
\end{itemize}
\end{tcolorbox}
\begin{tcolorbox}[enhanced,colback=white,colframe=poppink,boxrule=0.9pt,arc=3pt,title={QUALITIES},fonttitle=\bfseries\scriptsize,coltitle=white,colbacktitle=poppink,left=4pt,right=4pt,top=2pt,bottom=2pt,boxsep=1pt,toptitle=1.5pt,bottomtitle=1.5pt,halign=left]
\begin{itemize}[nosep,leftmargin=*,itemsep=1pt,font=\scriptsize]
\item commitment
\item dedication
\item reliability
\item empathy
\end{itemize}
\end{tcolorbox}
\end{tcbraster}

\legende{Carte mentale du champ lexical « Volunteering and community service » — À reproduire dans votre cahier de vocabulaire.}
\end{popfigure}

\courssub{Le saviez-vous ? : La culture du volontariat au Royaume-Uni}

\begin{savaistu}[Volunteering \& Employability au UK]
Au Royaume-Uni, le \eng{volunteering} est hautement valorisé sur les \eng{CVs} (Curriculum Vitae) et dans les dossiers \eng{UCAS} (admission universitaire). Les employeurs et universités recherchent des preuves de \eng{community service} ou de \eng{voluntary work} car ils démontrent \eng{soft skills} (compétences transverses) : \eng{teamwork}, \eng{leadership}, \eng{time management}, \eng{empathy}. Le \eng{National Citizen Service (NCS)} est un programme gouvernemental pour les 16-17 ans incluant du \eng{voluntary work}. Au Cameroun, le \eng{Service Civique National} ou l'engagement dans les \eng{Clubs Genre} des lycées sont les équivalents locaux à mettre en avant dans vos lettres de motivation pour les bourses (ex. \eng{Chevening}, \eng{Mastercard Foundation}).
\end{savaistu}

\courssub{Exercice résolu : Mettre en pratique}

\begin{exoresolu}[Choix du mot juste et transformation]
\textbf{Énoncé :} Complétez les phrases avec la forme correcte du mot entre parenthèses (\eng{volunteer}, \eng{voluntary}, \eng{voluntarily}). Traduisez ensuite.

1. Last year, my sister decided to \_\_\_\_\_\_\_\_\_\_\_\_ for a local NGO in Douala. (\textbf{volunteer}) \\
2. The association relies entirely on \_\_\_\_\_\_\_\_\_\_\_\_ contributions. (\textbf{voluntary}) \\
3. He \_\_\_\_\_\_\_\_\_\_\_\_ offered to drive the children to school. (\textbf{voluntarily}) \\
4. We need more \_\_\_\_\_\_\_\_\_\_\_\_ for the \_\_\_\_\_\_\_\_\_\_\_\_ campaign next week. (\textbf{volunteer} x2) \\
5. \_\_\_\_\_\_\_\_\_\_\_\_ work looks great on a university application. (\textbf{voluntary})

\tcblower
\textbf{Solution détaillée :}

1. \textbf{to volunteer} (verbe après \eng{decided to}). \eng{Last year, my sister decided to \textbf{volunteer} for a local NGO in Douala.} $\rightarrow$ « L'an dernier, ma sœur a décidé de faire du bénévolat pour une ONG locale à Douala. »
2. \textbf{voluntary} (adjectif qualificatif devant le nom \eng{contributions}). \eng{The association relies entirely on \textbf{voluntary} contributions.} $\rightarrow$ « L'association repose entièrement sur des contributions bénévoles / volontaires. »
3. \textbf{voluntarily} (adverbe modifiant le verbe \eng{offered}). \eng{He \textbf{voluntarily} offered to drive the children to school.} $\rightarrow$ « Il s'est porté volontaire / a volontairement proposé de conduire les enfants à l'école. »
4. \textbf{volunteers} (nom pluriel, sujet de \eng{are needed} implicite ou COD de \eng{need}) / \textbf{volunteer} (adjectif attributif ou nom composé \eng{volunteer campaign} mais ici \eng{volunteer campaign} est moins courant que \eng{voluntary campaign} ou \eng{volunteering campaign}. \textbf{Piège :} \eng{volunteer} peut être adjectif attributif = \eng{a volunteer firefighter}. Mais pour « campagne de bénévoles », on dit \eng{a volunteering campaign} ou \eng{a volunteer recruitment campaign}. Ici, le contexte \eng{volunteer campaign} est acceptable comme \eng{volunteer drive}). \eng{We need more \textbf{volunteers} for the \textbf{volunteer} campaign next week.} $\rightarrow$ « Nous avons besoin de plus de bénévoles pour la campagne de recrutement de bénévoles la semaine prochaine. »
5. \textbf{Voluntary} (adjectif devant \eng{work}). \eng{\textbf{Voluntary} work looks great on a university application.} $\rightarrow$ « Le travail bénévole fait bonne impression sur une dossier universitaire. »

\textbf{Rappel règle :} \eng{Volunteer} = Nom/Verbe | \eng{Voluntary} = Adjectif | \eng{Voluntarily} = Adverbe.
\end{exoresolu}

\courssub{Exercices d'entraînement (À faire dans votre cahier)}

\begin{exercice}[Entraînement lexical — Section 1]
\textbf{A. Vocabulary Matching.} Associez chaque terme anglais à sa définition française.
1. Outreach programme \quad 2. Commitment \quad 3. To empower \quad 4. Community centre \quad 5. Reliability
a) Engagement, dévouement ferme \quad b) Fiabilité \quad c) Programme de proximité / sensibilisation \quad d) Centre communautaire \quad e) Autonomiser, donner du pouvoir

\textbf{B. Translation.} Traduisez en anglais en utilisant le vocabulaire de la leçon.
1. Elle fait du bénévolat dans un centre communautaire tous les samedis.
2. Les militants plaident pour l'égalité des sexes.
3. Ce travail de terrain demande beaucoup de dévouement et d'empathie.
4. Il a donné de son temps volontairement pour la campagne de sensibilisation.
5. L'ONG organise des ateliers de formation pour autonomiser les femmes rurales.

\textbf{C. Error Correction.} Corrigez les erreurs (grammaire, vocabulaire, faux ami) dans les phrases suivantes.
1. She is a voluntary in a big charity.
2. We must make voluntary work to help the community.
3. He worked voluntary for two years.
4. The NGO makes workshops for girls.
5. She has a strong motivation and she is very reliable person.

\textbf{D. Prononciation.} Soulignez la syllabe accentuée (en majuscules) : volunTEER / volunTARY / volunTARily / CHArity / comMITment / acTIVist.
\end{exercice}

\begin{corrige}[Exercice Entraînement lexical — Section 1]
\textbf{A.} 1-c, 2-a, 3-e, 4-d, 5-b.
\textbf{B.}
1. \eng{She volunteers at a community centre every Saturday.} (ou \eng{does voluntary work at...})
2. \eng{Activists advocate for gender equality.}
3. \eng{This field work requires a lot of dedication and empathy.}
4. \eng{He voluntarily gave his time for the outreach campaign.}
5. \eng{The NGO organises training workshops to empower rural women.}
\textbf{C.}
1. \eng{She is a \textbf{volunteer} in a big charity.} (\eng{voluntary} = adjectif).
2. \eng{We must \textbf{do voluntary work} / \textbf{volunteer} to help the community.} (\eng{make} $\rightarrow$ \eng{do}).
3. \eng{He worked \textbf{voluntarily} for two years.} (Adverbe).
4. \eng{The NGO \textbf{organises/runs/holds} workshops for girls.} (\eng{make workshops} $\rightarrow$ collocation incorrecte).
5. \eng{She has strong motivation and \textbf{is a} very reliable person.} (Ajout article \eng{a} devant nom singulier).
\textbf{D.} volun\textbf{TEER} / \textbf{VOL}untary (ou vol\textbf{UN}tary selon dictionnaires, mais standard UK \textbf{VOL}untary) / vol\textbf{UN}tarily / \textbf{CHA}rity / com\textbf{MIT}ment / \textbf{AC}tivist.
\end{corrige}

\coursec{Champ lexical 2 : Gender equality and women's rights}

Dans la section précédente, nous avons découvert le vocabulaire du \cle{volunteering} et du \cle{community service} : \eng{to volunteer}, \eng{a charity}, \eng{to raise awareness}, \eng{to give back to the community}. Mais pour quelles causes les volontaires s'engagent-ils le plus souvent ? Au Cameroun comme dans les pays anglophones, l'une des causes majeures est la promotion de l'égalité entre les hommes et les femmes. Cette deuxième famille de mots est absolument centrale dans le chapitre \emph{Citizenship and Human Rights} : c'est elle qui vous permettra de comprendre les textes du Bac et de rédiger vos propres compositions sur les droits humains.

\courssub{Observer d'abord : un texte pour découvrir le lexique}

Lisons d'abord un court article de presse original. Soulignez mentalement tous les mots qui se rapportent à l'égalité entre hommes et femmes.

\begin{textebox}[Article — “Volunteers for Equality”, The Bamenda Herald]
In the North-West Region of Cameroon, a group of young volunteers has launched a campaign to promote \textbf{gender equality} in rural communities. Their association, called “Equal Chances”, believes that \textbf{women's rights} are human rights and that no society can develop while half of its population faces \textbf{discrimination}.

Every Saturday, the volunteers organise workshops in villages around Bamenda. They talk to traditional rulers, parents and teenagers about the \textbf{gender gap} that still exists in education: in some families, boys are sent to secondary school while girls are kept at home or pushed into \textbf{early marriage}. “We want to \textbf{challenge stereotypes},” explains Mirabel, a twenty-two-year-old volunteer. “Many people still believe that a girl's place is only in the kitchen. This kind of \textbf{prejudice} destroys the future of thousands of girls.”

The association also works with survivors of \textbf{gender-based violence}. Trained volunteers listen to the victims, guide them towards legal services and help them recover their \textbf{dignity}. “\textbf{Empowering women} is our main goal,” says Mirabel. “When a woman is educated and economically independent, the whole community benefits.”

The volunteers insist that their fight is not against men but against \textbf{inequality} and \textbf{injustice}. They invite fathers and husbands to their meetings, because real change requires everyone's participation. Their dream is simple: a Cameroon where every child, boy or girl, enjoys \textbf{equal opportunities}, where \textbf{girls' education} is a priority, and where \textbf{fairness} and \textbf{justice} are not just words in official speeches but a daily reality for all citizens.
\end{textebox}

\begin{definition}[Glossaire du texte]
\begin{itemize}
\item \eng{gender gap} (n.) : écart, fossé entre les sexes
\item \eng{early marriage} (n.) : mariage précoce
\item \eng{to challenge stereotypes} (v.) : remettre en question les stéréotypes
\item \eng{prejudice} (n.) : préjugé
\item \eng{survivor} (n.) : survivant(e), personne ayant survécu à des violences
\item \eng{dignity} (n.) : dignité
\item \eng{fairness} (n.) : équité, justice dans le traitement
\end{itemize}
\end{definition}

\textbf{Questions de compréhension}\par
\begin{enumerate}
\item What is the name of the volunteers' association and where does it work?
\item According to Mirabel, what prejudice destroys the future of many girls?
\item How does the association help survivors of gender-based violence?
\item Why do the volunteers invite men to their meetings?
\end{enumerate}

\begin{corrige}[Questions de compréhension]
1. \eng{The association is called “Equal Chances” and it works in villages around Bamenda, in the North-West Region of Cameroon.} 2. \eng{The prejudice that a girl's place is only in the kitchen destroys the future of many girls.} 3. \eng{Trained volunteers listen to the victims, guide them towards legal services and help them recover their dignity.} 4. \eng{They invite men because real change requires everyone's participation.}
\end{corrige}

\courssub{Les concepts clés : gender, equality, equity, rights}

\begin{definition}[Gender equality]
\eng{Gender equality means that men and women have the same rights, opportunities and responsibilities.} (L'égalité des sexes signifie que les hommes et les femmes ont les mêmes droits, opportunités et responsabilités.)
\end{definition}

\begin{attention}[Sex ou gender ?]
\eng{Sex} désigne le \cle{sexe biologique} (male / female) ; \eng{gender} désigne le \cle{genre social et culturel} (les rôles construits par la société). Ne traduisez jamais « égalité des sexes » par \eng{sex equality} : la seule expression correcte est \eng{gender equality}. De même, on dit \eng{gender gap}, \eng{gender-based violence}, \eng{gender roles} — jamais avec \eng{sex}.
\end{attention}

\begin{propriete}[Equality ou equity ?]
\cle{Equality} (\eng{gender equality}) signifie l'\textbf{égalité de droits} : donner la même chose à tous. \cle{Equity} (\eng{gender equity}) signifie l'\textbf{équité} : traiter chacun \emph{justement selon ses besoins}, ce qui peut exiger des mesures différentes (par exemple des bourses réservées aux filles pour rattraper un retard scolaire). Image classique : \emph{equality} donne la même paire de chaussures à tout le monde ; \emph{equity} donne à chacun une paire à sa pointure.
\end{propriete}

\begin{exemplebox}[Les concepts en phrases]
\eng{The NGO fights gender-based violence in rural areas.} (L'ONG lutte contre les violences basées sur le genre en milieu rural.)\par
\eng{Empowering women benefits the whole community.} (Autonomiser les femmes profite à toute la communauté.)\par
\eng{The government wants to close the gender gap in education.} (Le gouvernement veut combler l'écart entre les sexes dans l'éducation.)\par
\eng{Equity means giving girls the specific support they need.} (L'équité signifie donner aux filles le soutien spécifique dont elles ont besoin.)\par
\eng{Women's rights are human rights.} (Les droits des femmes sont des droits humains.)
\end{exemplebox}

\begin{propriete}[Gender-based violence (GBV)]
\eng{Gender-based violence} (GBV) désigne les violences commises contre une personne \textbf{en raison de son sexe ou de son genre} : violences physiques, psychologiques, sexuelles ou économiques. C'est le terme officiel utilisé par l'ONU et les ONG — à connaître absolument pour le Bac. On dit aussi \eng{violence against women} lorsque les victimes sont spécifiquement des femmes. Attention : \eng{violence} est \textbf{incompatable avec le pluriel} en anglais (nom indénombrable) : on ne dit jamais « violences ».
\end{propriete}

\courssub{Problèmes dénoncés et objectifs poursuivis}

\begin{center}
\begin{tabularx}{\linewidth}{|X|X|X|}\hline
\rowcolor{popblueL} English & Français & Exemple en contexte \\ \hline
discrimination & discrimination & \eng{Women face discrimination at work.} \\ \hline
sexism & sexisme & \eng{Sexism limits girls' ambitions.} \\ \hline
inequality & inégalité & \eng{They fight inequality in pay.} \\ \hline
stereotype & stéréotype & \eng{“Boys don't cry” is a stereotype.} \\ \hline
prejudice & préjugé & \eng{Prejudice against women leaders is common.} \\ \hline
harassment & harcèlement & \eng{She reported sexual harassment.} \\ \hline
early marriage & mariage précoce & \eng{Early marriage ends many girls' studies.} \\ \hline
exclusion & exclusion & \eng{The exclusion of widows must stop.} \\ \hline
empowerment & autonomisation & \eng{Women's empowerment changes societies.} \\ \hline
inclusion & inclusion & \eng{The inclusion of women in politics is vital.} \\ \hline
equal opportunities & chances égales & \eng{Every child deserves equal opportunities.} \\ \hline
girls' education & éducation des filles & \eng{Girls' education reduces poverty.} \\ \hline
dignity / justice / fairness & dignité / justice / équité & \eng{Victims deserve dignity and justice.} \\ \hline
\end{tabularx}
\end{center}

\begin{exemplebox}[Les verbes d'action : collocations essentielles]
\eng{to promote gender equality} (promouvoir l'égalité des sexes)\par
\eng{to empower women} (autonomiser les femmes)\par
\eng{to defend women's rights} (défendre les droits des femmes)\par
\eng{to fight discrimination} (lutter contre la discrimination)\par
\eng{to challenge stereotypes} (combattre les stéréotypes)\par
\eng{to advocate for girls' education} (plaider pour l'éducation des filles)
\end{exemplebox}

\begin{popfigure}
\begin{tikzpicture}[
  box/.style={draw=popline, rounded corners, fill=popcream, align=center, text width=4.2cm, minimum height=2.6cm, font=\small},
  goal/.style={draw=popline, rounded corners, fill=popgreenL, align=center, text width=4.2cm, minimum height=2.6cm, font=\small},
  arr/.style={-{Stealth[length=3mm]}, very thick, poporange}]
\node[box, align=center] (pb) at (0,0){\textbf{Problems to fight}\\[2pt] discrimination\\ sexism\\ gender-based violence\\ stereotypes};
\node[goal, align=center] (gl) at (8,0){\textbf{Goals to reach}\\[2pt] equality\\ empowerment\\ inclusion\\ equal opportunities};
\draw[arr] (pb) -- node[above, font=\small\itshape, text=popdark]{volunteering} (gl);
\end{tikzpicture}
\legende{Le volontariat transforme les problèmes en objectifs atteints.}
\end{popfigure}

\courssub{Prononciation et orthographe}

\begin{aretenir}[Prononciation des mots-clés]
\begin{itemize}
\item \eng{equality} : « i-KOUO-li-ti » — accent sur la 2\textsuperscript{e} syllabe (attention : \eng{equal} se dit « I-kouol », accent sur la 1\textsuperscript{re} syllabe).
\item \eng{discrimination} : « dis-kri-mi-NEI-cheun » — accent sur « nei ».
\item \eng{harassment} : « HA-reuss-ment » (britannique) ou « heu-RASS-ment » (américain) — les deux sont acceptées.
\item \eng{empowerment} : « im-PAOU-eu-ment » ; \eng{stereotype} : « STÉ-rio-taïpe ».
\end{itemize}
\end{aretenir}

\begin{exoresolu}[Complétez avec le mot du champ lexical]
Énoncé — Complétez chaque phrase avec le mot correct : \emph{gender equality, harassment, stereotype, empowerment, gender gap}.\par
1. The association trains women in business skills: this is women's \ldots\par
2. The \ldots{} between boys' and girls' school results is getting smaller.\par
3. The campaign denounces sexual \ldots{} in public transport.\par
4. Our club works to promote \ldots{} in our school.\par
5. The idea that girls cannot study science is just a \ldots
\tcblower
Solution détaillée — 1. \eng{empowerment} (la formation professionnelle autonomise les femmes). 2. \eng{gender gap} (l'écart entre les résultats des garçons et des filles). 3. \eng{harassment} (le harcèlement sexuel dans les transports ; nom indénombrable, sans article). 4. \eng{gender equality} (promouvoir l'égalité des sexes). 5. \eng{stereotype} (au singulier, car l'article indéfini \eng{a} précède : \eng{a stereotype}).
\end{exoresolu}

\begin{attention}[Pièges des francophones]
\begin{itemize}
\item Ne dites pas \eng{the equality of sexes} : dites \eng{gender equality}, sans article et sans préposition.
\item \eng{Discrimination} se construit avec \eng{against} : \eng{discrimination against women} (et non « discrimination of women »).
\item Orthographe : \eng{harassment} s'écrit avec \textbf{un seul r} et \textbf{deux s} (ni « harrassment » ni « harrasment ») ; \eng{prejudice} s'écrit sans \emph{d} après \emph{preju-} (ni « predjudice »).
\item \eng{To fight} peut se construire directement : \eng{to fight discrimination} ; la forme \eng{to fight against discrimination} est également correcte, mais ne dites jamais « to fight for discrimination » — \eng{to fight for} signifie « se battre en faveur de » : \eng{to fight for women's rights} (se battre pour les droits des femmes).
\item Apostrophes : \eng{women's rights} (les droits des femmes — \eng{women} est déjà pluriel, l'apostrophe précède le \emph{s}) ; \eng{girls' education} (l'éducation des filles — apostrophe après le \emph{s} du pluriel régulier).
\end{itemize}
\end{attention}

\begin{savaistu}[SDG 5 : Gender Equality]
L'objectif 5 des Objectifs de Développement Durable (SDG 5) de l'ONU est précisément \eng{Gender Equality} : « Achieve gender equality and empower all women and girls. » Ce repère est très souvent cité dans les sujets de Bac et dans les discours officiels — au Cameroun, de nombreuses ONG de volontaires alignent leurs actions sur les SDGs. Le connaître enrichira vos compositions.
\end{savaistu}

\begin{exercice}[À vous de jouer]
Traduisez en anglais : 1. Les volontaires défendent les droits des femmes dans les villages. 2. L'éducation des filles réduit les mariages précoces. 3. Nous devons combattre les stéréotypes et les préjugés. 4. L'autonomisation des femmes profite à toute la communauté.
\end{exercice}

\begin{corrige}[Exercice — traduction]
1. \eng{Volunteers defend women's rights in the villages.} 2. \eng{Girls' education reduces early marriages.} 3. \eng{We must challenge stereotypes and prejudices.} 4. \eng{Empowering women benefits the whole community.}
\end{corrige}

\coursec{Collocations et expressions idiomatiques}

Dans la section précédente, nous avons découvert le vocabulaire de l'égalité des sexes et des droits des femmes : \eng{gender equality}, \eng{women's rights}, \eng{empowerment}, \eng{discrimination}. Mais connaître des mots isolés ne suffit pas pour parler et écrire un anglais naturel. Un mot ne voyage jamais seul : il s'associe à des partenaires privilégiés. Ces associations habituelles s'appellent des \cle{collocations}. C'est elles qui font la différence entre un élève qui « traduit du français » et un élève qui « pense en anglais ».

\begin{definition}[Collocation]
Une \cle{collocation} est une combinaison habituelle et stable de deux mots ou plus (verbe + nom, adjectif + nom, verbe + préposition) que les anglophones emploient spontanément ensemble. On dit \eng{to raise awareness} et non \eng{to lift awareness} ; \eng{to fight against discrimination} et non \eng{to war discrimination}. La collocation ne se devine pas : elle s'apprend comme un bloc.
\end{definition}

\begin{textebox}[Reading text: “A Volunteer in Bamenda”]
Every Saturday morning, Neh, a nineteen-year-old student, puts on her blue T-shirt and joins a local association in Bamenda. She signed up as a volunteer six months ago, after attending a workshop on gender equality. Her team works hard to raise awareness about early marriage and to promote equality between boys and girls at school. Last term, they helped to raise funds to buy books for thirty girls who had dropped out of school.

Neh is not afraid to raise her voice. During community meetings, she speaks out against harassment and challenges stereotypes about what women can or cannot do. “Many people think a girl's place is in the kitchen,” she says. “We want to break the silence and give women a voice.”

The association also defends the rights of widows and supports victims of domestic violence. Its members advocate for change in front of local authorities and campaign for girls' education in remote villages. Thanks to their door-to-door visits, they reach out to rural women who never listen to the radio. “When we lend a helping hand to one woman, the whole family benefits,” Neh explains. “Little by little, we are bridging the gender gap. We know we can make a difference.”

Last month, Neh was invited to speak on a community radio programme. She encouraged young people to help out in their neighbourhoods. “You do not need money to volunteer,” she concluded. “You only need time, courage and love for your community.”
\end{textebox}

\textbf{Glossaire}\par
\begin{itemize}
\item \eng{to sign up} (phrasal verb) : s'inscrire, s'engager
\item \eng{a workshop} (n.) : un atelier (de formation)
\item \eng{to drop out of school} : abandonner l'école, décrocher
\item \eng{a widow} (n.) : une veuve
\item \eng{domestic violence} (n.) : la violence domestique
\item \eng{remote} (adj.) : reculé, isolé
\item \eng{door-to-door} (adj.) : du porte-à-porte
\item \eng{a neighbourhood} (n.) : un quartier
\end{itemize}

\textbf{Comprehension questions}\par
\begin{enumerate}
\item Why did Neh sign up as a volunteer?
\item What did the team buy with the funds they raised?
\item Who does the association reach out to, and how?
\item According to Neh, what do you need to become a volunteer?
\end{enumerate}

Observez le texte : chaque verbe important est accompagné d'un nom « partenaire » : \eng{raise awareness}, \eng{raise funds}, \eng{raise her voice}, \eng{promote equality}, \eng{defend the rights}, \eng{support victims}, \eng{challenge stereotypes}, \eng{break the silence}, \eng{bridge the gender gap}. Étudions maintenant ces familles de collocations une par une.

\courssub{Les collocations avec \eng{raise}}

\begin{propriete}[Trois collocations essentielles avec \eng{raise}]
\eng{To raise} (raised, raised) signifie « élever, soulever, augmenter ». Dans le langage du volontariat et des ONG, il forme trois collocations incontournables :
\begin{itemize}
\item \eng{to raise awareness (about/of something)} : sensibiliser (à quelque chose)
\item \eng{to raise funds / to raise money (for something)} : collecter des fonds (pour)
\item \eng{to raise one's voice (against something)} : élever la voix, protester (contre)
\end{itemize}
Notez la préposition : \eng{awareness \textbf{about/of}}, \eng{funds \textbf{for}}, \eng{voice \textbf{against}}.
\end{propriete}

\begin{exemplebox}[Exemples traduits]
\eng{Volunteers raise awareness about early marriage.} « Les volontaires sensibilisent au mariage précoce. »\par
\eng{The students raised funds for the girls' hostel.} « Les élèves ont collecté des fonds pour le foyer des filles. »\par
\eng{She raised her voice against injustice.} « Elle a élevé la voix contre l'injustice. »
\end{exemplebox}

\begin{popfigure}
\begin{center}
\begin{tikzpicture}[node distance=1.1cm]
\node[draw=popblue, fill=popblueL, rounded corners, thick, font=\bfseries\large, minimum width=2.6cm, minimum height=1cm] (raise) {RAISE};
\node[draw=poporange, fill=poporangeL, rounded corners, thick, right=3.2cm of raise, yshift=1.6cm, align=center, minimum width=4.6cm] (aw) {awareness\\ \emph{\small sensibiliser}};
\node[draw=popgreen, fill=popgreenL, rounded corners, thick, right=3.2cm of raise, align=center, minimum width=4.6cm] (fu) {funds / money\\ \emph{\small collecter des fonds}};
\node[draw=poppurple, fill=poppurpleL, rounded corners, thick, right=3.2cm of raise, yshift=-1.6cm, align=center, minimum width=4.6cm] (vo) {one's voice\\ \emph{\small élever la voix}};
\draw[-{Stealth}, very thick, poporange] (raise) -- (aw);
\draw[-{Stealth}, very thick, popgreen] (raise) -- (fu);
\draw[-{Stealth}, very thick, poppurple] (raise) -- (vo);
\end{tikzpicture}
\end{center}
\legende{Organigramme : le verbe \eng{raise} et ses trois collocations majeures.}
\end{popfigure}

\begin{methode}[Apprendre un verbe avec son nom habituel]
\begin{enumerate}
\item N'apprenez jamais un verbe seul : mémorisez le couple \cle{verbe + nom} comme un seul bloc, par exemple \eng{raise + awareness}.
\item Dans le langage des ONG, on dit \eng{to raise awareness} et non \eng{to increase awareness} ; \eng{increase} convient aux chiffres (\eng{increase the number of girls at school}) mais pas à la sensibilisation.
\item Faites des fiches : verbe à gauche, nom partenaire à droite, un exemple personnel en dessous.
\item Réemployez chaque collocation dans une phrase sur votre propre contexte (votre école, votre quartier) : c'est la répétition en contexte qui fixe la mémoire.
\end{enumerate}
\end{methode}

\begin{attention}[Le piège de « sensibiliser »]
\eng{To raise awareness} = sensibiliser. Ne dites \textbf{jamais} \eng{to sensibilize} : ce mot n'existe pas en anglais. Le verbe \eng{to sensitise} (britannique) / \eng{to sensitize} (américain) existe, mais il est rare et technique (médecine, photographie). Le francophone prudent utilisera \eng{to raise awareness about...} ou, à la rigueur, \eng{to sensitise people about...}. De même, « une campagne de sensibilisation » se dit \eng{an awareness campaign} ou \eng{an awareness-raising campaign}, jamais \eng{a sensibilization campaign}.
\end{attention}

\courssub{Les collocations avec \eng{promote}, \eng{defend}, \eng{fight} et \eng{combat}}

\begin{propriete}[Promouvoir, défendre, combattre]
\begin{itemize}
\item \eng{to promote equality / girls' education / peace} : promouvoir l'égalité, l'éducation des filles, la paix. Attention : \eng{promote} ne signifie pas « promouvoir un élève » au sens scolaire français (on dirait \eng{to move a pupil up to the next class}).
\item \eng{to defend rights / a cause} : défendre des droits, une cause. On défend quelqu'un \eng{against} une attaque : \eng{to defend women against violence}.
\item \eng{to fight \textbf{against} discrimination / to fight \textbf{for} equality} : lutter contre la discrimination, lutter pour l'égalité. La préposition change le sens !
\item \eng{to combat sexism / violence} : combattre le sexisme, la violence (plus formel que \eng{fight} ; très apprécié à l'écrit). \eng{Combat} se construit \textbf{sans préposition} : \eng{to combat sexism}, jamais \eng{to combat against sexism}.
\end{itemize}
\end{propriete}

\begin{exemplebox}[Exemples traduits]
\eng{Our club promotes equality between boys and girls.} « Notre club promeut l'égalité entre les garçons et les filles. »\par
\eng{Lawyers defend the rights of widows.} « Les avocats défendent les droits des veuves. »\par
\eng{They fight against discrimination in the workplace.} « Ils luttent contre la discrimination au travail. »\par
\eng{The government must combat sexism in textbooks.} « Le gouvernement doit combattre le sexisme dans les manuels scolaires. »
\end{exemplebox}

\courssub{Les collocations avec \eng{empower}, \eng{support}, \eng{advocate} et \eng{campaign}}

\begin{propriete}[Autonomiser, soutenir, plaider]
\begin{itemize}
\item \eng{to empower women / girls / young people} : autonomiser les femmes, leur donner les moyens d'agir. Nom correspondant : \eng{empowerment} (l'autonomisation).
\item \eng{to support victims / survivors / a project} : soutenir les victimes, les survivantes, un projet. Attention : \eng{to support} ne veut pas dire « supporter » au sens de « tolérer » (voir la section sur les faux amis).
\item \eng{to advocate \textbf{for} change / for women's rights} : plaider pour le changement. Le nom correspondant est \eng{an advocate} (un défenseur, une défenseure) : \eng{she is an advocate of women's rights} « elle est une défenseure des droits des femmes ».
\item \eng{to campaign \textbf{for} girls' education / \textbf{against} early marriage} : mener campagne pour / contre. Noms : \eng{a campaign} (une campagne), \eng{a campaigner} (un militant).
\end{itemize}
\end{propriete}

\begin{exemplebox}[Exemples traduits]
\eng{Education empowers women.} « L'éducation autonomise les femmes. »\par
\eng{The centre supports victims of domestic violence.} « Le centre soutient les victimes de violence domestique. »\par
\eng{They advocate for a change in the law.} « Ils plaident pour un changement de la loi. »\par
\eng{We campaign for girls' education in rural areas.} « Nous menons campagne pour l'éducation des filles en milieu rural. »
\end{exemplebox}

\courssub{Les collocations avec \eng{break}, \eng{challenge} et \eng{bridge}}

\begin{propriete}[Rompre, contester, combler]
\begin{itemize}
\item \eng{to break the silence (about/on something)} : rompre le silence (sur un sujet tabou). Verbe irrégulier : \eng{break — broke — broken}.
\item \eng{to challenge stereotypes / prejudices / authority} : contester, remettre en question les stéréotypes, les préjugés. Ici, \eng{challenge} ne signifie pas « défier quelqu'un en duel » mais « questionner une idée reçue ».
\item \eng{to bridge the gender gap} : combler, réduire l'écart entre les sexes. Nom : \eng{the gender gap} (l'écart hommes-femmes). On dit aussi \eng{to close the gap} ou \eng{to narrow the gap} (réduire l'écart).
\end{itemize}
\end{propriete}

\begin{exemplebox}[Exemples traduits]
\eng{Survivors are breaking the silence about abuse.} « Les survivantes rompent le silence sur les abus. »\par
\eng{The play challenges stereotypes about housewives.} « La pièce remet en question les stéréotypes sur les femmes au foyer. »\par
\eng{Scholarships help to bridge the gender gap in science.} « Les bourses aident à combler l'écart entre les sexes dans les sciences. »
\end{exemplebox}

\begin{savaistu}[To break the glass ceiling]
\eng{To break the glass ceiling} — « briser le plafond de verre » — désigne le fait, pour une femme, d'atteindre des postes de responsabilité jusque-là réservés aux hommes (directrice générale, ministre, rectrice). Le « plafond de verre » est invisible : personne ne dit ouvertement qu'une femme ne peut pas diriger, mais l'obstacle existe. Cette expression, très imagée, est appréciée dans une copie de Bac : \eng{More women are breaking the glass ceiling in Cameroonian politics.} « De plus en plus de femmes brisent le plafond de verre dans la politique camerounaise. »
\end{savaistu}

\courssub{Expressions idiomatiques et verbes à particule}

\begin{definition}[Expression idiomatique]
Une \cle{expression idiomatique} est un groupe de mots dont le sens global ne se déduit pas du sens de chaque mot. \eng{To lend a helping hand} ne parle ni de prêt ni d'une main : cela signifie « donner un coup de main ». Les idioms ne se traduisent pas mot à mot ; ils s'apprennent entiers.
\end{definition}

\begin{center}
\begin{tabularx}{\linewidth}{|X|X|X|}
\hline
\rowcolor{popblueL} \textbf{English idiom} & \textbf{Français} & \textbf{Exemple en contexte} \\
\hline
\eng{to give someone a voice} & donner la parole à quelqu'un & \eng{Radio gives rural women a voice.} \\
\hline
\eng{to stand up for one's rights} & défendre ses droits & \eng{She stood up for her rights at work.} \\
\hline
\eng{to make a difference} & faire une différence, changer les choses & \eng{One volunteer can make a difference.} \\
\hline
\eng{to lend a helping hand} & donner un coup de main & \eng{Neighbours lent a helping hand after the flood.} \\
\hline
\eng{to break the glass ceiling} & briser le plafond de verre & \eng{She broke the glass ceiling in banking.} \\
\hline
\end{tabularx}
\end{center}

\begin{exemplebox}[Exemple traduit]
\eng{We want to make a difference.} « Nous voulons faire une différence. »
\end{exemplebox}

\begin{propriete}[Quatre verbes à particule du volontariat]
Un \cle{verbe à particule} (\eng{phrasal verb}) est un verbe suivi d'une petite particule qui en change le sens. La particule fait partie intégrante du sens :
\begin{itemize}
\item \eng{to help out (with something)} : aider, donner un coup de main. \eng{Can you help out with the cooking?} « Peux-tu donner un coup de main pour la cuisine ? »
\item \eng{to speak out (against something)} : s'exprimer publiquement, dénoncer. \eng{She spoke out against harassment.} « Elle a pris la parole publiquement contre le harcèlement. »
\item \eng{to reach out \textbf{to} somebody} : aller vers quelqu'un, prendre contact avec quelqu'un. \eng{The campaign reached out to rural women.} « La campagne est allée vers les femmes rurales. »
\item \eng{to sign up (as a volunteer / for an activity)} : s'inscrire, s'engager. \eng{They signed up as volunteers last month.} « Ils se sont inscrits comme volontaires le mois dernier. »
\end{itemize}
Attention à l'ordre des mots : avec un pronom, on sépare le verbe et la particule quand le verbe est séparable : \eng{she helped them out} « elle les a aidés » ; mais \eng{reach out to} est inséparable : \eng{she reached out to them}, jamais \eng{she reached them out to}.
\end{propriete}

\begin{center}
\begin{tabularx}{\linewidth}{|X|X|X|}
\hline
\rowcolor{popgreenL} \textbf{Français (calque à éviter)} & \textbf{Anglais correct} & \textbf{Remarque} \\
\hline
sensibiliser & \eng{to raise awareness} & jamais \eng{to sensibilize} \\
\hline
collecter des fonds & \eng{to raise funds} & \eng{to collect money} est possible, mais \eng{raise funds} est la collocation des ONG \\
\hline
lutter contre & \eng{to fight against / to combat} & \eng{against} avec \eng{fight} ; aucune préposition avec \eng{combat} \\
\hline
mener campagne pour & \eng{to campaign for} & pas \eng{to make a campaign} \\
\hline
s'inscrire comme volontaire & \eng{to sign up as a volunteer} & pas \eng{to inscribe oneself} \\
\hline
donner un coup de main & \eng{to lend a helping hand / to help out} & pas de traduction mot à mot \\
\hline
\end{tabularx}
\end{center}

\begin{exoresolu}[Collocations : complétez avec le bon verbe]
Énoncé — Complete each sentence with the correct verb in the correct form: \eng{raise, promote, defend, break, challenge, reach out to, sign up, make}.
\begin{enumerate}
\item Last year, our association \ldots\ funds to build a library for girls.
\item The minister encouraged women to \ldots\ the silence about violence.
\item Fifty students \ldots\ as volunteers for the clean-up campaign.
\item Posters \ldots\ equality between boys and girls in our school.
\item Small actions can \ldots\ a big difference.
\end{enumerate}
\tcblower
Solution détaillée :
\begin{enumerate}
\item \eng{raised} — collocation \eng{to raise funds} ; « last year » impose le prétérit : \eng{Last year, our association raised funds to build a library for girls.} « L'an dernier, notre association a collecté des fonds pour construire une bibliothèque pour les filles. »
\item \eng{break} — collocation \eng{to break the silence} ; après \eng{encouraged women to}, on met la base verbale : \eng{...to break the silence about violence.} « ...à rompre le silence sur la violence. »
\item \eng{signed up} — verbe à particule \eng{to sign up as volunteers} ; prétérit : \eng{signed up}. « Cinquante élèves se sont inscrits comme volontaires pour la campagne de nettoyage. »
\item \eng{promote} — collocation \eng{to promote equality} ; présent simple, sujet pluriel : \eng{Posters promote equality between boys and girls in our school.} « Les affiches promeuvent l'égalité entre les garçons et les filles dans notre école. »
\item \eng{make} — idiom \eng{to make a difference} ; après \eng{can}, base verbale : \eng{Small actions can make a big difference.} « De petites actions peuvent faire une grande différence. »
\end{enumerate}
\end{exoresolu}

\begin{attention}[Erreurs fréquentes des francophones]
\begin{itemize}
\item Ne dites pas \eng{to do a difference} : l'idiom exige \eng{make} : \eng{to \textbf{make} a difference}.
\item Pour « lutter contre », la forme la plus sûre à l'examen est \eng{to fight \textbf{against} discrimination}. (\eng{To fight discrimination} sans préposition se rencontre dans la presse américaine, mais la forme avec \eng{against} est recommandée.)
\item Ne dites pas \eng{to elevate awareness} ni \eng{to lift funds} : seul \eng{raise} convient.
\item \eng{To advocate} ne prend jamais \eng{about} : on dit \eng{to advocate \textbf{for} change}.
\item \eng{To combat} est transitif direct : \eng{to combat violence}, jamais \eng{to combat against violence}.
\item Prononciation : \eng{volunteer} se prononce « vo-leun-tièr » (accent sur la dernière syllabe) ; \eng{awareness} se prononce « a-wéèr-ness » ; \eng{challenge} se prononce « tcha-lindj ».
\end{itemize}
\end{attention}

\begin{experience}[Production guidée : la réunion de l'association]
Par groupes de quatre, préparez en cinq minutes le discours d'ouverture de la réunion d'une association de volontaires de votre lycée. Consigne : utilisez au moins \textbf{six} collocations ou idioms de cette section, par exemple : \eng{raise awareness}, \eng{raise funds}, \eng{sign up}, \eng{speak out}, \eng{make a difference}, \eng{lend a helping hand}, \eng{bridge the gender gap}. Un porte-parole présente le discours à la classe ; les autres groupes cochent sur une grille les collocations entendues. Modèle de début : \eng{Dear friends, thank you for signing up as volunteers. This year, we want to raise awareness about...}
\end{experience}

\begin{exercice}[À vous de jouer]
\begin{enumerate}
\item Traduisez en anglais : « Les volontaires ont rompu le silence et ont sensibilisé le village au mariage précoce. »
\item Trouvez la collocation correcte : \eng{to (do / make / give) a difference} ; \eng{to (raise / rise / lift) funds} ; \eng{to (speak / say / tell) out against violence}.
\item Rédigez trois phrases sur une association réelle ou imaginaire de votre ville, avec \eng{campaign for}, \eng{reach out to} et \eng{stand up for}.
\end{enumerate}
\end{exercice}

\begin{corrige}[Exercice : Collocations et expressions idiomatiques]
\begin{enumerate}
\item \eng{The volunteers broke the silence and raised awareness about early marriage in the village.} (Prétérit irrégulier : \eng{break — broke} ; collocation : \eng{raise awareness about}.)
\item \eng{to \textbf{make} a difference} ; \eng{to \textbf{raise} funds} ; \eng{to \textbf{speak} out against violence}.
\item Exemple de production attendue : \eng{The association “Girls First” campaigns for girls' education in Douala. Its members reach out to poor families in the markets. They stand up for the rights of every child.} « L'association “Girls First” mène campagne pour l'éducation des filles à Douala. Ses membres vont vers les familles pauvres dans les marchés. Ils défendent les droits de chaque enfant. »
\end{enumerate}
\end{corrige}

\begin{aretenir}[L'essentiel de la section]
\begin{itemize}
\item Une collocation s'apprend en bloc : \eng{raise awareness}, \eng{raise funds}, \eng{raise one's voice}.
\item \eng{promote equality}, \eng{defend rights}, \eng{fight against discrimination}, \eng{combat sexism} (sans préposition).
\item \eng{empower women}, \eng{support victims}, \eng{advocate for change}, \eng{campaign for girls' education}.
\item \eng{break the silence} (\eng{break — broke — broken}), \eng{challenge stereotypes}, \eng{bridge the gender gap}, \eng{break the glass ceiling}.
\item Idioms : \eng{give someone a voice}, \eng{stand up for one's rights}, \eng{make a difference}, \eng{lend a helping hand}.
\item Phrasal verbs : \eng{help out}, \eng{speak out}, \eng{reach out to}, \eng{sign up}.
\item Jamais \eng{to sensibilize} : dites \eng{to raise awareness}.
\end{itemize}
\end{aretenir}

\coursec{Faux amis et pièges des francophones}

Dans la section précédente, nous avons appris à combiner correctement les mots du bénévolat et de l'égalité des genres grâce aux \cle{collocations} (\eng{to raise awareness}, \eng{to empower women}, \eng{to promote equality}). Mais connaître les bons mots ne suffit pas : il faut aussi savoir lesquels \emph{éviter}. Le français et l'anglais partagent des milliers de mots d'origine latine qui se ressemblent… sans toujours signifier la même chose. Ces mots trompeurs s'appellent les \cle{faux amis} (\eng{false friends} ou \eng{false cognates}). Dans le domaine du volontariat et de la promotion de l'égalité, ils sont particulièrement nombreux : on « assiste » à des ateliers, on « anime » des campagnes, on « sensibilise » les communautés, on suit des « formations »… Autant de pièges que cette section va désamorcer un à un.

\courssub{Qu'est-ce qu'un faux ami ? Observation d'abord}

\begin{definition}[Faux ami]
Un \cle{faux ami} est un mot anglais qui ressemble par la forme à un mot français, mais dont le sens est différent — totalement ou partiellement. Il existe aussi des \cle{calques} : des expressions françaises traduites mot à mot en anglais, qui produisent des phrases incorrectes ou incompréhensibles pour un anglophone.
\end{definition}

Observez ces deux phrases écrites par des élèves de Terminale. L'une est correcte, l'autre contient un faux ami. Laquelle ?

\eng{1. I assisted to a workshop on women's rights last Saturday.}

\eng{2. I attended a workshop on women's rights last Saturday.}

La phrase 1 est \textbf{fausse} : le verbe \eng{to assist} ne signifie pas « assister à » mais « aider ». La phrase 2 est correcte : « assister à » (être présent) se dit \eng{to attend}. Voilà le mécanisme du faux ami : la ressemblance des formes (\eng{assist} / « assister ») nous pousse à l'erreur.

\begin{attention}[Le piège n° 1 : to assist]
\eng{I assisted to a workshop} est \textbf{FAUX} — doublement : mauvais verbe \emph{et} mauvaise préposition. Dites :
\begin{itemize}
\item \eng{I attended a workshop.} « J'ai assisté à un atelier. » (\eng{to attend} + nom, \textbf{sans} préposition)
\item \eng{I assisted the trainer.} « J'ai aidé le formateur. » (\eng{to assist} + personne = aider)
\end{itemize}
\end{attention}

\courssub{Les grands faux amis du thème : étude détaillée}

\textbf{1. To assist (aider) $\neq$ assister à (to attend)}\par
\eng{To assist} est un verbe formel qui signifie \textbf{aider} quelqu'un. « Assister à » un événement (être présent) se traduit par \eng{to attend}, qui se construit \textbf{sans préposition}.

\eng{Our volunteers assist women in rural areas.} « Nos bénévoles aident les femmes en milieu rural. »

\eng{She attended a training session on women's rights.} « Elle a assisté à une formation sur les droits des femmes. »

\eng{Two hundred people attended the conference in Yaoundé.} « Deux cents personnes ont assisté à la conférence à Yaoundé. »

\textbf{2. Sensible (raisonnable) $\neq$ sensible (fragile) = sensitive}\par
En anglais, \eng{sensible} signifie « raisonnable, plein de bon sens ». Le français « sensible » (émotif, fragile, réceptif) se traduit par \eng{sensitive}. Prononciation : \eng{sensible} « sé-ni-beul » ; \eng{sensitive} « sé-ni-ti-ve ».

\eng{It would be sensible to consult the community first.} « Il serait raisonnable de consulter d'abord la communauté. »

\eng{He is very sensitive to gender issues.} « Il est très sensible aux questions de genre. »

\eng{Domestic violence is a sensitive topic.} « La violence domestique est un sujet délicat. »

\textbf{3. To achieve (accomplir, atteindre) $\neq$ achever (to complete, to finish)}\par
\eng{To achieve} signifie « accomplir, atteindre » un objectif, un résultat — un mot central du vocabulaire des projets. Le français « achever » (terminer) se dit \eng{to complete} ou \eng{to finish}.

\eng{The project achieved its goals.} « Le projet a atteint ses objectifs. »

\eng{She completed her training in June.} « Elle a achevé sa formation en juin. »

\textbf{4. Un engagement (promesse, dévouement) = a commitment ; engagement (fiançailles) = an engagement}\par
Le mot anglais \eng{engagement} existe, mais il désigne surtout les \textbf{fiançailles} ou un rendez-vous professionnel. « Un engagement » au sens de promesse, de dévouement à une cause, se dit \eng{a commitment}.

\eng{Her commitment to the association is total.} « Son engagement envers l'association est total. »

\eng{Their engagement was announced last month.} « Leurs fiançailles ont été annoncées le mois dernier. »

\textbf{5. Une formation $\neq$ a formation}\par
« Une formation » (stage, apprentissage, session de formation) se traduit par \eng{a training} ou \eng{a training session}. En anglais, \eng{formation} désigne une formation \emph{géologique}, la constitution d'une équipe ou d'un gouvernement — jamais un stage. Attention : \eng{training} est un nom indénombrable dans ce sens ; on dit \eng{two training sessions}, pas \eng{two trainings}.

\eng{The NGO organised a training session for young leaders.} « L'ONG a organisé une formation pour de jeunes leaders. »

\eng{The formation of a new committee took place in Bamenda.} « La formation d'un nouveau comité a eu lieu à Bamenda. »

\textbf{6. Actuel (current) $\neq$ actual (réel, véritable)}\par
\eng{Actual} signifie « réel, véritable, effectif ». « Actuel » se traduit par \eng{current} ou \eng{present}. Même piège pour \eng{actually} (« en fait, en réalité ») $\neq$ « actuellement » (\eng{currently, at present}).

\eng{The current president of the association is a woman.} « La présidente actuelle de l'association est une femme. »

\eng{The actual number of volunteers is higher than expected.} « Le nombre réel de bénévoles est plus élevé que prévu. »

\eng{She is currently working in Douala.} « Elle travaille actuellement à Douala. »

\textbf{7. Autres faux amis fréquents dans ce thème}\par
Quelques pièges supplémentaires à connaître pour l'expression écrite :

\begin{center}
\begin{tabularx}{\linewidth}{|X|X|X|}
\hline
\rowcolor{popblueL} \textbf{Mot français} & \textbf{Faux ami anglais} & \textbf{Traduction correcte} \\
\hline
une opportunité, une occasion & an opportunity & \textbf{an opportunity} (vrai ami) — mais \eng{opportunity} $\neq$ \eng{opportunism} \\
\hline
un rapport (compte rendu) & a report & \textbf{a report} (vrai ami) — attention : \eng{rapport} (maths) = \eng{ratio} \\
\hline
un agenda (ordre du jour) & an agenda & \textbf{an agenda} (vrai ami) — mais \eng{diary} = agenda personnel (carnet de rendez-vous) \\
\hline
une campagne & a campaign & \textbf{a campaign} (vrai ami) — ne pas confondre avec \eng{campaign} (militaire) \\
\hline
un stage & an internship / a placement & \textbf{an internship} (stage pro), \textbf{a work placement} — \eng{a course} = un cours \\
\hline
\end{tabularx}
\end{center}

\courssub{Tableau récapitulatif des faux amis principaux}

\begin{center}
\begin{tabularx}{\linewidth}{|X|X|X|}
\hline
\rowcolor{popblueL} \textbf{Mot français} & \textbf{Faux ami anglais (à éviter)} & \textbf{Traduction correcte} \\
\hline
assister à (un atelier) & to assist (= aider) & to attend (a workshop) \\
\hline
sensible (fragile, réceptif) & sensible (= raisonnable) & sensitive \\
\hline
une formation (stage) & a formation (= géologique, d'une équipe) & a training / a training session \\
\hline
actuel & actual (= réel, véritable) & current / present \\
\hline
un engagement (à une cause) & an engagement (= fiançailles) & a commitment \\
\hline
achever (terminer) & to achieve (= atteindre, accomplir) & to complete / to finish \\
\hline
actuellement & actually (= en fait) & currently / at present \\
\hline
\end{tabularx}
\end{center}

\begin{popfigure}
\begin{tikzpicture}[node distance=1.2cm and 1.5cm]
\node[draw=popdark, fill=popblueL, rounded corners, font=\bfseries, align=center] (titre) {Faux amis : attention !};
\node[draw=popblue, fill=popblueL, rounded corners, align=center] (a) at (-5,-2) {assister à\\ \eng{to attend}};
\node[draw=poporange, fill=poporangeL, rounded corners, align=center] (b) at (-1.5,-2) {sensible (fragile)\\ \eng{sensitive}};
\node[draw=popgreen, fill=popgreenL, rounded corners, align=center] (c) at (2,-2) {une formation\\ \eng{a training}};
\node[draw=poppurple, fill=poppurpleL, rounded corners, align=center] (d) at (5.5,-2) {actuel / actuellement\\ \eng{current} / \eng{currently}};
\node[draw=popgold, fill=popgoldL, rounded corners, align=center] (e) at (-5,-4) {un engagement\\ \eng{a commitment}};
\node[draw=poppink, fill=poppinkL, rounded corners, align=center] (f) at (-1.5,-4) {achever\\ \eng{to complete}};
\node[draw=popdark, fill=popcream, rounded corners, align=center] (g) at (2,-4) {animer un atelier\\ \eng{to run/lead}};
\node[draw=popmuted, fill=popcream, rounded corners, align=center] (h) at (5.5,-4) {sensibiliser\\ \eng{to raise awareness}};
\draw[-{Stealth}, thick, popblue] (titre) -- (a);
\draw[-{Stealth}, thick, poporange] (titre) -- (b);
\draw[-{Stealth}, thick, popgreen] (titre) -- (c);
\draw[-{Stealth}, thick, poppurple] (titre) -- (d);
\draw[-{Stealth}, thick, popgold] (titre) -- (e);
\draw[-{Stealth}, thick, poppink] (titre) -- (f);
\draw[-{Stealth}, thick, popdark] (titre) -- (g);
\draw[-{Stealth}, thick, popmuted] (titre) -- (h);
\end{tikzpicture}
\legende{Carte mentale : le mot français et sa traduction anglaise correcte. En beige : les calques à remplacer.}
\end{popfigure}

\courssub{Les calques du français : des expressions qui n'existent pas en anglais}

Au-delà des mots isolés, le danger vient aussi des \cle{expressions traduites mot à mot}. Deux calques reviennent constamment dans les copies camerounaises sur ce thème :

\begin{propriete}[Calques à bannir]
\begin{itemize}
\item « Animer un atelier / une campagne » ne se dit \textbf{pas} \eng{to animate a workshop} (\eng{to animate} = donner vie à un dessin animé, ou encourager une foule). Dites : \eng{to lead a workshop}, \eng{to run a workshop}, \eng{to facilitate a workshop}, \eng{to conduct a campaign}.
\item « Sensibiliser (quelqu'un à quelque chose) » ne se dit \textbf{pas} \eng{to sensitize} (ce verbe existe en anglais technique/médical = « rendre sensible, allergiser », mais pas dans le sens « faire prendre conscience »). Dites : \eng{to raise awareness (of/about)}, \eng{to make people aware of}, \eng{to educate people about}.
\end{itemize}
\end{propriete}

\begin{exemplebox}[Les calques corrigés]
\eng{They run workshops in schools.} « Ils animent des ateliers dans les écoles. »

\eng{Our association raises awareness about girls' education.} « Notre association sensibilise à l'éducation des filles. »

\eng{Volunteers led a campaign to raise awareness of women's rights in Buea.} « Des bénévoles ont animé une campagne de sensibilisation aux droits des femmes à Buea. »

\eng{The facilitator conducted a training session on gender-based violence.} « L'animateur a mené une session de formation sur les violences basées sur le genre. »
\end{exemplebox}

\begin{methode}[La stratégie anti-faux amis]
\begin{enumerate}
\item Face à un mot anglais qui ressemble au français, \textbf{méfiez-vous} : la ressemblance est un signal d'alarme, pas une garantie.
\item \textbf{Vérifiez le sens dans le contexte} : la phrase a-t-elle un sens logique avec ce mot ?
\item En cas de doute, \textbf{préférez un mot simple et sûr} : \eng{to help} (aider), \eng{to join} (rejoindre, participer à), \eng{to lead} (animer), \eng{to teach people about} (sensibiliser à).
\item Tenez un \textbf{carnet personnel de faux amis} : notez chaque piège rencontré avec un exemple complet (phrase anglaise + traduction).
\end{enumerate}
\end{methode}

\begin{exoresolu}[Corrigez les erreurs]
Énoncé — Les phrases suivantes, extraites de copies d'élèves, contiennent chacune un faux ami ou un calque. Corrigez-les.
\begin{enumerate}
\item \eng{I assisted to a formation on gender equality.}
\item \eng{The actual government promotes women's rights.}
\item \eng{They animate workshops to sensitize the population.}
\item \eng{She is a very sensible girl; she cries easily.}
\item \eng{The project achieved its end in December.}
\end{enumerate}
\tcblower
Solution détaillée :
\begin{enumerate}
\item Double erreur : \eng{assisted to} $\rightarrow$ \eng{attended} ; \eng{a formation} $\rightarrow$ \eng{a training session}. Correction : \eng{I attended a training session on gender equality.} « J'ai assisté à une formation sur l'égalité des genres. »
\item \eng{Actual} signifie « réel » ; « actuel » se dit \eng{current}. Correction : \eng{The current government promotes women's rights.} « Le gouvernement actuel promeut les droits des femmes. »
\item Deux calques : \eng{animate} $\rightarrow$ \eng{run/lead} ; \eng{sensitize} $\rightarrow$ \eng{raise awareness among}. Correction : \eng{They run workshops to raise awareness among the population.} « Ils animent des ateliers pour sensibiliser la population. »
\item « Sensible » (fragile, émotive) se dit \eng{sensitive} ; \eng{sensible} signifie « raisonnable ». Correction : \eng{She is a very sensitive girl; she cries easily.} « C'est une fille très sensible ; elle pleure facilement. »
\item \eng{Achieve} ne s'emploie pas avec \eng{end} ; on \eng{achieves goals/objectives/results}. Pour « terminer », on utilise \eng{complete} ou \eng{finish}. Correction : \eng{The project was completed in December.} ou \eng{The project ended in December.} « Le projet s'est achevé en décembre. »
\end{enumerate}
\end{exoresolu}

\begin{attention}[Récapitulatif des pièges à éviter à l'écrit]
\begin{itemize}
\item Jamais de préposition après \eng{to attend} : \eng{attend a meeting}, pas \eng{attend to a meeting} (\eng{attend to} = s'occuper de).
\item \eng{To assist} + personne = aider ; jamais « assister à » un événement.
\item \eng{Actually} = « en fait, en réalité », pas « actuellement » (\eng{currently, at present}).
\item \eng{A training} est indénombrable dans ce sens : on dit \eng{two training sessions}, pas \eng{two trainings} en anglais standard.
\item \eng{To animate} = faire un dessin animé ; pour un atelier : \eng{run}, \eng{lead}, \eng{facilitate}.
\item \eng{To sensitize} = rendre allergique / hypersensible (médical) ; pour « sensibiliser » : \eng{raise awareness}.
\item \eng{Commitment} = engagement (cause) ; \eng{engagement} = fiançailles ou contrat (booking).
\end{itemize}
\end{attention}

\begin{savaistu}[Les faux amis au Bac]
Au baccalauréat, l'usage correct de \eng{to attend} et de \eng{to raise awareness} — au lieu des calques du français \eng{to assist to} et \eng{to sensitize} — est un critère de distinction dans la grille d'expression écrite. Les correcteurs repèrent immédiatement ces deux erreurs, très fréquentes chez les candidats francophones. Les éviter, c'est montrer une maîtrise réelle de l'anglais authentique : un point fort qui peut faire basculer votre note de « bien » à « très bien ».
\end{savaistu}

\begin{exercice}[À vous de jouer]
Traduisez en anglais, en évitant soigneusement les faux amis et calques :
\begin{enumerate}
\item J'ai assisté à un atelier sur les droits des femmes samedi dernier.
\item Son engagement envers la communauté est remarquable.
\item Les bénévoles animent des séances de formation dans les villages.
\item La présidente actuelle de l'association est très sensible aux questions de genre.
\item Le projet a atteint tous ses objectifs cette année.
\item Actuellement, nous sensibilisons les parents à l'éducation des filles.
\item Il a achevé son stage au mois de juin.
\end{enumerate}
\end{exercice}

\begin{corrige}[Exercice « À vous de jouer »]
\begin{enumerate}
\item \eng{I attended a workshop on women's rights last Saturday.}
\item \eng{Her commitment to the community is remarkable.}
\item \eng{The volunteers run training sessions in the villages.}
\item \eng{The current president of the association is very sensitive to gender issues.}
\item \eng{The project achieved all its goals this year.}
\item \eng{Currently, we are raising awareness among parents about girls' education.}
\item \eng{He completed his internship in June.}
\end{enumerate}
\end{corrige}

\begin{aretenir}[L'essentiel de la section]
Les faux amis sont des mots anglais qui ressemblent au français mais ont un sens différent : \eng{to assist} = aider ($\neq$ assister à = \eng{to attend}) ; \eng{sensible} = raisonnable ($\neq$ sensible = \eng{sensitive}) ; \eng{actual} = réel ($\neq$ actuel = \eng{current}) ; \eng{a formation} $\neq$ une formation (\eng{a training}) ; \eng{an engagement} = fiançailles ($\neq$ un engagement = \eng{a commitment}) ; \eng{to achieve} = atteindre ($\neq$ achever = \eng{to complete/finish}) ; \eng{actually} = en fait ($\neq$ actuellement = \eng{currently}). Les calques comme \eng{to animate a workshop} et \eng{to sensitize} doivent être remplacés par \eng{to run/lead a workshop} et \eng{to raise awareness}. En cas de doute, choisissez toujours le mot simple et sûr.
\end{aretenir}

\coursec{Mots de la même famille, prononciation et orthographe}

Après avoir débusqué les faux amis et les calques du français, il reste une dernière étape pour véritablement \cle{maîtriser} le vocabulaire du volontariat et de l'égalité des sexes : apprendre à \cle{former} les mots à partir d'une même racine, à les \cle{prononcer} correctement et à les \cle{orthographier} sans faute. C'est une excellente nouvelle pour vous : l'anglais, comme le français, construit ses familles de mots avec des préfixes et des suffixes. Si vous connaissez \eng{equal}, vous pouvez deviner \eng{equality}, \eng{unequal}, \eng{inequality} et \eng{equally}. Une seule racine apprise, cinq mots gagnés !

\courssub{Observer avant de mémoriser : un texte témoin}

Lisez d'abord ce court texte et soulignez mentalement tous les mots qui semblent appartenir à la même « famille ».

\begin{textebox}[Volunteering for equality in Bamenda]
Every Saturday, a group of young \textbf{volunteers} meets at the community hall in Bamenda. Their work is entirely \textbf{voluntary}: nobody pays them, and nobody forces them. Their main \textbf{activity} is to fight \textbf{discrimination} against girls in schools. ``Gender inequality remains a major problem in our region,'' says their leader, a committed \textbf{activist}. ``Girls and boys are \textbf{equal}, and they must be treated \textbf{equally}. Discrimination against women is illegal, yet it still happens every day.'' The group organises workshops on women's \textbf{empowerment}. ``Women's empowerment is our main goal,'' she adds. ``We want to \textbf{empower} girls through education, because education is the most \textbf{powerful} weapon against injustice. Our \textbf{commitment} is total, and we will \textbf{act} until real \textbf{equality} is achieved.''
\end{textebox}

\textbf{Glossaire :} \emph{volunteer} (n., volontaire, bénévole) ; \emph{voluntary} (adj., bénévole, volontaire) ; \emph{discrimination} (n., discrimination) ; \emph{activist} (n., militant(e)) ; \emph{empowerment} (n., autonomisation) ; \emph{to empower} (v., autonomiser, donner du pouvoir à) ; \emph{commitment} (n., engagement) ; \emph{to achieve} (v., atteindre, réaliser).

\textbf{Observation.} Dans ce texte, combien de mots dérivent de \eng{equal} ? De \eng{volunteer} ? De \eng{power} ? Vous voyez le principe : une racine + un affixe = un nouveau mot, avec souvent une nouvelle nature grammaticale.

\courssub{Les grandes familles de mots du thème}

\begin{definition}[Famille de mots]
Une \cle{famille de mots} (\eng{word family}) est l'ensemble des mots construits sur une même racine, au moyen de préfixes (placés avant) et de suffixes (placés après). Chaque membre de la famille a sa propre nature grammaticale : nom, verbe, adjectif ou adverbe.
\end{definition}

\textbf{1. La famille de EQUAL}

\begin{center}
\begin{tabularx}{\linewidth}{|l|l|X|}
\hline
\rowcolor{popblueL} Mot & Nature & Traduction \\
\hline
equal & adjectif & égal(e) \\
\hline
equality & nom & l'égalité \\
\hline
equally & adverbe & également, de façon égale \\
\hline
unequal & adjectif & inégal(e) \\
\hline
inequality & nom & l'inégalité \\
\hline
to equalise & verbe & égaliser, rendre égal \\
\hline
\end{tabularx}
\end{center}

\eng{Boys and girls have equal rights.} « Les garçons et les filles ont des droits égaux. »
\eng{The law should equalise opportunities for everyone.} « La loi devrait égaliser les chances pour tous. »

\textbf{2. La famille de VOLUNTEER}

\begin{center}
\begin{tabularx}{\linewidth}{|l|l|X|}
\hline
\rowcolor{popblueL} Mot & Nature & Traduction \\
\hline
volunteer & nom / verbe & un(e) bénévole / faire du bénévolat, se porter volontaire \\
\hline
voluntary & adjectif & bénévole, volontaire \\
\hline
voluntarily & adverbe & bénévolement, volontairement \\
\hline
involuntary & adjectif & involontaire \\
\hline
\end{tabularx}
\end{center}

\eng{She volunteers at a women's centre in Douala.} « Elle fait du bénévolat dans un centre pour femmes à Douala. »
\eng{He apologised for his involuntary mistake.} « Il s'est excusé pour son erreur involontaire. »

\textbf{3. La famille de DISCRIMINATE}

\begin{center}
\begin{tabularx}{\linewidth}{|l|l|X|}
\hline
\rowcolor{popblueL} Mot & Nature & Traduction \\
\hline
to discriminate & verbe & discriminer, faire preuve de discrimination \\
\hline
discrimination & nom & la discrimination \\
\hline
discriminatory & adjectif & discriminatoire \\
\hline
\end{tabularx}
\end{center}

Attention à la construction du verbe : \eng{to discriminate \textbf{against} somebody} (discriminer quelqu'un) et \eng{to discriminate \textbf{between} two things} (faire la distinction entre deux choses).
\eng{It is wrong to discriminate against girls.} « Il est mal de discriminer les filles. »
\eng{Discriminatory laws must be abolished.} « Les lois discriminatoires doivent être abolies. »

\textbf{4. La famille de EMPOWER}

\begin{center}
\begin{tabularx}{\linewidth}{|l|l|X|}
\hline
\rowcolor{popblueL} Mot & Nature & Traduction \\
\hline
to empower & verbe & autonomiser, donner du pouvoir à \\
\hline
empowerment & nom & l'autonomisation \\
\hline
powerful & adjectif & puissant(e) \\
\hline
powerless & adjectif & impuissant(e), sans pouvoir \\
\hline
\end{tabularx}
\end{center}

\eng{Education empowers women.} « L'éducation autonomise les femmes. »
\eng{Without education, many girls feel powerless.} « Sans éducation, beaucoup de filles se sentent impuissantes. »

\textbf{5. La famille de ACT}

\begin{center}
\begin{tabularx}{\linewidth}{|l|l|X|}
\hline
\rowcolor{popblueL} Mot & Nature & Traduction \\
\hline
to act & verbe & agir \\
\hline
action & nom & l'action \\
\hline
active & adjectif & actif(-ive) \\
\hline
activist & nom & un(e) militant(e) \\
\hline
activity & nom & une activité \\
\hline
\end{tabularx}
\end{center}

\eng{We must act now to protect women's rights.} « Nous devons agir maintenant pour protéger les droits des femmes. »
\eng{She is a well-known human rights activist.} « C'est une militante des droits humains bien connue. »

\courssub{Les suffixes qui fabriquent des noms}

\begin{propriete}[Les suffixes -tion, -ment, -ity, -ness]
Quatre suffixes permettent de transformer un verbe ou un adjectif en \cle{nom} :
\begin{itemize}
\item \textbf{-tion} (après un verbe) : \eng{promote} $\rightarrow$ \eng{promotion} ; \eng{educate} $\rightarrow$ \eng{education} ; \eng{discriminate} $\rightarrow$ \eng{discrimination}.
\item \textbf{-ment} (après un verbe) : \eng{commit} $\rightarrow$ \eng{commitment} ; \eng{empower} $\rightarrow$ \eng{empowerment} ; \eng{achieve} $\rightarrow$ \eng{achievement}.
\item \textbf{-ity} (après un adjectif) : \eng{equal} $\rightarrow$ \eng{equality} ; \eng{dignified} $\rightarrow$ \eng{dignity} ; \eng{active} $\rightarrow$ \eng{activity}.
\item \textbf{-ness} (après un adjectif) : \eng{fair} $\rightarrow$ \eng{fairness} ; \eng{selfless} $\rightarrow$ \eng{selflessness} ; \eng{aware} $\rightarrow$ \eng{awareness}.
\end{itemize}
À l'inverse, les préfixes \textbf{un-}, \textbf{in-} et \textbf{dis-} marquent le contraire : \eng{equal} $\rightarrow$ \eng{unequal} ; \eng{equality} $\rightarrow$ \eng{inequality} ; \eng{agree} $\rightarrow$ \eng{disagree} ; \eng{fair} $\rightarrow$ \eng{unfair}.
\end{propriete}

\begin{attention}[Orthographe : trois pièges à connaître par cœur]
\begin{itemize}
\item \eng{commitment} prend \textbf{deux m} : on double le \emph{t} de \eng{commit} avant d'ajouter \emph{-ment} (commit + t + ment). De même : \eng{committed}, \eng{committing}.
\item \eng{harassment} s'écrit avec \textbf{un seul r} et \textbf{deux s} : h-a-r-a-s-s-m-e-n-t. Ne copiez pas le français « harcèlement ».
\item \eng{volunteer} se termine en \textbf{-eer}, comme \eng{engineer} et \eng{career} — jamais « volonteer » ni « voluntair ».
\end{itemize}
\end{attention}

\courssub{Prononciation et accentuation : l'accent qui voyage}

En anglais, chaque mot possède une \cle{syllabe accentuée} (prononcée plus fort). Piège capital : quand on ajoute un suffixe, l'accent \emph{se déplace} souvent. C'est une différence majeure avec le français, où l'accent tombe presque toujours sur la dernière syllabe.

\begin{center}
\begin{tabularx}{\linewidth}{|l|l|X|}
\hline
\rowcolor{popblueL} Mot & Accentuation (syllabe forte en majuscules) & Prononciation approchée (français) \\
\hline
volunteer & volunTEER & « vo-leun-TIR » \\
\hline
equal & EQual & « I-kouol » \\
\hline
equality & eQUAlity & « i-KOUO-li-ti » \\
\hline
discriminate & DIScriminate & « DISS-kri-mi-neït » \\
\hline
discrimination & discrimiNAtion & « diss-kri-mi-NEI-cheune » \\
\hline
empower & emPOWER & « èm-PA-ou-eu » \\
\hline
empowerment & emPOWERment & « èm-PA-ou-eu-mente » \\
\hline
advocacy & ADvocacy & « AD-vo-keu-ssi » \\
\hline
\end{tabularx}
\end{center}

\begin{methode}[Retenir la prononciation en trois étapes]
\begin{enumerate}
\item Marquez la syllabe accentuée en majuscules dans votre cahier : volunTEER, eQUAlity, discrimiNAtion, emPOWERment, ADvocacy.
\item Répétez chaque mot \textbf{trois fois à voix haute}, en frappant dans vos mains sur la syllabe accentuée.
\item Réutilisez le mot dans une phrase complète : l'accent se mémorise mieux dans un contexte que dans une liste isolée.
\end{enumerate}
\end{methode}

\courssub{L'arbre de dérivation : visualiser la famille de EQUAL}

\begin{popfigure}
\begin{center}
\begin{tikzpicture}[
  grow=down,
  level distance=1.6cm,
  level 1/.style={sibling distance=3.4cm},
  every node/.style={draw, rounded corners, align=center, font=\small, inner sep=3pt},
  edge from parent/.style={draw=popblue, thick}
]
\node[fill=popblue, text=white, font=\bfseries, align=center]{EQUAL\\ \footnotesize égal}
  child { node[fill=popgreenL, align=center]{equality (n)\\ \footnotesize l'égalité} }
  child { node[fill=poporangeL, align=center]{unequal (adj)\\ \footnotesize inégal} }
  child { node[fill=poppinkL, align=center]{inequality (n)\\ \footnotesize l'inégalité} }
  child { node[fill=popgoldL, align=center]{equally (adv)\\ \footnotesize également} };
\end{tikzpicture}
\end{center}
\legende{Arbre de dérivation de la racine EQUAL : chaque branche ajoute un affixe et change la nature du mot.}
\end{popfigure}

\begin{savaistu}[D'où vient le « em- » de empower ?]
Le préfixe anglais \emph{em-} (ou \emph{en-}) signifie « mettre dans ». Ainsi, \eng{empower} signifie littéralement « mettre dans le pouvoir », c'est-à-dire \emph{donner le pouvoir à quelqu'un} — d'où la traduction « autonomiser ». Le même préfixe se retrouve dans \eng{encourage} (mettre du courage), \eng{enrich} (mettre dans la richesse, enrichir) et \eng{enable} (rendre capable). Ce préfixe est un héritage du vieux français, entré en anglais après la conquête normande de 1066 !
\end{savaistu}

\begin{exoresolu}[Complétez avec le bon membre de la famille]
Énoncé. Complétez chaque phrase avec la forme correcte du mot entre parenthèses.
\begin{enumerate}
\item Men and women should be paid \dots\ for the same job. (equal)
\item \dots\ against women is forbidden by law. (discriminate)
\item The association fights for women's \dots\ in rural areas. (power)
\item She works as a \dots\ at the community centre. (voluntary)
\item Gender \dots\ is still visible in many schools. (equal)
\end{enumerate}
\tcblower
Solution détaillée.
\begin{enumerate}
\item \textbf{equally} : il faut un adverbe pour modifier le verbe \eng{paid} (payés \emph{de façon égale}).
\item \textbf{Discrimination} : il faut un nom, sujet de la phrase ; on ajoute le suffixe \emph{-tion} au verbe.
\item \textbf{empowerment} : il faut un nom après le possessif \eng{women's} ; \eng{power} devient le verbe \eng{empower}, puis le nom \eng{empowerment} avec \emph{-ment}.
\item \textbf{volunteer} : il faut un nom de personne après l'article \eng{a} ; \eng{voluntary} est un adjectif et ne convient pas ici.
\item \textbf{inequality} : le contexte (« still visible », un problème) exige le contraire d'\eng{equality} ; on ajoute le préfixe \emph{in-}.
\end{enumerate}
\end{exoresolu}

\begin{attention}[Les erreurs typiques des francophones]
\begin{itemize}
\item Ne dites pas \eng{the equality of chances} : le calque du français est fautif. Dites \eng{equal opportunities} ou \eng{equal chances}.
\item \eng{Volunteer} comme nom désigne la \emph{personne} ; le \emph{travail} se dit \eng{voluntary work} ou \eng{volunteering}, jamais « a volunteering » au sens de « un bénévolat ».
\item Attention au pluriel : \eng{activity} devient \eng{activities} (le \emph{y} devient \emph{ies}), comme \eng{equality} qui reste d'ailleurs presque toujours au singulier.
\item Ne prononcez pas \eng{volunteer} à la française avec l'accent sur la première syllabe : c'est volun-\textbf{TEER}, sur la dernière !
\end{itemize}
\end{attention}

\begin{exemplebox}[Les mots de la leçon en phrases]
\eng{Gender inequality remains a major problem.} « L'inégalité des sexes reste un problème majeur. »
\eng{Discrimination against women is illegal.} « La discrimination à l'égard des femmes est illégale. »
\eng{Women's empowerment is our main goal.} « L'autonomisation des femmes est notre objectif principal. »
\eng{All citizens should be treated equally before the law.} « Tous les citoyens devraient être traités de façon égale devant la loi. »
\eng{Their commitment to voluntary work is admirable.} « Leur engagement dans le travail bénévole est admirable. »
\end{exemplebox}

\begin{exercice}[À vous de jouer]
\begin{enumerate}
\item Construisez la famille complète de \eng{fair} (adjectif, nom, adjectif contraire, nom contraire) et de \eng{power} (verbe en \emph{em-}, nom en \emph{-ment}, adjectif positif, adjectif négatif).
\item Soulignez la syllabe accentuée : volunteer, equality, discrimination, empowerment, activist, illegal.
\item Traduisez en anglais : « Les bénévoles luttent contre les lois discriminatoires et promeuvent l'égalité des chances. »
\end{enumerate}
\end{exercice}

\begin{corrige}[Exercice « À vous de jouer »]
\begin{enumerate}
\item \eng{fair} (adj.) $\rightarrow$ \eng{fairness} (n.) $\rightarrow$ \eng{unfair} (adj. contraire) $\rightarrow$ \eng{unfairness} (n. contraire). \eng{power} $\rightarrow$ \eng{to empower} (v.) $\rightarrow$ \eng{empowerment} (n.) $\rightarrow$ \eng{powerful} / \eng{powerless} (adj.).
\item volun\textbf{TEER} ; e\textbf{QUA}lity ; discrimi\textbf{NA}tion ; em\textbf{POWER}ment ; \textbf{AC}tivist ; i\textbf{LLE}gal.
\item \eng{Volunteers fight against discriminatory laws and promote equal opportunities.}
\end{enumerate}
\end{corrige}

\begin{aretenir}[L'essentiel de la section]
Une racine anglaise se décline en famille grâce aux suffixes de noms (\emph{-tion}, \emph{-ment}, \emph{-ity}, \emph{-ness}) et aux préfixes du contraire (\emph{un-}, \emph{in-}, \emph{dis-}). En ajoutant un suffixe, l'accent tonique se déplace souvent : \eng{EQual} mais \eng{eQUAlity}. Enfin, surveillez l'orthographe des mots pièges : \eng{commitment} (deux m), \eng{harassment} (un r, deux s), \eng{volunteer} (en \emph{-eer}). Maîtriser ces mécanismes, c'est multiplier votre vocabulaire par cinq sans effort supplémentaire.
\end{aretenir}

\coursec{Le lexique en contexte : lecture et production guidée}

Dans la section précédente, nous avons étudié les familles de mots (\eng{equal, equality, unequal}), la prononciation (\eng{volunteer} « vo-lon-tire ») et l'orthographe (\eng{awareness}, \eng{empowerment}). Il est temps de réunir toutes ces connaissances : un vocabulaire ne vit vraiment que dans un texte. C'est exactement ce que l'on vous demande au Baccalauréat : comprendre un document authentique, y repérer le lexique du thème, puis le réemployer dans une production écrite personnelle. Suivez le parcours complet : lecture, compréhension, repérage, dialogue, production.

\courssub{Texte support : lecture guidée}

\begin{methode}[Lire efficacement un texte au Bac]
\begin{enumerate}
\item Lisez les questions \textbf{avant} le texte : elles vous disent ce que vous cherchez.
\item Repérez les mots du champ lexical (\eng{volunteer, equality, awareness, workshop, campaign}) : ils signalent les paragraphes clés.
\item Justifiez chaque réponse par une \textbf{courte citation} du texte entre guillemets.
\item Ne traduisez jamais mot à mot : cherchez le sens global de la phrase.
\end{enumerate}
\end{methode}

\begin{textebox}[A Volunteer in Kumbo]
Every weekend, Ngong, 19, volunteers for a local NGO that promotes gender equality. With her team, she runs workshops in schools to raise awareness about girls' education and to challenge stereotypes. Last year, their campaign reached out to more than 500 families and helped 30 girls return to school. “Volunteering is unpaid, but it gives a voice to women and it makes a real difference,” she says with pride. Ngong first joined the organisation after attending a community meeting in her village. At the beginning, some parents refused to send their daughters to the workshops, but Ngong and her teammates visited them at home and explained the benefits of education. Today, she is training five new volunteers. “Empowering women is not a job for one person,” she adds. “It is a duty for the whole community.”
\end{textebox}

\begin{definition}[Glossaire du texte]
\begin{itemize}
\item \cle{to run a workshop} : animer un atelier
\item \cle{to reach out to} : aller vers, toucher (un public)
\item \cle{unpaid} : non rémunéré
\item \cle{to challenge stereotypes} : combattre les stéréotypes
\item \cle{to give a voice to} : donner la parole à
\item \cle{a duty} : un devoir
\end{itemize}
\end{definition}

\courssub{Compréhension du texte}

\begin{exoresolu}[Compréhension écrite — modèle Bac]
\textbf{A. True or False? Justify with a quotation.}
\begin{enumerate}
\item Ngong is paid for her work.
\item The campaign helped girls go back to school.
\end{enumerate}
\textbf{B. Choose the correct answer.}
Ngong trains \textbf{a)} five teachers \textbf{b)} five new volunteers \textbf{c)} 500 families.
\textbf{C. Answer in a short sentence.}
What did the campaign achieve last year?
\tcblower
\textbf{A.} 1. \textbf{False} — “Volunteering is unpaid.” 2. \textbf{True} — “helped 30 girls return to school.”
\textbf{B.} b) five new volunteers (“she is training five new volunteers”).
\textbf{C.} It reached out to more than 500 families and helped 30 girls return to school.
\end{exoresolu}

\begin{exemplebox}[Questions types et réponses modèles]
\eng{Why did some parents refuse to send their daughters?} « Pourquoi certains parents refusaient-ils d'envoyer leurs filles ? » — \eng{Because of stereotypes about girls' education.} « À cause des stéréotypes sur l'éducation des filles. »

\eng{What does Ngong do at the weekend?} « Que fait Ngong le week-end ? » — \eng{She volunteers for a local NGO and runs workshops in schools.} « Elle fait du bénévolat pour une ONG locale et anime des ateliers dans les écoles. »
\end{exemplebox}

\courssub{Repérage lexical dans le texte}

\begin{exercice}[Soulignez dans le texte]
\begin{enumerate}
\item Trois \textbf{collocations} (ex. verbe + nom).
\item Deux \textbf{mots de la même famille} que \eng{equal} ou \eng{volunteer}.
\end{enumerate}
\end{exercice}

\begin{corrige}[Repérage lexical]
\textbf{Collocations} : \eng{runs workshops} (animer des ateliers), \eng{raise awareness} (sensibiliser), \eng{makes a real difference} (faire une vraie différence). On pouvait aussi citer \eng{challenge stereotypes} et \eng{reached out to families}.

\textbf{Familles de mots} : \eng{volunteers} (verbe) et \eng{volunteering} (nom) ; \eng{gender equality} dérive de \eng{equal} ; \eng{empowering} appartient à la famille de \eng{power} (\eng{power, empower, empowerment}).
\end{corrige}

\begin{center}
\begin{tabularx}{\linewidth}{|X|X|X|}
\hline
\rowcolor{popblueL} Collocation & Traduction & Emploi dans le texte \\
\hline
\eng{run a workshop} & animer un atelier & \eng{she runs workshops in schools} \\
\hline
\eng{raise awareness} & sensibiliser & \eng{to raise awareness about girls' education} \\
\hline
\eng{make a difference} & faire une différence & \eng{it makes a real difference} \\
\hline
\eng{reach out to} & aller vers & \eng{reached out to more than 500 families} \\
\hline
\eng{challenge stereotypes} & combattre les stéréotypes & \eng{to challenge stereotypes} \\
\hline
\end{tabularx}
\end{center}

\courssub{Dialogue modèle : un entretien avec un coordinateur d'ONG}

\begin{textebox}[Interview at the NGO office]
\textbf{Coordinator:} Good morning. Why do you want to join our organisation?

\textbf{Student:} Good morning, sir. Because I want to defend women's rights and make a difference in my community.

\textbf{Coordinator:} Have you ever done any volunteer work before?

\textbf{Student:} Yes. Last year, I helped to raise awareness about girls' education in my school club.

\textbf{Coordinator:} Very good. Are you ready to work without pay, at weekends?

\textbf{Student:} Of course. Volunteering is unpaid, but it gives a voice to women. I am also ready to run workshops and to reach out to families in the villages.

\textbf{Coordinator:} Excellent. Welcome to the team! You will start next Saturday.
\end{textebox}

\begin{exemplebox}[Phrases clés du dialogue, traduites]
\eng{Why do you want to join our organisation?} « Pourquoi voulez-vous rejoindre notre organisation ? »

\eng{Have you ever done any volunteer work before?} « Avez-vous déjà fait du bénévolat auparavant ? » (present perfect + \eng{ever} pour l'expérience)

\eng{Are you ready to work without pay?} « Êtes-vous prêt à travailler sans salaire ? »

\eng{Welcome to the team!} « Bienvenue dans l'équipe ! »
\end{exemplebox}

\begin{experience}[Jeu de rôle : l'entretien du volontaire]
Par groupes de deux : l'un joue le coordinateur d'une ONG de Bamenda, l'autre un élève de Tle qui veut devenir volontaire. Le coordinateur pose au moins quatre questions (\eng{Why...? Have you ever...? Are you ready to...? What can you do?}). Le candidat répond en utilisant au moins cinq mots du lexique de la leçon (\eng{volunteer, gender equality, raise awareness, empower, make a difference, unpaid}). Puis inversez les rôles.
\end{experience}

\courssub{Production guidée : “Why I want to volunteer for gender equality”}

\begin{methode}[Rédiger un paragraphe argumenté en 3 étapes]
\begin{enumerate}
\item \textbf{Introduction} : exprimez votre motivation (\eng{I want to volunteer because...}).
\item \textbf{Body} : deux actions concrètes avec des collocations (\eng{raise awareness, empower women, run workshops}).
\item \textbf{Conclusion} : le résultat attendu (\eng{make a difference, give a voice to women}).
\end{enumerate}
Exigences : 8 à 10 lignes, au moins \textbf{6 mots du lexique} et \textbf{2 collocations}.
\end{methode}

\begin{popfigure}
\begin{tikzpicture}[node distance=0.9cm, every node/.style={align=center}]
\node[draw=popblue, fill=popblueL, rounded corners, thick, minimum width=3.4cm, minimum height=1.3cm, align=center] (intro){\textbf{Introduction}\\ \emph{motivation}\\ \eng{I want to volunteer...}};
\node[draw=poporange, fill=poporangeL, rounded corners, thick, minimum width=3.4cm, minimum height=1.3cm, right=of intro, align=center] (body){\textbf{Body}\\ \emph{actions}\\ \eng{raise awareness}\\ \eng{empower women}};
\node[draw=popgreen, fill=popgreenL, rounded corners, thick, minimum width=3.4cm, minimum height=1.3cm, right=of body, align=center] (conc){\textbf{Conclusion}\\ \emph{résultat}\\ \eng{make a difference}};
\draw[-{Stealth}, very thick, popdark] (intro) -- (body);
\draw[-{Stealth}, very thick, popdark] (body) -- (conc);
\end{tikzpicture}
\legende{Plan en trois étapes du paragraphe « Why I want to volunteer for gender equality ».}
\end{popfigure}

\begin{exemplebox}[Paragraphe modèle]
\eng{I want to volunteer for gender equality because every girl in my community deserves a chance to study. Too many girls leave school early, and I believe volunteering can change this situation. First, I would like to raise awareness about girls' education by running workshops in schools and villages. Then, I want to empower women by helping them to speak for themselves and to defend their rights. Volunteering is unpaid, but it gives a voice to women who are often silent. If we all lend a helping hand, we can challenge stereotypes together. I am sure that my small actions will make a real difference in my community.}

« Je veux faire du bénévolat pour l'égalité des genres parce que chaque fille de ma communauté mérite une chance d'étudier. Trop de filles quittent l'école tôt, et je crois que le bénévolat peut changer cette situation. D'abord, je voudrais sensibiliser à l'éducation des filles en animant des ateliers dans les écoles et les villages. Ensuite, je veux autonomiser les femmes en les aidant à s'exprimer et à défendre leurs droits. Le bénévolat n'est pas rémunéré, mais il donne la parole aux femmes souvent réduites au silence. Si nous donnons tous un coup de main, nous pouvons combattre les stéréotypes ensemble. Je suis sûr que mes petites actions feront une vraie différence dans ma communauté. »

Lexique réemployé : \eng{volunteer, gender equality, community, raise awareness, run workshops, empower women, unpaid, give a voice to, lend a helping hand, challenge stereotypes, make a difference}.
\end{exemplebox}

\courssub{Stratégie de réemploi du lexique dans une copie de Bac}

\begin{methode}[Réemployer le lexique dans un essai de Bac]
\begin{enumerate}
\item \textbf{Introduction} : un mot du champ lexical dès la première phrase (\eng{Gender equality is one of the major challenges of our society.}).
\item \textbf{Argument 1 illustré} : idée + collocation + exemple concret (\eng{Volunteers can raise awareness by running workshops, as Ngong does in Kumbo.}).
\item \textbf{Argument 2 illustré} : variez le vocabulaire (\eng{Empowering women also means giving them a voice in decision-making.}).
\item \textbf{Conclusion} : une expression idiomatique forte (\eng{If everyone lends a helping hand, we will make a real difference.}).
\end{enumerate}
\end{methode}

\begin{attention}[Ne répétez pas “help” dix fois]
Dans vos productions, variez le lexique : \eng{to help} est correct mais pauvre s'il revient sans cesse. Utilisez selon le contexte \eng{to support} (soutenir), \eng{to empower} (autonomiser), \eng{to assist} (assister), \eng{to lend a helping hand} (donner un coup de main). De même, évitez le calque \eng{*to make volunteering} : on dit \eng{to do volunteer work} ou \eng{to volunteer}. Et rappelez-vous : \eng{a volunteer} est un bénévole, jamais « un volontaire » au sens militaire.
\end{attention}

\begin{exercice}[À vous d'écrire]
Rédigez votre propre paragraphe \eng{Why I want to volunteer for gender equality} (8 à 10 lignes). Utilisez au moins six mots du lexique de la leçon et deux collocations différentes de celles du modèle. Soulignez-les dans votre copie pour vérification.
\end{exercice}

\begin{aretenir}[L'essentiel de la section]
\begin{itemize}
\item Lisez les questions avant le texte et justifiez chaque réponse par une citation.
\item Les collocations (\eng{run a workshop, raise awareness, make a difference}) donnent du naturel à votre anglais.
\item Structurez toute production en trois temps : motivation, actions, résultat.
\item Variez le lexique et méfiez-vous des calques du français.
\end{itemize}
\end{aretenir}

\newpage

\coursec{Activité d'intégration}

\coursec{Lesson 39 — Vocabulary: Words and expressions related to volunteering in gender equality promotion activities}

\courssub{Objectifs de l'activité}

À l'issue de cette activité d'intégration, l'élève sera capable de :
\begin{itemize}
\item réemployer à l'oral et à l'écrit le vocabulaire du \cle{volontariat} (\eng{volunteer, community service, commitment, to get involved}) et de \cle{l'égalité des sexes} (\eng{gender equality, women's rights, to empower, discrimination}) dans une situation de communication authentique ;
\item produire une courte \cle{lettre de motivation} (\eng{cover letter}) correcte, utilisant au moins huit mots du lexique, trois collocations et un mot de la même famille que \eng{equal} ou \eng{volunteer} ;
\item conduire et passer un \cle{entretien de recrutement} (\eng{job interview}) en anglais, avec une prononciation soignée ;
\item évaluer la production d'un camarade à l'aide d'une grille simple et justifier un choix en anglais ;
\item éviter les \cle{faux amis} et les calques du français (\eng{to assist} ne veut pas dire « assister à », \eng{sensible} n'est pas « sensible »).
\end{itemize}

\courssub{Situation}

Le \emph{Gender Equality Club} du lycée de Mfoundi, à Yaoundé, travaille en partenariat avec une ONG britannique, \emph{Equal Chances UK}, qui finance chaque année des projets scolaires de promotion de l'égalité entre les filles et les garçons : clubs de lecture pour les filles, campagnes contre le harcèlement, séances de sensibilisation dans les quartiers. Cette année, l'ONG recherche \textbf{deux volontaires élèves de Terminale} pour coordonner la campagne \emph{“Girls in School, Bright Futures”}. Une \emph{Recruitment Day} est organisée : les candidats envoient d'abord une lettre de motivation, puis passent un entretien en anglais devant les deux coordinateurs de l'ONG.

Voici l'annonce officielle affichée au tableau du club :

\begin{textebox}[Equal Chances UK — Volunteer Recruitment Notice]
\textbf{VOLUNTEERS NEEDED!}\par
Equal Chances UK, in partnership with the Gender Equality Club of Lycée du Mfoundi, Yaoundé, is recruiting \textbf{two student volunteers} for the campaign “Girls in School, Bright Futures”.\par
\textbf{Your mission:} raise awareness about girls' education, help organise workshops, and empower young women in your community.\par
\textbf{We are looking for} committed, reliable and open-minded students who want to make a difference and fight against gender discrimination.\par
\textbf{How to apply:} send a short cover letter (about 10 lines) explaining your motivation, then attend a two-minute interview in English.\par
\textit{Deadline: Friday, 5 p.m. — Equal opportunities for boys and girls!}
\end{textebox}

\begin{definition}[Glossaire de l'annonce]
\begin{itemize}
\item \textbf{committed} (adj.) : engagé, dévoué
\item \textbf{reliable} (adj.) : fiable, sur qui on peut compter
\item \textbf{open-minded} (adj.) : ouvert d'esprit, tolérant
\item \textbf{workshop} (n.) : atelier, séance de travail pratique
\item \textbf{deadline} (n.) : date limite, échéance
\item \textbf{raise awareness} (loc. v.) : sensibiliser, faire prendre conscience
\item \textbf{empower} (v.) : donner du pouvoir à, rendre autonome, renforcer
\item \textbf{make a difference} (loc. v.) : changer les choses, avoir un impact
\end{itemize}
\end{definition}

\courssub{Consignes}

Travail par groupes de quatre : deux élèves jouent les \textbf{coordinateurs} de l'ONG, deux élèves jouent les \textbf{candidats}.

\begin{enumerate}
\item \textbf{Compréhension de l'annonce.} Lisez l'avis de recrutement. Relevez : (a) deux collocations du lexique de la leçon, (b) trois qualités attendues, (c) la mission des volontaires. Les coordinateurs vérifient que tout le groupe a compris la mission.
\item \textbf{Préparation de l'entretien (coordinateurs).} Rédigez cinq questions d'entretien en anglais. Exemples : \eng{Why do you want to volunteer for this campaign?} ; \eng{What qualities do you have?} ; \eng{Have you ever taken part in community service?} ; \eng{How would you raise awareness in your neighbourhood?} ; \eng{What does gender equality mean to you?}
\item \textbf{Lettre de motivation (candidats).} Rédigez individuellement une lettre d'environ dix lignes. Elle doit contenir : au moins \textbf{huit mots du lexique} de la leçon, \textbf{trois collocations} (par exemple \eng{raise awareness, empower women, make a difference, fight discrimination, promote gender equality}) et \textbf{un mot de la même famille} que \eng{equal} (\eng{equality, unequal, equally}) ou \eng{volunteer} (\eng{voluntary, voluntarily, volunteering}).
\item \textbf{Relecture croisée.} Échangez les lettres dans le groupe. Soulignez les mots du lexique employés, corrigez les faux amis et les calques du français, vérifiez l'orthographe (\eng{volunteer} avec un \emph{-eer}, \eng{committee} avec deux \emph{m} et deux \emph{t}).
\item \textbf{Entretien oral.} Chaque candidat passe un entretien de deux minutes. Soignez la prononciation : \eng{volunteer} se prononce « vo-lon-\textbf{tier} » (accent sur la dernière syllabe), \eng{women} se prononce « oui-min ».
\item \textbf{Évaluation et délibération.} Les coordinateurs remplissent la grille ci-dessous, puis annoncent le volontaire retenu en justifiant leur choix en anglais : \eng{We have selected ... because her letter showed real commitment and she used the vocabulary correctly.}
\end{enumerate}

\begin{center}
\begin{tabularx}{\linewidth}{|X|c|c|c|}
\hline
\rowcolor{popblueL} Critère & Excellent (2) & Correct (1) & À revoir (0) \\
\hline
Lexique du thème (8 mots minimum) & & & \\
\hline
Collocations correctes (3 minimum) & & & \\
\hline
Prononciation et fluidité & & & \\
\hline
Absence de faux amis et de calques & & & \\
\hline
\end{tabularx}
\end{center}

\courssub{Ressources à mobiliser}

\begin{aretenir}[Boîte à outils de l'activité]
\begin{itemize}
\item \textbf{Lexique du volontariat :} \eng{to volunteer, a volunteer, voluntary work, community service, commitment, to get involved, to take part in, to join an NGO, to dedicate one's time}.
\item \textbf{Lexique de l'égalité :} \eng{gender equality, women's rights, to empower women, discrimination, stereotypes, equal opportunities, to promote, to defend}.
\item \textbf{Collocations :} \eng{raise awareness, make a difference, fight (against) discrimination, promote gender equality, empower women, break stereotypes, apply for a position}.
\item \textbf{Familles de mots :} \eng{equal} $\rightarrow$ \eng{equality, inequality, equally} ; \eng{volunteer} $\rightarrow$ \eng{voluntary, voluntarily, volunteering} ; \eng{discriminate} $\rightarrow$ \eng{discrimination, discriminatory}.
\item \textbf{Structures utiles :} \eng{I am writing to apply for...} ; \eng{I would like to volunteer because...} ; \eng{I am deeply committed to...} ; \eng{I believe that every girl deserves...} ; \eng{Thank you for considering my application.}
\end{itemize}
\end{aretenir}

\begin{attention}[Pièges à éviter pendant l'activité]
\begin{itemize}
\item Ne dites pas \eng{*I want to make volunteering} : dites \eng{I want to do some voluntary work} ou \eng{I want to volunteer}.
\item \eng{To assist} signifie « aider » ; « assister à un atelier » se dit \eng{to attend a workshop}.
\item \eng{Sensible} signifie « raisonnable, sensé » ; « sensible (émotif) » se dit \eng{sensitive}.
\item Pas de calque : « je suis engagé » se dit \eng{I am committed}, jamais \eng{*I am engaged} (qui signifie « fiancé »).
\item \eng{Gender} = genre (social), pas \eng{sex} (biologique) ; \eng{gender equality} = égalité entre les sexes / égalité de genre.
\end{itemize}
\end{attention}

\courssub{Proposition de production et corrigé commenté}

\begin{exoresolu}[Modèle de lettre de motivation]
Rédigez la lettre de motivation d'un candidat à la campagne “Girls in School, Bright Futures” (environ dix lignes).
\tcblower
\textbf{Modèle de production :}\par
\eng{Dear Sir or Madam,}\par
\eng{I am writing to apply for the position of student volunteer in the campaign “Girls in School, Bright Futures”. I am a final-year student at Lycée du Mfoundi and I am deeply committed to promoting gender equality in my community. Last year, I took part in a voluntary project to raise awareness about girls' education in my neighbourhood, and this experience showed me that one person can really make a difference. I believe we must fight against discrimination and break the stereotypes that prevent girls from staying in school. As a volunteer, I would like to help empower young women and defend equal opportunities for all. I am reliable, hard-working and open-minded. Thank you for considering my application. I look forward to hearing from you.}\par
\eng{Yours faithfully,}\par
\eng{Brenda Etoundi}\par
\textbf{Commentaire des choix de langue :}
\begin{itemize}
\item \textbf{Lexique du thème (plus de 8 mots) :} \eng{volunteer, committed, gender equality, voluntary project, awareness, discrimination, stereotypes, empower, equal opportunities} — tous issus des champs lexicaux de la leçon.
\item \textbf{Collocations respectées :} \eng{apply for the position}, \eng{raise awareness}, \eng{make a difference}, \eng{fight against discrimination}, \eng{break stereotypes}, \eng{empower young women}.
\item \textbf{Mot de la même famille :} \eng{voluntary} (adjectif dérivé de \eng{volunteer}) et \eng{equality} (nom dérivé de \eng{equal}).
\item \textbf{Formules de lettre formelle :} ouverture \eng{Dear Sir or Madam}, formule de politesse \eng{Yours faithfully}, structure d'introduction \eng{I am writing to apply for...}.
\item \textbf{Faux amis évités :} \eng{committed} (et non \eng{*engaged}), \eng{took part in} (et non \eng{*assisted}).
\item \textbf{Grammaire :} \eng{committed to promoting} (préposition + gérondif), \eng{help empower} (infinitif sans \eng{to} après \eng{help}), \eng{look forward to hearing} (préposition + gérondif).
\end{itemize}
\end{exoresolu}

\courssub{Bilan de l'activité}

Cette \emph{Recruitment Day} a permis de mobiliser l'ensemble des savoir-faire de la leçon dans une tâche finale authentique et motivante. Sur le plan de la \textbf{compréhension}, les élèves ont exploité un document réel (l'avis de recrutement) pour en extraire les informations utiles. Sur le plan du \textbf{lexique}, ils ont réemployé les deux champs lexicaux (volontariat et égalité des sexes), les collocations et les familles de mots dans une production écrite contrainte puis dans un échange oral spontané. Sur le plan de la \textbf{correction}, la relecture croisée et la grille d'évaluation ont développé l'attention aux faux amis, à l'orthographe et à la prononciation. Enfin, l'activité a une dimension \textbf{citoyenne} forte : elle relie l'apprentissage de l'anglais à l'éducation à l'égalité et à l'engagement communautaire, au cœur du chapitre \emph{Citizenship and Human Rights}. Une prolongation possible : afficher les meilleures lettres au tableau du club ou enregistrer les entretiens pour une auto-évaluation en classe.

\newpage

\coursec{Méthodes et exercices résolus}

\courssub{Savoir-faire 1 : Identifier et employer le vocabulaire du thème}

\begin{methode}[Identifier et employer le vocabulaire du thème]
Pour maîtriser le champ lexical du \cle{volunteering} et de la \cle{gender equality}, procède ainsi :
\begin{enumerate}
\item \textbf{Repère le thème} du texte : cherche les mots-clés (\eng{volunteer}, \eng{community service}, \eng{equal rights}, \eng{empowerment}, \eng{discrimination}).
\item \textbf{Classe les mots} en familles : personnes (\eng{a volunteer, an activist, a beneficiary}), actions (\eng{to raise awareness, to empower, to advocate}), lieux et institutions (\eng{a charity, an NGO, a community centre}).
\item \textbf{Vérifie la nature grammaticale} du mot avant de l'employer : \eng{voluntary} (adjectif) $\neq$ \eng{volunteer} (nom ou verbe) $\neq$ \eng{voluntarily} (adverbe).
\item \textbf{Choisis la bonne préposition} : \eng{to volunteer \textbf{for} an organisation}, \eng{to be involved \textbf{in} a project}, \eng{to fight \textbf{against} discrimination}, \eng{equal \textbf{to}}.
\item \textbf{Réemploie le mot dans ta propre phrase} liée au contexte camerounais : \eng{Students in Yaoundé volunteer for a local NGO that promotes girls' education.} « Des élèves de Yaoundé sont bénévoles dans une ONG locale qui promeut l'éducation des filles. »
\end{enumerate}
\end{methode}

\begin{definition}[Les mots-clés du thème]
\begin{itemize}
\item \eng{a volunteer} (nom) : un bénévole, un volontaire ; \eng{to volunteer} (verbe) : se porter volontaire, faire du bénévolat.
\item \eng{voluntary work} : le travail bénévole ; \eng{community service} : le service à la communauté.
\item \eng{gender equality} : l'égalité des sexes ; \eng{equal rights} : l'égalité des droits.
\item \eng{to empower} : autonomiser, donner les moyens d'agir ; \eng{empowerment} : l'autonomisation.
\item \eng{to raise awareness about} : sensibiliser à ; \eng{to campaign for/against} : mener campagne pour/contre.
\item \eng{discrimination} : la discrimination ; \eng{to advocate} : plaider pour, défendre.
\end{itemize}
\end{definition}

\begin{exoresolu}[Vocabulary in context — gap filling]
\textbf{Énoncé.} Complete the following passage with words from the box: \emph{volunteers, awareness, empower, gender, community, discrimination}.

« Last year, our school launched a project to promote \ldots{} equality in the \ldots{}. Twenty \ldots{} visited villages near Bafoussam to raise \ldots{} about girls' right to education. Their goal was to \ldots{} women and to fight against all forms of \ldots{}. »
\tcblower
\textbf{Solution détaillée.}
\begin{itemize}
\item \textbf{Blank 1 — gender} : la leçon porte sur l'égalité des sexes ; la collocation correcte est \eng{gender equality}. On ne dit jamais \eng{genre equality}.
\item \textbf{Blank 2 — community} : \eng{in the community} est la collocation standard pour « dans la communauté ».
\item \textbf{Blank 3 — volunteers} : le sujet de \eng{visited} doit être des personnes ; le déterminant \eng{twenty} exige le pluriel \eng{volunteers}.
\item \textbf{Blank 4 — awareness} : collocation figée \eng{to raise awareness about} = « sensibiliser à ». Attention : \eng{awareness} est indénombrable, jamais de pluriel.
\item \textbf{Blank 5 — empower} : après \eng{to} (infinitif de but), il faut la base verbale ; \eng{to empower women} = « autonomiser les femmes ».
\item \textbf{Blank 6 — discrimination} : \eng{to fight against discrimination} ; indénombrable, sans article.
\end{itemize}
\textbf{Réponse rédigée :} « Last year, our school launched a project to promote \textbf{gender} equality in the \textbf{community}. Twenty \textbf{volunteers} visited villages near Bafoussam to raise \textbf{awareness} about girls' right to education. Their goal was to \textbf{empower} women and to fight against all forms of \textbf{discrimination}. »
\textbf{Conclusion :} toujours vérifier la nature du mot attendu (nom, verbe, adjectif) et la collocation avant de remplir un blanc.
\end{exoresolu}

\courssub{Savoir-faire 2 : Réemployer le lexique à l'oral et à l'écrit}

\begin{methode}[Réemployer le lexique à l'oral et à l'écrit]
Pour produire un paragraphe ou une prise de parole sur le volontariat :
\begin{enumerate}
\item \textbf{Planifie en trois temps} : présentation de l'action (\eng{I volunteered for\ldots}), description des activités (\eng{We organised workshops, campaigned for\ldots}), bilan et sentiments (\eng{It was a rewarding experience}).
\item \textbf{Utilise des connecteurs} : \eng{first, then, moreover, as a result, finally}.
\item \textbf{Varie le lexique} : remplace \eng{help} par \eng{to support, to assist, to give a hand to} ; remplace \eng{good} par \eng{rewarding, meaningful, valuable}.
\item \textbf{Accorde les temps} : pour raconter une expérience passée, utilise le prétérit (\eng{I joined, we organised}) ; pour un fait général, le présent (\eng{Volunteering develops leadership skills}).
\item \textbf{Conclus par une opinion} avec \eng{in my opinion, I believe that\ldots} suivi d'un argument.
\end{enumerate}
\end{methode}

\begin{exemplebox}[Collocations utiles pour la composition]
\begin{center}
\begin{tabularx}{\linewidth}{|X|X|}\hline
\rowcolor{popblueL} \textbf{Collocation} & \textbf{Traduction} \\ \hline
\eng{to do voluntary work} & faire du bénévolat \\ \hline
\eng{to join an association / a club} & adhérer à une association / un club \\ \hline
\eng{to raise awareness about a problem} & sensibiliser à un problème \\ \hline
\eng{to campaign for equal pay} & mener campagne pour l'égalité salariale \\ \hline
\eng{to fight against discrimination} & lutter contre la discrimination \\ \hline
\eng{to empower women and girls} & autonomiser les femmes et les filles \\ \hline
\eng{to make a difference} & faire une différence, changer les choses \\ \hline
\eng{a rewarding experience} & une expérience enrichissante \\ \hline
\end{tabularx}
\end{center}
\end{exemplebox}

\begin{exoresolu}[Guided writing — a short paragraph]
\textbf{Énoncé.} Write a paragraph of about 80 words describing a volunteering activity you took part in to promote gender equality. Use at least five words from the lesson: \emph{volunteer, raise awareness, empower, campaign, rewarding}.
\tcblower
\textbf{Solution détaillée.}
\textbf{Étape 1 — Plan :} (a) où et quand ; (b) les actions menées ; (c) le bilan personnel.
\textbf{Étape 2 — Choix des temps :} le récit d'une expérience terminée appelle le prétérit simple : \eng{volunteered, organised, campaigned, learnt}.
\textbf{Étape 3 — Intégration des mots imposés :} vérifier que chaque mot est employé avec sa bonne construction : \eng{to volunteer \textbf{for}}, \eng{to raise awareness \textbf{about}}, \eng{to campaign \textbf{for}}.
\textbf{Paragraphe modèle :}
« Last holidays, I \textbf{volunteered} for a women's association in Douala. We organised workshops to \textbf{raise awareness} about girls' education and we \textbf{campaigned} for equal access to secondary school. I also helped to \textbf{empower} young mothers by teaching them basic computer skills. It was a truly \textbf{rewarding} experience because I learnt that small actions can change lives. In my opinion, every student should volunteer at least once. »
(83 mots — les cinq mots imposés sont employés avec les bonnes prépositions.)
\textbf{Conclusion :} un bon paragraphe = plan clair + temps corrects + lexique varié + opinion finale.
\end{exoresolu}

\courssub{Savoir-faire 3 : Éviter les faux amis et les calques du français}

\begin{methode}[Éviter les faux amis et les calques du français]
Réflexes à adopter face aux pièges de traduction dans ce thème :
\begin{enumerate}
\item \textbf{Méfie-toi des mots transparents} : \eng{to assist} signifie « aider » (assister à une réunion = \eng{to attend a meeting}) ; \eng{actually} = « en fait » (actuellement = \eng{currently}).
\item \textbf{Ne traduis jamais mot à mot} : « sensibiliser » ne se dit pas \eng{sensibilize} mais \eng{to raise awareness} ; « une œuvre de charité » = \eng{a charity} et non \eng{a charity work}.
\item \textbf{Vérifie le statut du mot} : \eng{a volunteer} est la personne ; l'adjectif est \eng{voluntary} (\eng{voluntary work} = « travail bénévole »). La forme \eng{volunteer work} existe en anglais américain, mais au Bac, privilégie la forme britannique \eng{voluntary work}.
\item \textbf{Attention aux prépositions calquées} : \eng{to participate \textbf{in}} (pas \eng{to}), \eng{to be interested \textbf{in}}, \eng{to depend \textbf{on}}, \eng{to listen \textbf{to}}.
\item \textbf{En cas de doute}, reformule avec un mot simple et sûr plutôt qu'un faux ami risqué.
\end{enumerate}
\end{methode}

\begin{attention}[Faux amis à connaître absolument]
\begin{center}
\begin{tabularx}{\linewidth}{|X|X|X|}\hline
\rowcolor{popblueL} \textbf{Français} & \textbf{English} & \textbf{Piège} \\ \hline
assister à (une réunion) & \eng{to attend (a meeting)} & \eng{to assist} = aider \\ \hline
sensibiliser & \eng{to raise awareness} & \eng{sensibilize} n'existe pas \\ \hline
actuellement & \eng{currently, at present} & \eng{actually} = en fait \\ \hline
une œuvre de charité & \eng{a charity} & pas \eng{a charity work} \\ \hline
participer à & \eng{to participate in} & pas \eng{participate to} \\ \hline
s'engager dans un projet & \eng{to get involved in a project} & \eng{involved} + \eng{in} \\ \hline
\end{tabularx}
\end{center}
\end{attention}

\begin{exoresolu}[Correct the mistakes]
\textbf{Énoncé.} Each sentence contains one mistake (false friend, wrong preposition or calque from French). Correct it.
\begin{enumerate}
\item Many students assisted the meeting on women's rights yesterday.
\item The association wants to sensibilize the population about gender violence.
\item She is volunteer in an NGO that defends children's rights.
\item We participated to a campaign for equal pay.
\end{enumerate}
\tcblower
\textbf{Solution détaillée.}
\begin{enumerate}
\item \eng{assisted} $\rightarrow$ \eng{attended}. Faux ami : \eng{to assist} = « aider » ; « assister à » (être présent) = \eng{to attend}. Phrase corrigée : \eng{Many students \textbf{attended} the meeting on women's rights yesterday.}
\item \eng{to sensibilize} $\rightarrow$ \eng{to raise awareness among}. Le verbe français « sensibiliser » n'a pas d'équivalent direct ; la collocation anglaise est \eng{to raise awareness among the population about gender violence}.
\item \eng{is volunteer} $\rightarrow$ \eng{is \textbf{a} volunteer}. \eng{Volunteer} est un nom comptable : il exige l'article indéfini. Autre solution : \eng{she \textbf{volunteers for} an NGO} (verbe).
\item \eng{participated to} $\rightarrow$ \eng{participated \textbf{in}}. Calque du français « participer à » : en anglais, la préposition est \eng{in}.
\end{enumerate}
\textbf{Conclusion :} avant d'écrire, demande-toi toujours si le mot anglais existe vraiment avec ce sens et quelle préposition il gouverne.
\end{exoresolu}

\courssub{Savoir-faire 4 : Exploiter les familles de mots, la prononciation et l'orthographe}

\begin{methode}[Familles de mots, prononciation et orthographe]
Pour réussir les exercices de formation de mots (\emph{word formation}) :
\begin{enumerate}
\item \textbf{Identifie la nature du mot manquant} grâce au contexte : après un article ou un adjectif $\rightarrow$ nom ; après un auxiliaire $\rightarrow$ verbe ; devant un nom $\rightarrow$ adjectif ; pour modifier un verbe $\rightarrow$ adverbe.
\item \textbf{Mémorise les suffixes clés} du thème : -\emph{tion} (\eng{promotion, discrimination, education}), -\emph{ment} (\eng{empowerment, development}), -\emph{ity} (\eng{equality, solidarity}), -\emph{ist} (\eng{activist, feminist}), -\emph{ship} (\eng{leadership, membership}).
\item \textbf{Attention aux changements de base} : \eng{equal} $\rightarrow$ \eng{equality} ; \eng{voluntary} $\rightarrow$ \eng{voluntarily} ; \eng{to discriminate} $\rightarrow$ \eng{discrimination}.
\item \textbf{Soigne la prononciation} : \eng{volunteer} se prononce « vo-lon-TIER » (accent sur la dernière syllabe) ; \eng{charity} « TCHA-ri-ti » ; \eng{women} se dit « oui-mine » (et non « ouo-mène »).
\item \textbf{Vérifie l'orthographe} : \eng{committee} (deux m, deux t, deux e), \eng{benefit} $\rightarrow$ \eng{benefited} (un seul t en britannique courant ; \eng{benefitted} est aussi admis), \eng{equal} $\rightarrow$ \eng{equally}.
\end{enumerate}
\end{methode}

\begin{propriete}[Familles de mots du thème]
\begin{center}
\begin{tabular}{|l|l|l|l|}\hline
\rowcolor{popblueL} \textbf{Verbe} & \textbf{Nom} & \textbf{Adjectif} & \textbf{Personne} \\ \hline
\eng{to volunteer} & \eng{volunteering} & \eng{voluntary} & \eng{a volunteer} \\ \hline
\eng{to equal} & \eng{equality} & \eng{equal} & --- \\ \hline
\eng{to empower} & \eng{empowerment} & \eng{empowered} & --- \\ \hline
\eng{to discriminate} & \eng{discrimination} & \eng{discriminatory} & --- \\ \hline
\eng{to promote} & \eng{promotion} & \eng{promotional} & \eng{a promoter} \\ \hline
\eng{to develop} & \eng{development} & \eng{developed/developing} & \eng{a developer} \\ \hline
\eng{to lead} & \eng{leadership} & \eng{leading} & \eng{a leader} \\ \hline
\eng{---} & \eng{charity} & \eng{charitable} & --- \\ \hline
\end{tabular}
\end{center}
\end{propriete}

\begin{exoresolu}[Word formation]
\textbf{Énoncé.} Complete each sentence with the correct form of the word in brackets.
\begin{enumerate}
\item \ldots{} between men and women is a fundamental human right. (equal)
\item The \ldots{} of girls through education is a priority for our NGO. (empower)
\item She works as a \ldots{} for a charity in Bamenda. (voluntary)
\item Gender \ldots{} remains a serious problem in some workplaces. (discriminate)
\end{enumerate}
\tcblower
\textbf{Solution détaillée.}
\begin{enumerate}
\item \textbf{Equality} : le mot est sujet de la phrase et suivi du verbe \eng{is} ; il faut un nom. \eng{Equal} (adjectif) $+$ suffixe -\emph{ity} $\rightarrow$ \eng{equality}. Majuscule en tête de phrase.
\item \textbf{empowerment} : après l'article \eng{the} et devant \eng{of}, il faut un nom ; le verbe \eng{empower} $+$ -\emph{ment} $\rightarrow$ \eng{empowerment} (« l'autonomisation »).
\item \textbf{volunteer} : après \eng{as a}, il faut un nom de personne. \eng{Voluntary} est un adjectif ; la personne est \eng{a volunteer}. Phrase : \eng{She works as a volunteer for a charity in Bamenda.}
\item \textbf{discrimination} : après l'adjectif \eng{gender}, il faut un nom ; \eng{to discriminate} $\rightarrow$ \eng{discrimination} (la base verbale perd le -\emph{e} final devant -\emph{ion}).
\end{enumerate}
\textbf{Conclusion :} la stratégie gagnante = repérer la nature grammaticale exigée, puis appliquer le bon suffixe en vérifiant l'orthographe de la base.
\end{exoresolu}

\begin{savaistu}[Le bénévolat, une tradition anglophone]
Dans les pays anglophones, le volontariat (\eng{volunteering}) fait partie du parcours scolaire et professionnel : de nombreux lycéens britanniques et américains accomplissent du \eng{community service} avant l'université, et les employeurs apprécient les candidats qui mentionnent des activités bénévoles dans leur CV. Au Cameroun, de nombreuses ONG (\eng{NGOs} = \emph{non-governmental organisations}) et associations de jeunes, à Yaoundé, Douala, Buea ou Bamenda, s'appuient sur des volontaires pour promouvoir l'éducation des filles et lutter contre les mariages précoces (\eng{early marriages}).
\end{savaistu}

\courssub{Savoir-faire 5 : Comprendre un texte et répondre aux questions}

\begin{methode}[Répondre aux questions de compréhension]
Face aux questions de reading du Bac sur un texte de ce thème :
\begin{enumerate}
\item \textbf{Lis d'abord les questions}, puis le texte : tu sauras quels détails repérer.
\item \textbf{Localise la réponse} dans le texte grâce aux mots-clés de la question (\eng{why, when, who, how many}).
\item \textbf{Réponds avec tes propres mots} : ne recopie jamais une longue phrase du texte ; reformule en une phrase complète avec sujet et verbe.
\item \textbf{Pour les questions de vocabulaire} (\emph{Find in the text a word meaning\ldots}), donne le mot exact du texte, sans le modifier.
\item \textbf{Pour les questions vrai/faux}, justifie toujours par une courte citation entre guillemets anglais “ \ldots ”.
\end{enumerate}
\end{methode}

\begin{exoresolu}[Reading comprehension]
\textbf{Énoncé.} Read the text and answer the questions in your own words.

\emph{« In 2022, a group of students from Government High School Limbe created the “Equal Future Club”. Every Saturday, the members volunteer in nearby villages. They organise talks on girls' education, distribute books and campaign against early marriage. “Volunteering has taught me leadership and respect for others,” says Clarisse, the club president. The club now has fifty members, both boys and girls. »}

\begin{enumerate}
\item When was the Equal Future Club created?
\item What do the members do every Saturday? (Mention two activities.)
\item Find in the text a word meaning “a person who works without being paid”.
\item True or false? Justify: Only girls can join the club.
\end{enumerate}
\tcblower
\textbf{Solution détaillée.}
\begin{enumerate}
\item La question porte sur la date : le texte dit \eng{in 2022}. Réponse complète : \eng{The Equal Future Club was created in 2022.}
\item Deux activités parmi celles listées : \eng{Every Saturday, the members volunteer in nearby villages: they organise talks on girls' education and distribute books.} (On pouvait aussi citer \eng{campaign against early marriage}.) Reformulation acceptée, verbes au présent car il s'agit d'une habitude (\eng{every Saturday}).
\item Le mot exact du texte : \eng{volunteer} (dans \eng{the members volunteer}). On ne change ni la forme ni l'orthographe.
\item \textbf{False.} Justification par citation : the text says “the club now has fifty members, \textbf{both boys and girls}”, so boys can join too.
\end{enumerate}
\textbf{Conclusion :} une réponse complète, reformulée, avec les bons temps et une justification citée entre guillemets, rapporte tous les points.
\end{exoresolu}

\begin{attention}[Les 5 erreurs qui coûtent des points au Bac]
\begin{enumerate}
\item \textbf{Faux ami \eng{to assist}} : écrire \eng{I assisted the workshop} pour « j'ai assisté à l'atelier » est faux ; dites \eng{I \textbf{attended} the workshop}. \eng{To assist} = « aider ».
\item \textbf{Calque « sensibiliser »} : \eng{sensibilize} n'existe pas en anglais courant ; la seule formule correcte est \eng{to \textbf{raise awareness} about}.
\item \textbf{Article oublié devant \eng{volunteer}} : \eng{She is volunteer} est incorrect ; nom comptable oblige : \eng{She is \textbf{a} volunteer} ou bien le verbe \eng{She \textbf{volunteers for}\ldots}
\item \textbf{Prépositions calquées} : \eng{participate \textbf{to}}, \eng{listen the radio}, \eng{depend \textbf{of}} sont fautifs ; retenez \eng{participate \textbf{in}}, \eng{listen \textbf{to}}, \eng{depend \textbf{on}}, \eng{involved \textbf{in}}, \eng{fight \textbf{against}}.
\item \textbf{Orthographe et prononciation} : \eng{volonteer} (faute : \eng{volunteer}), \eng{commitee} (faute : \eng{committee}), et \eng{women} prononcé « ouo-mène » au lieu de « oui-mine » pénalisent l'écrit comme l'oral.
\end{enumerate}
\end{attention}

\begin{exercice}[Pour s'entraîner]
\begin{enumerate}
\item Complete with the correct preposition: \eng{She volunteers \ldots{} a charity. They are involved \ldots{} a project. We fight \ldots{} discrimination. He depends \ldots{} his team.}
\item Put the words in brackets in the correct form: \eng{\ldots{} (equal) of opportunity is essential. The \ldots{} (promote) of women's rights is our goal. She is a committed \ldots{} (activism).}
\item Translate into English: « Les élèves de notre lycée font du bénévolat pour sensibiliser la population à l'égalité des sexes. »
\item Write five sentences about a volunteering activity in your town, using \eng{volunteer, campaign, raise awareness, empower, rewarding}.
\end{enumerate}
\end{exercice}

\begin{corrige}[Exercice « Pour s'entraîner »]
\begin{enumerate}
\item \eng{She volunteers \textbf{for} a charity. They are involved \textbf{in} a project. We fight \textbf{against} discrimination. He depends \textbf{on} his team.}
\item \eng{\textbf{Equality} of opportunity is essential. The \textbf{promotion} of women's rights is our goal. She is a committed \textbf{activist}.}
\item \eng{The students of our high school do voluntary work to raise awareness among the population about gender equality.} (Variante acceptée : \eng{Our school students volunteer to raise awareness about gender equality in the community.})
\item Exemple de production : \eng{Last term, I volunteered for a youth club in Bafoussam. We campaigned for girls' education and organised debates to raise awareness about early marriage. I helped to empower young girls by tutoring them in English. It was a rewarding experience that taught me leadership.}
\end{enumerate}
\end{corrige}

\newpage

\coursec{Exercices}

\courssub{Je vérifie mes connaissances}

\begin{exercice}[QCM : le bon mot]
Pour chaque phrase, choisis la bonne réponse (a, b ou c).
\begin{enumerate}
\item Many young people in Yaoundé \cle{volunteer} to help girls stay in school.
\begin{itemize}
\item[a.] volunteer \item[b.] volonteer \item[c.] volontar
\end{itemize}
\item The NGO launched an \cle{awareness} campaign against early marriage.
\begin{itemize}
\item[a.] awareness \item[b.] awarness \item[c.] awaireness
\end{itemize}
\item She has a strong \cle{commitment} to women's rights.
\begin{itemize}
\item[a.] committment \item[b.] commitment \item[c.] comitment
\end{itemize}
\item The association fights \cle{discrimination} against girls in rural areas.
\begin{itemize}
\item[a.] discrimination \item[b.] discrimmination \item[c.] descrimination
\end{itemize}
\end{enumerate}
\end{exercice}

\begin{exercice}[Vrai ou faux ?]
Dis si les affirmations suivantes sont vraies (V) ou fausses (F), puis justifie ta réponse en une phrase complète en anglais.
\begin{enumerate}
\item A \eng{volunteer} is always paid for his or her work.
\item \eng{Gender equality} means that men and women should have the same rights and opportunities.
\item \eng{To raise awareness} means to inform people about a problem.
\item \eng{Empowerment} means giving people more power and control over their own lives.
\end{enumerate}
\end{exercice}

\begin{exercice}[Vocabulaire : trouve le mot]
Complète chaque définition avec le mot anglais qui convient, choisi dans la liste : \eng{outreach} — \eng{gender gap} — \eng{advocacy} — \eng{community service} — \eng{role model}.
\begin{enumerate}
\item Work done for free to help people in your area : \dotfill
\item The difference between men and women in education, jobs or salaries : \dotfill
\item Activities that bring information and help to people in remote villages : \dotfill
\item Public support for an idea or a cause, for example women's rights : \dotfill
\item A person whose behaviour is an example for others to follow : \dotfill
\end{enumerate}
\end{exercice}

\begin{exercice}[La bonne forme du mot]
Complète chaque phrase avec la forme correcte du mot entre parenthèses.
\begin{enumerate}
\item Gender \dotfill (equal) is a human right.
\item Her \dotfill (commit) to the NGO is remarkable.
\item They fight \dotfill (discriminate) against girls.
\item The \dotfill (empower) of rural women is our main goal.
\item Many \dotfill (volunteer) joined the campaign last year.
\end{enumerate}
\end{exercice}

\courssub{Je m'entraîne}

\begin{exercice}[Traduction]
Traduis les phrases suivantes en anglais.
\begin{enumerate}
\item Beaucoup de jeunes font du bénévolat pour promouvoir l'égalité entre les sexes.
\item L'association organise une campagne de sensibilisation dans les écoles de Douala.
\item Les filles doivent avoir les mêmes chances que les garçons.
\item Nous devons lutter contre les violences faites aux femmes.
\item Son engagement envers la communauté est admirable.
\end{enumerate}
\end{exercice}

\begin{exercice}[Appariement : les collocations]
Associe chaque verbe (1--5) à la collocation qui convient (a--e).

\begin{center}
\begin{tabularx}{\linewidth}{|X|X|}
\hline
\rowcolor{popblueL} Verbe & Collocation \\
\hline
1. to raise & a. women and girls \\
\hline
2. to empower & b. against violence \\
\hline
3. to fight & c. the gender gap \\
\hline
4. to bridge & d. awareness \\
\hline
5. to speak out & e. discrimination \\
\hline
\end{tabularx}
\end{center}
\end{exercice}

\begin{exercice}[Texte à trous : une campagne de sensibilisation]
Complète le texte suivant avec les mots de la liste : \eng{volunteer} — \eng{awareness} — \eng{empowerment} — \eng{discrimination} — \eng{outreach} — \eng{commitment}.

\begin{textebox}[A campaign in Bamenda]
Last month, our school club decided to \dotfill{} in a local campaign for girls' education. The NGO organised an \dotfill{} campaign in three villages near Bamenda. During the \dotfill{} activities, we visited families and explained why sending girls to school matters. We also spoke against all forms of \dotfill{}, especially early marriage. The \dotfill{} of women is the main objective of the project. Thanks to the \dotfill{} of all the members, more than two hundred families received our message.
\end{textebox}
\end{exercice}

\begin{exercice}[Faux amis : choisis la bonne traduction]
Pour chaque expression française, choisis la bonne traduction anglaise (a ou b).
\begin{enumerate}
\item assister à une réunion : a. \eng{to assist a meeting} / b. \eng{to attend a meeting}
\item une formation : a. \eng{a formation} / b. \eng{a training course}
\item une sensibilisation : a. \eng{a sensitisation} / b. \eng{an awareness campaign}
\item une opportunité : a. \eng{an opportunity} / b. \eng{an opportuneness}
\item soutenir une cause : a. \eng{to support a cause} / b. \eng{to sustain a cause}
\end{enumerate}
\end{exercice}

\courssub{Je me prépare au Bac}

\begin{exercice}[Compréhension de l'écrit — type Bac]
Lis le texte suivant, puis réponds aux questions en anglais, avec des phrases complètes.

\begin{textebox}[Volunteering for equality in Cameroon]
In many towns and villages of Cameroon, young people are giving their time to promote gender equality. In Buea, a group of university students created an association called “Equal Future”. Every Saturday, its members visit secondary schools and talk to girls about their right to education. They also meet parents and traditional leaders to discuss early marriage, which still prevents many girls from finishing school.

The volunteers receive no salary, but they say the work changes their lives too. “I have learned to speak in public and to defend my ideas,” explains Clarisse, a nineteen-year-old volunteer. “When a girl tells me she will stay in school because of our visit, I feel proud.”

The association also trains women in small business management, because economic independence is an important part of empowerment. Last year, “Equal Future” helped more than fifty women start income-generating activities, from selling vegetables to tailoring. The members believe that real change begins in the community, one family at a time.
\end{textebox}

\textbf{A. Comprehension questions}
\begin{enumerate}
\item What do the members of “Equal Future” do every Saturday?
\item According to the text, what problem prevents many girls from finishing school?
\item Are the volunteers paid for their work? Justify your answer with a sentence from the text.
\item What has Clarisse learned from her volunteer work?
\item Why does the association train women in small business management?
\end{enumerate}

\textbf{B. Vocabulary in context}

Trouve dans le texte un mot ou une expression qui signifie :
\begin{enumerate}
\item equality between men and women (paragraph 1)
\item to stop something from happening (paragraph 1)
\item money received for work (paragraph 2)
\item activities that bring in money (paragraph 3)
\end{enumerate}
\end{exercice}

\begin{exercice}[Traduction — type Bac]
Traduis en français le passage suivant, extrait du texte étudié :

\begin{textebox}[Passage to translate]
“The volunteers receive no salary, but they say the work changes their lives too. ‘I have learned to speak in public and to defend my ideas,' explains Clarisse, a nineteen-year-old volunteer. ‘When a girl tells me she will stay in school because of our visit, I feel proud.'”
\end{textebox}
\end{exercice}

\begin{exercice}[Expression écrite — type Bac]
\textbf{Topic :} Your school wants to recruit volunteers to promote gender equality in your town. Write a letter of application to the coordinator of the local NGO “Equal Future”, 25 Avenue de la Réunification, Yaoundé. Introduce yourself, explain why you want to volunteer, mention your qualities and your availability.

\medskip
\textbf{Consignes :}
\begin{itemize}
\item Write between \textbf{150 and 180 words}.
\item Use the layout of a formal letter (addresses, date, salutation, closing formula).
\item Use at least \textbf{five words or expressions} from the lesson : \eng{volunteer}, \eng{awareness campaign}, \eng{empowerment}, \eng{commitment}, \eng{gender equality}, \eng{outreach}, \eng{to fight discrimination}.
\end{itemize}
\end{exercice}

\newpage

\coursec{Corrigés des exercices}

\begin{corrige}[Exercice 1 — QCM : le bon mot]
\begin{enumerate}
\item \textbf{a. volunteer} — Orthographe correcte : \eng{volunteer} (« vo-lon-tier »). Attention au double \emph{e} final ; \eng{volonteer} est un calque du français « volontaire ».
\item \textbf{a. awareness} — On écrit \eng{awareness} avec un seul \emph{a} après \emph{w} : \eng{aware} + suffixe \emph{-ness}. « Sensibilisation » se dit \eng{awareness}, pas \emph{awarness}.
\item \textbf{b. commitment} — Un seul \emph{t} après \emph{commit} : \eng{commitment}. Le verbe \eng{to commit} double le \emph{t} devant \emph{-ing} et \emph{-ed} (\eng{committing}, \eng{committed}), mais pas devant \emph{-ment}.
\item \textbf{a. discrimination} — Un seul \emph{m} : \eng{discrimination} (du latin \emph{discriminare}). Ne pas écrire \emph{descrimination} : le préfixe est \emph{dis-}, pas \emph{des-}.
\end{enumerate}
\textbf{Astuce orthographique :} en anglais, les suffixes \emph{-ment}, \emph{-ness}, \emph{-tion} sont très fréquents pour former des noms abstraits : \eng{commit $\rightarrow$ commitment}, \eng{aware $\rightarrow$ awareness}, \eng{discriminate $\rightarrow$ discrimination}.
\end{corrige}

\begin{corrige}[Exercice 2 — Vrai ou faux ?]
\begin{enumerate}
\item \textbf{F} — \eng{A volunteer works for free; he or she receives no salary.} (Un bénévole travaille gratuitement ; il ne reçoit pas de salaire.)
\item \textbf{V} — \eng{Gender equality means that men and women have the same rights and the same opportunities.} (L'égalité entre les sexes signifie que les hommes et les femmes ont les mêmes droits et les mêmes chances.)
\item \textbf{V} — \eng{To raise awareness means to inform people about a problem so that they can act.} (Sensibiliser, c'est informer les gens d'un problème pour qu'ils agissent.)
\item \textbf{V} — \eng{Empowerment means giving people more power and control over their own lives.} (L'autonomisation consiste à donner aux gens plus de pouvoir et de contrôle sur leur propre vie.)
\end{enumerate}
\textbf{Point de vigilance :} la justification doit être une phrase complète avec sujet et verbe conjugué ; une simple répétition de la définition du cours est acceptée si elle est correcte.
\end{corrige}

\begin{corrige}[Exercice 3 — Vocabulaire : trouve le mot]
\begin{enumerate}
\item \textbf{community service} — travail gratuit au service de son quartier ou de son village.
\item \textbf{gender gap} — l'écart entre les hommes et les femmes (éducation, emploi, salaires). Prononciation : « djen-deur gap ».
\item \textbf{outreach} — activités de proximité qui vont vers les populations éloignées. Faux ami : ce n'est pas « atteindre » ; c'est « action de terrain / sensibilisation de proximité ».
\item \textbf{advocacy} — plaidoyer, défense publique d'une cause. Prononciation : « a-dvo-keu-ci ».
\item \textbf{role model} — modèle, exemple à suivre. Attention : on dit \eng{a role model for girls}, avec la préposition \emph{for}.
\end{enumerate}
\end{corrige}

\begin{corrige}[Exercice 4 — La bonne forme du mot]
\begin{enumerate}
\item \textbf{equality} — \eng{Gender equality is a human right.} Nom formé avec le suffixe \emph{-ity} : \eng{equal $\rightarrow$ equality}.
\item \textbf{commitment} — \eng{Her commitment to the NGO is remarkable.} Verbe \eng{to commit} + \emph{-ment}. Note la préposition : \eng{commitment to}.
\item \textbf{discrimination} — \eng{They fight discrimination against girls.} Nom avec \emph{-tion} ; préposition \emph{against}.
\item \textbf{empowerment} — \eng{The empowerment of rural women is our main goal.} \eng{to empower $\rightarrow$ empowerment} ; préposition \emph{of}.
\item \textbf{volunteers} — \eng{Many volunteers joined the campaign last year.} Nom comptable au pluriel après \eng{many} : \emph{volunteer} + \emph{s}.
\end{enumerate}
\textbf{Rappel de méthode :} identifier d'abord la nature du mot attendu (nom, verbe, adjectif) grâce à la place dans la phrase : après un article ou un possessif (\eng{gender}, \eng{her}, \eng{the}), il faut un nom.
\end{corrige}

\begin{corrige}[Exercice 5 — Traduction]
\begin{enumerate}
\item \eng{Many young people volunteer to promote gender equality.} — « faire du bénévolat » = \eng{to volunteer} ou \eng{to do voluntary work} ; « promouvoir » = \eng{to promote}.
\item \eng{The association organises an awareness campaign in the schools of Douala.} — « campagne de sensibilisation » = \eng{awareness campaign} (jamais \emph{sensibilisation campaign}).
\item \eng{Girls must have the same opportunities as boys.} — « les mêmes ... que » = \eng{the same ... as}. On peut aussi dire \eng{the same chances as boys}.
\item \eng{We must fight against violence against women.} — « lutter contre » = \eng{to fight against} ; « les violences faites aux femmes » = \eng{violence against women} (\emph{violence} est ici indénombrable, sans \emph{s}).
\item \eng{Her commitment to the community is admirable.} — « engagement envers » = \eng{commitment to}.
\end{enumerate}
\textbf{Points de vigilance :} ne pas calquer « faire du bénévolat » par \emph{to make volunteering} ; ne pas traduire « les mêmes que » par \emph{the same that}.
\end{corrige}

\begin{corrige}[Exercice 6 — Appariement : les collocations]
\begin{enumerate}
\item \textbf{1 -- d : to raise awareness} — sensibiliser (littéralement « élever la conscience »).
\item \textbf{2 -- a : to empower women and girls} — autonomiser les femmes et les filles.
\item \textbf{3 -- e : to fight discrimination} — lutter contre la discrimination.
\item \textbf{4 -- c : to bridge the gender gap} — réduire / combler l'écart entre les sexes. Image du pont (\eng{bridge}) qui relie deux rives.
\item \textbf{5 -- b : to speak out against violence} — dénoncer publiquement la violence. \eng{To speak out against} = prendre la parole pour dénoncer.
\end{enumerate}
\textbf{À retenir :} ces cinq collocations sont des « blocs » à mémoriser entiers ; elles sont très fréquentes dans les textes du Bac sur les droits humains.
\end{corrige}

\begin{corrige}[Exercice 7 — Texte à trous : une campagne de sensibilisation]
Texte complété :
\begin{enumerate}
\item \textbf{volunteer} — après \eng{decided to}, il faut la base verbale : \eng{decided to volunteer}.
\item \textbf{awareness} — \eng{an awareness campaign} : article indéfini \emph{an} devant une voyelle.
\item \textbf{outreach} — \eng{during the outreach activities} : les activités de terrain.
\item \textbf{discrimination} — \eng{spoke against all forms of discrimination}.
\item \textbf{empowerment} — \eng{the empowerment of women is the main objective}.
\item \textbf{commitment} — \eng{thanks to the commitment of all the members}.
\end{enumerate}
\textbf{Traduction du texte complété :} « Le mois dernier, notre club scolaire a décidé de faire du bénévolat dans une campagne locale pour l'éducation des filles. L'ONG a organisé une campagne de sensibilisation dans trois villages près de Bamenda. Pendant les activités de terrain, nous avons rendu visite aux familles et expliqué pourquoi il est important d'envoyer les filles à l'école. Nous avons aussi dénoncé toutes les formes de discrimination, en particulier le mariage précoce. L'autonomisation des femmes est l'objectif principal du projet. Grâce à l'engagement de tous les membres, plus de deux cents familles ont reçu notre message. »
\end{corrige}

\begin{corrige}[Exercice 8 — Faux amis : choisis la bonne traduction]
\begin{enumerate}
\item \textbf{b. to attend a meeting} — \eng{to assist} signifie « aider » ; « assister à » (être présent) se dit \eng{to attend}.
\item \textbf{b. a training course} — \eng{formation} n'existe pas en anglais dans ce sens ; on dit \eng{training} ou \eng{a training course}.
\item \textbf{b. an awareness campaign} — \eng{sensitisation} est un anglicisme rare et maladroit ; l'anglais authentique dit \eng{awareness campaign} ou \eng{awareness-raising}.
\item \textbf{a. an opportunity} — \eng{opportuneness} n'existe pas ; attention, \eng{opportunity} signifie « occasion favorable », pas « opportunité » au sens de « caractère opportun ».
\item \textbf{a. to support a cause} — \eng{to sustain} signifie « maintenir, subvenir aux besoins » ; « soutenir une cause » = \eng{to support a cause}.
\end{enumerate}
\textbf{Point de vigilance :} les mots transparents (\emph{assist}, \emph{formation}, \emph{sensitisation}, \emph{sustain}) sont les pièges classiques du Bac : toujours vérifier le sens réel du mot anglais avant de l'employer.
\end{corrige}

\begin{corrige}[Exercice 9 — Compréhension de l'écrit — type Bac]
\textbf{A. Comprehension questions}
\begin{enumerate}
\item \eng{Every Saturday, the members of “Equal Future” visit secondary schools and talk to girls about their right to education. They also meet parents and traditional leaders to discuss early marriage.}
\item \eng{Early marriage prevents many girls from finishing school.}
\item \eng{No, they are not paid. The text says: “The volunteers receive no salary.”}
\item \eng{Clarisse has learned to speak in public and to defend her ideas.}
\item \eng{The association trains women in small business management because economic independence is an important part of empowerment.}
\end{enumerate}

\textbf{B. Vocabulary in context}
\begin{enumerate}
\item \eng{gender equality} (paragraph 1)
\item \eng{to prevent} (paragraph 1) — \eng{to prevent somebody from doing something}.
\item \eng{salary} (paragraph 2)
\item \eng{income-generating activities} (paragraph 3)
\end{enumerate}

\textbf{Barème indicatif (10 points) :} Partie A : 5 $\times$ 1,5 pt (phrase complète, information exacte, anglais correct) ; Partie B : 4 $\times$ 0,5 pt. Retirer 0,25 pt par faute de grammaire ou de mot mal recopié.

\textbf{Points de vigilance :} répondre par des phrases complètes, jamais par un seul mot ; reprendre le vocabulaire du texte sans le déformer ; à la question 3, la citation doit être exacte.
\end{corrige}

\begin{corrige}[Exercice 10 — Traduction — type Bac]
\textbf{Proposition de traduction :}

« Les bénévoles ne reçoivent aucun salaire, mais ils affirment que ce travail change aussi leur propre vie. “J'ai appris à parler en public et à défendre mes idées”, explique Clarisse, une bénévole de dix-neuf ans. “Quand une fille me dit qu'elle va rester à l'école grâce à notre visite, je suis fière.” »

\textbf{Barème indicatif (5 points) :} 1 pt par segment correctement traduit (5 segments), avec tolérance pour les formulations équivalentes.

\textbf{Points de vigilance :}
\begin{itemize}
\item \eng{receive no salary} = « ne reçoivent aucun salaire » (négation portée par \emph{no}, pas de double négation en français).
\item \eng{changes their lives too} = « change aussi leur vie » ; ne pas oublier \emph{too}.
\item \eng{a nineteen-year-old volunteer} = « une bénévole de dix-neuf ans » : l'adjectif composé anglais se traduit par « de ... ans ».
\item \eng{because of our visit} = « grâce à notre visite » (cause positive) ; \eng{I feel proud} = « je suis fière » (présent de vérité générale).
\end{itemize}
\end{corrige}

\begin{corrige}[Exercice 11 — Expression écrite — type Bac]
\textbf{Proposition de rédaction modèle (environ 165 mots) :}

\begin{textebox}[Model letter]
15 Rue des Hibiscus\\
Biyem-Assi, Yaoundé\\
15th May 2025

The Coordinator\\
“Equal Future” NGO\\
25 Avenue de la Réunification\\
Yaoundé

Dear Sir or Madam,

I am writing to apply for a volunteer position in your association. I am a seventeen-year-old student in Terminale at Lycée de Biyem-Assi, and I strongly believe in gender equality.

Last year, I took part in an awareness campaign against early marriage in my neighbourhood, and I saw how much girls need support to stay in school. I want to volunteer with “Equal Future” because I admire your commitment to the empowerment of women and girls.

I am hard-working, patient and I speak English and French fluently, which is useful for outreach activities in both communities. I am available every Saturday and during the school holidays.

I would be honoured to help you fight discrimination and bridge the gender gap in our town. Thank you for considering my application.

Yours faithfully,

Amina Mbarga
\end{textebox}

\textbf{Barème indicatif (10 points) :}
\begin{itemize}
\item Mise en page de la lettre formelle (adresses, date, formules) : 2 pts.
\item Respect des consignes (150--180 mots, 5 mots du lexique) : 2 pts.
\item Correction grammaticale et orthographique : 3 pts.
\item Cohérence, clarté et richesse du contenu : 3 pts.
\end{itemize}

\textbf{Points de vigilance :}
\begin{itemize}
\item Date anglaise : \eng{15th May 2025} ou \eng{May 15th, 2025}, sans « le ».
\item Formules : \eng{Dear Sir or Madam} $\rightarrow$ \eng{Yours faithfully} ; si le nom du destinataire est connu : \eng{Dear Mr X} $\rightarrow$ \eng{Yours sincerely}.
\item Écrire \eng{I am writing to apply for...} (présent continu), pas \emph{I write for...}.
\item Ne pas traduire « je suis disponible » par \emph{I am disposable} : dire \eng{I am available}.
\item Employer le lexique de la leçon de façon naturelle, sans forcer.
\end{itemize}
\end{corrige}

\newpage

\coursec{Fiche bilan}

\begin{popfigure}
\begin{center}\tcbox[colback=popblue,colframe=popblue,coltext=white,arc=9pt,boxsep=4pt,left=6pt,right=6pt]{\bfseries\small Volunteering for Gender Equality}\end{center}
\vspace{-2pt}
\begin{tcbraster}[raster columns=3,raster equal height=rows,raster column skip=6pt,raster row skip=6pt,enhanced,blankest]
\begin{tcolorbox}[enhanced,colback=poporange,colframe=poporange,coltext=white,boxrule=0.9pt,arc=3pt,left=4pt,right=4pt,top=3pt,bottom=3pt,boxsep=1pt,halign=left]\bfseries\scriptsize Volunteering \& community service\end{tcolorbox}
\begin{tcolorbox}[enhanced,colback=popgreen,colframe=popgreen,coltext=white,boxrule=0.9pt,arc=3pt,left=4pt,right=4pt,top=3pt,bottom=3pt,boxsep=1pt,halign=left]\bfseries\scriptsize Gender equality \& women's rights\end{tcolorbox}
\begin{tcolorbox}[enhanced,colback=poppurple,colframe=poppurple,coltext=white,boxrule=0.9pt,arc=3pt,left=4pt,right=4pt,top=3pt,bottom=3pt,boxsep=1pt,halign=left]\bfseries\scriptsize Collocations \& idioms\end{tcolorbox}
\begin{tcolorbox}[enhanced,colback=poppink,colframe=poppink,coltext=white,boxrule=0.9pt,arc=3pt,left=4pt,right=4pt,top=3pt,bottom=3pt,boxsep=1pt,halign=left]\bfseries\scriptsize False friends \& traps\end{tcolorbox}
\begin{tcolorbox}[enhanced,colback=popgold,colframe=popgold,coltext=white,boxrule=0.9pt,arc=3pt,left=4pt,right=4pt,top=3pt,bottom=3pt,boxsep=1pt,halign=left]\bfseries\scriptsize Word families \& spelling\end{tcolorbox}
\begin{tcolorbox}[enhanced,colback=popblue!60,colframe=popblue!60,coltext=white,boxrule=0.9pt,arc=3pt,left=4pt,right=4pt,top=3pt,bottom=3pt,boxsep=1pt,halign=left]\bfseries\scriptsize Using the lexis in context\end{tcolorbox}
\end{tcbraster}

\legende{Carte mentale de la leçon : le volontariat pour l'égalité des sexes}
\end{popfigure}

\begin{bilan}
\textbf{1. Lexique clé : volunteering and community service}
\begin{itemize}
\item \eng{a volunteer} : un bénévole, un volontaire ; \eng{to volunteer (for)} : se porter volontaire (pour)
\item \eng{voluntary work / community service} : travail bénévole / service communautaire
\item \eng{an NGO (non-governmental organisation)} : une ONG ; \eng{a charity} : une œuvre caritative
\item \eng{to raise awareness (of/about)} : sensibiliser (à) ; \eng{an awareness campaign} : une campagne de sensibilisation
\item \eng{to raise funds} : collecter des fonds ; \eng{a fundraiser} : une collecte de fonds
\item \eng{to help out / to lend a hand} : donner un coup de main ; \eng{to make a difference} : changer les choses
\item \eng{to commit oneself to a cause} : s'engager pour une cause ; \eng{commitment} : l'engagement
\end{itemize}

\textbf{2. Lexique clé : gender equality and women's rights}
\begin{itemize}
\item \eng{gender equality} : l'égalité des sexes ; \eng{gender gap} : l'écart entre les sexes
\item \eng{women's rights} : les droits des femmes ; \eng{to empower women} : autonomiser les femmes ; \eng{empowerment} : l'autonomisation
\item \eng{equal pay} : salaire égal ; \eng{equal opportunities} : l'égalité des chances
\item \eng{discrimination (against)} : la discrimination (envers) ; \eng{sexism} : le sexisme ; \eng{a stereotype} : un stéréotype
\item \eng{to fight / to combat gender-based violence} : lutter contre les violences basées sur le genre
\item \eng{girls' education} : l'éducation des filles ; \eng{to promote / to advocate} : promouvoir / défendre
\end{itemize}

\textbf{3. Collocations et expressions à connaître par cœur}
\begin{center}
\begin{tabularx}{\linewidth}{|X|X|}
\hline
\rowcolor{popblueL} Collocation anglaise & Traduction française \\
\hline
\eng{to volunteer one's time} & donner bénévolement de son temps \\
\hline
\eng{to take part in a campaign} & participer à une campagne \\
\hline
\eng{to break down stereotypes} & briser les stéréotypes \\
\hline
\eng{to bridge the gender gap} & réduire l'écart entre les sexes \\
\hline
\eng{to stand up for one's rights} & défendre ses droits \\
\hline
\eng{to give women a voice} & donner la parole aux femmes \\
\hline
\eng{to achieve gender equality} & parvenir à l'égalité des sexes \\
\hline
\end{tabularx}
\end{center}

\textbf{4. Faux amis et pièges à éviter}
\begin{itemize}
\item \eng{to achieve} = réaliser (un objectif), jamais « achever » au sens de finir un travail.
\item \eng{sensible} = raisonnable ; « sensible » se dit \eng{sensitive}.
\item \eng{actually} = en fait ; « actuellement » se dit \eng{currently}.
\item \eng{a volunteer} se prononce « vo-lon-TIRE » (accent sur la dernière syllabe) ; \eng{voluntary} se prononce « VO-lon-tri ».
\item On dit \eng{to be involved IN}, \eng{to participate IN}, \eng{to consist OF} : attention aux prépositions.
\end{itemize}

\textbf{5. Familles de mots}
\begin{center}
\begin{tabularx}{\linewidth}{|X|X|X|}
\hline
\rowcolor{popblueL} Nom & Verbe & Adjectif \\
\hline
\eng{equality} & \eng{to equalise} & \eng{equal} \\
\hline
\eng{discrimination} & \eng{to discriminate} & \eng{discriminatory} \\
\hline
\eng{empowerment} & \eng{to empower} & \eng{empowered} \\
\hline
\eng{volunteer} & \eng{to volunteer} & \eng{voluntary} \\
\hline
\eng{promotion} & \eng{to promote} & \eng{promotional} \\
\hline
\end{tabularx}
\end{center}

\textbf{6. Méthode express}
\begin{enumerate}
\item Repérer le thème du texte (volontariat, droits des femmes) et souligner les mots-clés.
\item Deviner le sens des mots inconnus grâce au contexte et à la famille de mots.
\item Réemployer au moins cinq mots du lexique dans toute production orale ou écrite.
\item Vérifier prépositions et collocations avant de rendre la copie.
\end{enumerate}
\end{bilan}

\begin{aretenir}[Checklist avant le Bac]
\begin{itemize}
\item[$\square$] Je sais définir \eng{volunteering}, \eng{community service} et \eng{NGO} en anglais.
\item[$\square$] Je connais au moins dix mots du champ lexical de l'égalité des sexes.
\item[$\square$] Je sais employer \eng{to raise awareness of} et \eng{to raise funds} dans une phrase.
\item[$\square$] Je maîtrise les collocations \eng{break down stereotypes}, \eng{bridge the gender gap}, \eng{stand up for one's rights}.
\item[$\square$] Je ne confonds plus \eng{achieve} et « achever », \eng{actually} et « actuellement ».
\item[$\square$] Je place correctement l'accent dans \eng{volunteer} et \eng{voluntary}.
\item[$\square$] Je connais les familles de mots : \eng{equal / equality / equalise}, \eng{empower / empowerment}.
\item[$\square$] J'utilise les bonnes prépositions : \eng{involved in}, \eng{participate in}, \eng{discrimination against}.
\item[$\square$] Je sais écrire un court paragraphe sur une action bénévole dans ma communauté.
\item[$\square$] Je peux parler une minute, sans notes, d'une campagne de sensibilisation à l'égalité des sexes.
\end{itemize}
\end{aretenir}

\findecours

\end{document}
